\documentclass[11pt,a4paper]{article}
\usepackage[utf8]{inputenc}
\usepackage[spanish,es-nodecimaldot,es-noquoting]{babel}
\usepackage{amsmath,amssymb,amsfonts,amsthm,mathtools}
\usepackage{geometry}
\geometry{top=2.2cm, bottom=2.2cm, left=2.2cm, right=2.2cm}
\usepackage[table]{xcolor}
\usepackage{fancyhdr}
\usepackage{enumitem}
\usepackage{titlesec}
\usepackage{tcolorbox}
\tcbuselibrary{skins,breakable}
\usepackage{tabularx}
\usepackage{array}
\usepackage{hyperref}

\usepackage{tikz}
\usetikzlibrary{arrows.meta,calc,patterns,decorations.markings,3d}
\usepackage{pgfplots}
\pgfplotsset{compat=1.17}
\usepgfplotslibrary{fillbetween}

% Paleta de colores académica
\definecolor{navyblue}{RGB}{16, 44, 87}
\definecolor{tealblue}{RGB}{20, 85, 120}
\definecolor{darkslate}{RGB}{35, 45, 55}
\definecolor{forestgreen}{RGB}{24, 94, 60}
\definecolor{crimsonred}{RGB}{140, 25, 35}
\definecolor{amberorange}{RGB}{180, 90, 10}

\definecolor{bgblue}{RGB}{245, 248, 253}
\definecolor{bggreen}{RGB}{244, 251, 246}
\definecolor{bgred}{RGB}{254, 245, 245}
\definecolor{bgyellow}{RGB}{255, 251, 240}
\definecolor{bgslate}{RGB}{248, 249, 251}

\hypersetup{
    colorlinks=true,
    linkcolor=navyblue,
    citecolor=tealblue,
    urlcolor=tealblue,
    pdfauthor={Antigravity - Guía de Estudio UNI},
    pdftitle={Teoría Rigurosa y Práctica - PC05 CM2A1 (2026-I / 2026-II)}
}

% Encabezados y pie de página
\pagestyle{fancy}
\fancyhf{}
\renewcommand{\headrulewidth}{0.6pt}
\fancyhead[L]{\small\scshape UNI $\cdot$ Facultad de Ciencias $\cdot$ Escuela de Matemática}
\fancyhead[R]{\small\scshape CM2A1: Cálculo Avanzado (Teoría PC05)}
\fancyfoot[C]{\thepage}
\renewcommand{\footrulewidth}{0.4pt}

% Estilo de Cajas con tcolorbox
\newtcolorbox{definicionbox}[1]{
    enhanced, breakable, colback=bgblue, colframe=navyblue,
    fonttitle=\bfseries\color{white}, title={Definición: #1},
    arc=3mm, boxrule=0.9pt, left=4mm, right=4mm, top=3mm, bottom=3mm, drop shadow=black!10
}

\newtcolorbox{teoremabox}[1]{
    enhanced, breakable, colback=bgblue!70!white, colframe=tealblue,
    fonttitle=\bfseries\color{white}, title={Teorema / Proposición: #1},
    arc=3mm, boxrule=0.9pt, left=4mm, right=4mm, top=3mm, bottom=3mm, drop shadow=black!10
}

\newtcolorbox{ejemplobox}[1]{
    enhanced, breakable, colback=bggreen, colframe=forestgreen,
    fonttitle=\bfseries\color{white}, title={Ejemplo Mínimo Ilustrativo: #1},
    arc=2mm, boxrule=0.8pt, left=4mm, right=4mm, top=3mm, bottom=3mm
}

\newtcolorbox{ejerciciobox}[1]{
    enhanced, breakable, colback=bgslate, colframe=darkslate,
    fonttitle=\bfseries\color{white}, title={Ejercicio Resuelto Paso a Paso: #1},
    arc=2mm, boxrule=0.9pt, left=4mm, right=4mm, top=3mm, bottom=3mm, drop shadow=black!10
}

\newtcolorbox{errorbox}[1]{
    enhanced, breakable, colback=bgred, colframe=crimsonred,
    fonttitle=\bfseries\color{white}, title={¡Cuidado! Errores Frecuentes y Trampas Conceptuales: #1},
    arc=2mm, boxrule=0.8pt, left=4mm, right=4mm, top=3mm, bottom=3mm
}

\newtcolorbox{observacionbox}[1]{
    enhanced, breakable, colback=bgyellow, colframe=amberorange,
    fonttitle=\bfseries\color{white}, title={Observación Metodológica: #1},
    arc=2mm, boxrule=0.8pt, left=4mm, right=4mm, top=3mm, bottom=3mm
}

% Formato de secciones
\titleformat{\section}{\Large\bfseries\color{navyblue}}{\thesection.}{0.8em}{}[\titlerule]
\titleformat{\subsection}{\large\bfseries\color{tealblue}}{\thesubsection.}{0.6em}{}


% Cajas Adicionales para Compendio PC05
\newtcolorbox{traductorbox}[1]{
    enhanced, breakable, colback=bgslate, colframe=navyblue,
    fonttitle=\bfseries\color{white}, title={Traductor de Jerga Formal: #1},
    arc=2mm, boxrule=0.8pt, left=3mm, right=3mm, top=3mm, bottom=3mm, drop shadow=black!10
}
\newtcolorbox{faseunobox}[1]{
    enhanced, breakable, colback=bgyellow, colframe=amberorange,
    fonttitle=\bfseries\color{white}, title={Fase 1: Borrador Heurístico e Intuición (#1)},
    arc=2mm, boxrule=0.8pt, left=3mm, right=3mm, top=3mm, bottom=3mm
}
\newtcolorbox{fasedosbox}[1]{
    enhanced, breakable, colback=bgblue, colframe=tealblue,
    fonttitle=\bfseries\color{white}, title={Fase 2: Demostración Formal Rigurosa (#1)},
    arc=2mm, boxrule=0.8pt, left=3mm, right=3mm, top=3mm, bottom=3mm
}
\newtcolorbox{nivelbox}[2]{
    enhanced, breakable, colback=white, colframe=darkslate,
    fonttitle=\bfseries\color{white}, title={Ejercicio #1 \quad $\cdot$ \quad #2},
    arc=2mm, boxrule=0.8pt, left=3mm, right=3mm, top=3mm, bottom=3mm, drop shadow=black!10
}
\begin{document}

% Portada / Cabecera Principal
\begin{center}
    {\large \scshape Universidad Nacional de Ingeniería}\\[0.15cm]
    {\large \scshape Facultad de Ciencias --- Escuela Profesional de Matemática}\\[0.35cm]
    {\LARGE \bfseries \color{navyblue} Fundamentos Teóricos y Métodos de Solución:\\Quinta Práctica Calificada (PC05)}\\[0.25cm]
    {\large \textbf{Curso:} Cálculo Diferencial e Integral Avanzado (CM2A1) \quad $\cdot$ \quad \textbf{Nivel:} Pregrado en Matemática}\\[0.2cm]
    {\small \textbf{Cobertura del Sílabo Dosificado:} Sesiones 25 a 30}
\end{center}
\vspace{-0.2cm}
\hrule height 1.2pt
\vspace{0.4cm}

\begin{abstract}
Este compendio proporciona el marco teórico riguroso, analítico y didáctico necesario para dominar con solidez los fundamentos evaluados en la \textbf{Quinta Práctica Calificada (PC05)} del curso CM2A1 de la Escuela Profesional de Matemática de la Universidad Nacional de Ingeniería. Abarca desde la teoría de integrales dobles y triples sobre regiones elementales y conjuntos de contenido nulo, el poderoso Teorema de Fubini para el cálculo de volúmenes e integrales iteradas, hasta los sistemas de coordenadas polares, cilíndricas y esféricas. Incluye el análisis exhaustivo del Teorema de Cambio de Variable y generalizaciones a $\mathbb{R}^n$. Cada sección integra definiciones formales con nivel de rigor matemático puro, ejemplos mínimos transparentes, ejercicios tipo examen de semestres recientes (2024-II, 2025-I) resueltos minuciosamente paso a paso sin omitir detalles algebraicos ni topológicos, y una sección crítica sobre errores comunes y trampas conceptuales típicas en evaluaciones UNI.
\end{abstract}

\tableofcontents
\vspace{0.5cm}
\hrule
\vspace{0.6cm}

\part*{Parte I: Regiones elementales y contenido nulo}

\section{Regiones Elementales en $\mathbb{R}^2$}

Para formalizar el estudio de la integración múltiple en dimensiones superiores, es imperativo primero dotar de rigor a los dominios sobre los cuales estas integrales son evaluadas. En el plano cartesiano $\mathbb{R}^2$, los dominios más tratables geométricamente y analíticamente son las denominadas regiones elementales.

\begin{definicionbox}{Regiones Elementales en $\mathbb{R}^2$}
Una región $D \subset \mathbb{R}^2$ se denomina \textbf{región de Tipo 1} (o elemental respecto al eje $y$) si puede ser descrita en la forma:
\[
D = \left\{ (x,y) \in \mathbb{R}^2 \mid a \leq x \leq b, \; \phi_1(x) \leq y \leq \phi_2(x) \right\}
\]
donde $a, b \in \mathbb{R}$ con $a < b$, y $\phi_1, \phi_2 : [a,b] \to \mathbb{R}$ son funciones continuas tales que $\phi_1(x) \leq \phi_2(x)$ para todo $x \in [a,b]$.

Análogamente, una región $D \subset \mathbb{R}^2$ se denomina \textbf{región de Tipo 2} (o elemental respecto al eje $x$) si puede ser expresada como:
\[
D = \left\{ (x,y) \in \mathbb{R}^2 \mid c \leq y \leq d, \; \psi_1(y) \leq x \leq \psi_2(y) \right\}
\]
donde $c, d \in \mathbb{R}$ con $c < d$, y $\psi_1, \psi_2 : [c,d] \to \mathbb{R}$ son funciones continuas tales que $\psi_1(y) \leq \psi_2(y)$ para todo $y \in [c,d]$.
\end{definicionbox}

\begin{observacionbox}{Regiones de ambos tipos}
Es común encontrar regiones que satisfacen simultáneamente las condiciones de ser de Tipo 1 y de Tipo 2. Por ejemplo, todo disco cerrado o rectángulo con lados paralelos a los ejes coordenados es una región de ambos tipos. La flexibilidad en la elección del tipo de región es crucial para simplificar el cálculo de integrales dobles mediante el Teorema de Fubini.
\end{observacionbox}

\begin{ejemplobox}{Determinación de tipos de regiones}
Considere la región $D$ delimitada por la parábola $y = x^2$ y la recta $y = 2x$.
Para describirla como una región de Tipo 1, primero hallamos las intersecciones de las curvas igualando las ordenadas:
\[
x^2 = 2x \implies x(x-2) = 0 \implies x = 0 \quad \lor \quad x = 2.
\]
En el intervalo $[0, 2]$, se cumple que $x^2 \leq 2x$. Por lo tanto, $D$ como región de Tipo 1 es:
\[
D = \{ (x,y) \in \mathbb{R}^2 \mid 0 \leq x \leq 2, \; x^2 \leq y \leq 2x \}.
\]
Para describirla como una región de Tipo 2, despejamos $x$ en función de $y$. De $y = x^2$ para $x \geq 0$ obtenemos $x = \sqrt{y}$, y de $y = 2x$ obtenemos $x = y/2$. El intervalo de variación para $y$ es desde $y=0$ hasta $y=4$ (obtenido evaluando en los puntos de intersección). En $[0, 4]$, se tiene $y/2 \leq \sqrt{y}$. Por ende, como región de Tipo 2, tenemos:
\[
D = \left\{ (x,y) \in \mathbb{R}^2 \mid 0 \leq y \leq 4, \; \frac{y}{2} \leq x \leq \sqrt{y} \right\}.
\]
\begin{center}
\begin{tikzpicture}
\begin{axis}[
    axis lines = middle,
    xlabel = {$x$},
    ylabel = {$y$},
    xmin = -0.5, xmax = 2.5,
    ymin = -0.5, ymax = 4.5,
    domain = 0:2,
    samples = 100,
    width = 8cm,
    height = 8cm
]
\addplot [name path=A, blue, thick] {2*x};
\addplot [name path=B, red, thick] {x^2};
\addplot [gray!30] fill between[of=A and B];
\node[blue, left] at (axis cs: 1.5, 3) {$y=2x$};
\node[red, right] at (axis cs: 1.5, 2.25) {$y=x^2$};
\end{axis}
\end{tikzpicture}
\end{center}
\end{ejemplobox}

\begin{errorbox}{Descomposición de regiones}
Un error frecuente al tratar con regiones delimitadas por curvas cerradas complejas o regiones con "agujeros" (no simplemente conexas) es asumir que toda región medible de Jordan puede describirse mediante un único conjunto de desigualdades del Tipo 1 o Tipo 2. 
En general, una región arbitraria debe descomponerse en una unión finita de subregiones $D_i$, cada una de las cuales es elemental, satisfaciendo que sus interiores son disjuntos por pares ($\mathring{D}_i \cap \mathring{D}_j = \emptyset$ para $i \neq j$).
\end{errorbox}

\section{Sólidos Elementales en $\mathbb{R}^3$}

La generalización de estas ideas al espacio tridimensional introduce la noción de sólidos elementales, definidos en base a sus proyecciones sobre los planos coordenados.

\begin{definicionbox}{Sólidos Elementales}
Un sólido $E \subset \mathbb{R}^3$ es considerado un \textbf{sólido elemental} si puede expresarse de una de las tres formas siguientes:

\textbf{Tipo 1 (Proyección en el plano $xy$):}
\[
E = \left\{ (x,y,z) \in \mathbb{R}^3 \mid (x,y) \in D_{xy}, \; u_1(x,y) \leq z \leq u_2(x,y) \right\}
\]
donde $D_{xy}$ es la proyección de $E$ sobre el plano $xy$ (y se asume que $D_{xy}$ es una región elemental en $\mathbb{R}^2$), y $u_1, u_2 : D_{xy} \to \mathbb{R}$ son funciones continuas.

\textbf{Tipo 2 (Proyección en el plano $xz$):}
\[
E = \left\{ (x,y,z) \in \mathbb{R}^3 \mid (x,z) \in D_{xz}, \; v_1(x,z) \leq y \leq v_2(x,z) \right\}
\]
donde $D_{xz}$ es la proyección de $E$ sobre el plano $xz$ (región elemental en $\mathbb{R}^2$), y $v_1, v_2 : D_{xz} \to \mathbb{R}$ son funciones continuas.

\textbf{Tipo 3 (Proyección en el plano $yz$):}
\[
E = \left\{ (x,y,z) \in \mathbb{R}^3 \mid (y,z) \in D_{yz}, \; w_1(y,z) \leq x \leq w_2(y,z) \right\}
\]
donde $D_{yz}$ es la proyección de $E$ sobre el plano $yz$ (región elemental en $\mathbb{R}^2$), y $w_1, w_2 : D_{yz} \to \mathbb{R}$ son funciones continuas.
\end{definicionbox}

\begin{ejemplobox}{Sólido de revolución y curva de intersección}
Sea $E$ la región delimitada por el paraboloide $z = x^2 + y^2$ y el plano $z = 4$.
Para determinar la proyección $D_{xy}$, igualamos las superficies para encontrar su curva de intersección:
\[
x^2 + y^2 = 4.
\]
Esta curva, en el plano $z=4$, se proyecta sobre el plano $xy$ como la circunferencia $x^2 + y^2 = 4$. Por lo tanto, el dominio de proyección $D_{xy}$ es el disco cerrado:
\[
D_{xy} = \{ (x,y) \in \mathbb{R}^2 \mid x^2 + y^2 \leq 4 \}.
\]
El sólido $E$ puede describirse como uno de Tipo 1:
\[
E = \{ (x,y,z) \in \mathbb{R}^3 \mid (x,y) \in D_{xy}, \; x^2 + y^2 \leq z \leq 4 \}.
\]
Describiendo $D_{xy}$ como región de Tipo 1 en $\mathbb{R}^2$:
\[
E = \{ (x,y,z) \in \mathbb{R}^3 \mid -2 \leq x \leq 2, \; -\sqrt{4-x^2} \leq y \leq \sqrt{4-x^2}, \; x^2+y^2 \leq z \leq 4 \}.
\]
\begin{center}
\begin{tikzpicture}
\begin{axis}[
    view={60}{30},
    xlabel={$x$}, ylabel={$y$}, zlabel={$z$},
    zmin=0, zmax=5,
    colormap/viridis
]
\addplot3[
    surf,
    domain=-2:2,
    domain y=-2:2,
    restrict expr to domain={x^2+y^2}{0:4},
    samples=30
] {x^2 + y^2};
\addplot3[
    surf,
    domain=-2:2,
    domain y=-2:2,
    restrict expr to domain={x^2+y^2}{0:4},
    opacity=0.5,
    fill=gray
] {4};
\end{axis}
\end{tikzpicture}
\end{center}
\end{ejemplobox}

\section{Conjuntos de Contenido Nulo}

Un concepto fundamental en la teoría de la integración de Riemann en $\mathbb{R}^n$ es el de conjunto de medida nula y, de forma más restrictiva, el de conjunto de contenido nulo. En nuestro contexto, dado que las integrales de Riemann operan sobre dominios compactos, el contenido nulo provee la caracterización exacta de los conjuntos que no contribuyen al valor de la integral.

\begin{definicionbox}{Contenido Nulo}
Un subconjunto $A \subset \mathbb{R}^n$ tiene \textbf{contenido nulo} si, para todo $\varepsilon > 0$, existe un número \emph{finito} de rectángulos cerrados $R_1, R_2, \dots, R_k$ en $\mathbb{R}^n$ tales que:
\begin{enumerate}
    \item $A \subset \bigcup_{i=1}^k R_i$.
    \item $\sum_{i=1}^k \operatorname{vol}(R_i) < \varepsilon$.
\end{enumerate}
donde $\operatorname{vol}(R)$ denota el volumen $n$-dimensional del rectángulo $R$.
\end{definicionbox}

\begin{observacionbox}{Relación con la medida nula de Lebesgue}
Mientras que la medida nula permite un recubrimiento contable infinito de rectángulos, el contenido nulo exige un recubrimiento finito. Todo conjunto de contenido nulo tiene medida nula, pero el recíproco no es cierto en general (por ejemplo, $\mathbb{Q} \cap [0,1]$ tiene medida nula por ser numerable, pero su contenido no es nulo pues su clausura es el intervalo $[0,1]$ compacto de volumen 1, y cualquier recubrimiento finito cerrado de $\mathbb{Q} \cap [0,1]$ recubre inevitablemente a todo $[0,1]$). No obstante, para conjuntos \textbf{compactos}, ambas nociones coinciden por el Teorema de Heine-Borel.
\end{observacionbox}

\begin{teoremabox}{Gráfica de una función continua}
Sea $f: [a,b] \to \mathbb{R}$ una función continua. Entonces, el grafo de $f$, definido como:
\[
G(f) = \{ (x, f(x)) \in \mathbb{R}^2 \mid x \in [a,b] \}
\]
tiene contenido nulo en $\mathbb{R}^2$.
\end{teoremabox}

\textbf{Demostración:}
Sea $\varepsilon > 0$ un número real arbitrario. Debemos demostrar la existencia de un recubrimiento finito de $G(f)$ mediante rectángulos cuya suma de áreas sea estrictamente menor a $\varepsilon$.

Dado que $f$ es continua en el intervalo cerrado y acotado $[a,b]$, por el Teorema de Heine-Cantor, $f$ es \textbf{uniformemente continua} en $[a,b]$.
La definición de continuidad uniforme nos asegura que para el valor $\varepsilon' = \frac{\varepsilon}{2(b-a)} > 0$, existe un $\delta > 0$ tal que para cualesquiera $x, x' \in [a,b]$:
\[
|x - x'| < \delta \implies |f(x) - f(x')| < \varepsilon' = \frac{\varepsilon}{2(b-a)}.
\]

Elijamos una partición regular $P = \{x_0, x_1, \dots, x_N\}$ del intervalo $[a,b]$ de modo que la norma de la partición sea menor que $\delta$. Es decir, dividimos $[a,b]$ en $N$ subintervalos de igual longitud $\Delta x = \frac{b-a}{N}$, eligiendo $N \in \mathbb{N}$ suficientemente grande tal que $\Delta x < \delta$. Tenemos entonces:
\[
x_i = a + i\frac{b-a}{N}, \quad \text{para } i=0, 1, \dots, N.
\]
Consideremos el $i$-ésimo subintervalo $I_i = [x_{i-1}, x_i]$ para $i=1, \dots, N$.
Dado que $f$ es continua, alcanza sus valores máximo $M_i$ y mínimo $m_i$ en $I_i$ por el Teorema de Weierstrass. Es decir, existen $c_i, d_i \in I_i$ tales que $M_i = f(c_i)$ y $m_i = f(d_i)$.
Como ambos puntos $c_i, d_i$ pertenecen a $I_i$, su distancia es a lo sumo la longitud del intervalo:
\[
|c_i - d_i| \leq |x_i - x_{i-1}| = \Delta x < \delta.
\]
Aplicando la continuidad uniforme, inferimos que la oscilación de $f$ en el intervalo $I_i$ está acotada:
\[
M_i - m_i = |f(c_i) - f(d_i)| < \frac{\varepsilon}{2(b-a)}.
\]
Ahora, cubrimos la porción del grafo $G(f)$ correspondiente al intervalo $I_i$ con un rectángulo $R_i$. Definimos:
\[
R_i = I_i \times [m_i, M_i] = [x_{i-1}, x_i] \times [m_i, M_i].
\]
Es evidente que para todo $x \in I_i$, $(x, f(x)) \in R_i$. Luego, la unión finita de estos $N$ rectángulos cubre completamente al grafo de $f$:
\[
G(f) \subset \bigcup_{i=1}^N R_i.
\]
El área de cada rectángulo $R_i$ es:
\[
\operatorname{area}(R_i) = (x_i - x_{i-1})(M_i - m_i) = \Delta x (M_i - m_i).
\]
Calculamos la suma total de las áreas del recubrimiento:
\begin{align*}
\sum_{i=1}^N \operatorname{area}(R_i) &= \sum_{i=1}^N \Delta x (M_i - m_i) \\
&< \sum_{i=1}^N \Delta x \left( \frac{\varepsilon}{2(b-a)} \right) \\
&= \frac{\varepsilon}{2(b-a)} \sum_{i=1}^N \Delta x.
\end{align*}
Observamos que la suma de las longitudes de los subintervalos es precisamente la longitud total del intervalo original:
\[
\sum_{i=1}^N \Delta x = \sum_{i=1}^N (x_i - x_{i-1}) = x_N - x_0 = b - a.
\]
Sustituyendo esto en la desigualdad:
\[
\sum_{i=1}^N \operatorname{area}(R_i) < \frac{\varepsilon}{2(b-a)} \cdot (b - a) = \frac{\varepsilon}{2} < \varepsilon.
\]
Hemos encontrado un recubrimiento finito mediante rectángulos cerrados cuya suma de áreas es estrictamente menor que $\varepsilon$. Dado que $\varepsilon > 0$ fue elegido arbitrariamente, se concluye formalmente que $G(f)$ tiene contenido nulo. \qed

\begin{ejerciciobox}{Fronteras de regiones elementales}
Con base en el teorema anterior, se puede demostrar que toda región elemental en $\mathbb{R}^2$ de Tipo 1 o Tipo 2 tiene una frontera de contenido nulo. 
En efecto, si $D = \{ (x,y) \mid a \leq x \leq b, \phi_1(x) \leq y \leq \phi_2(x) \}$, la frontera $\partial D$ está compuesta por las gráficas de las funciones continuas $\phi_1, \phi_2$ en $[a,b]$ y, posiblemente, por segmentos verticales en $x=a$ y $x=b$.
Los segmentos verticales son degenerados y claramente tienen contenido nulo. Por el Teorema demostrado, los grafos de $\phi_1$ y $\phi_2$ también tienen contenido nulo. Dado que la unión finita de conjuntos de contenido nulo tiene contenido nulo, $\partial D$ tiene contenido nulo. 
\textit{Ejercicio para el lector:} Proveer los detalles algebraicos para acotar los segmentos de recta verticales y formalizar la unión de cubrimientos.
\end{ejerciciobox}
\part*{Parte II: Integral en regiones generales y Teorema de Fubini}

\section{Integral Doble sobre Regiones Generales}

En la Parte I se estableció la teoría de integración de Darboux y Riemann para funciones acotadas definidas en rectángulos cerrados $R \subset \mathbb{R}^2$. Sin embargo, en las aplicaciones geométricas y físicas (como el cálculo de centros de masa, momentos de inercia o probabilidades), los dominios de integración raramente son rectangulares. A continuación, extenderemos la noción de integral doble a regiones acotadas más generales.

\subsection{Definición mediante extensión por cero}

Sea $D \subset \mathbb{R}^2$ un conjunto acotado y sea $f: D \to \mathbb{R}$ una función acotada. Dado que $D$ es acotado, existe algún rectángulo cerrado $R = [a,b] \times [c,d]$ tal que $D \subset R$. 

\begin{definicionbox}{Extensión por cero y Función Característica}{}
    La \textbf{función característica} (o indicatriz) del conjunto $D$, denotada por $\chi_D : \mathbb{R}^2 \to \mathbb{R}$, se define como:
    \[
    \chi_D(x,y) = \begin{cases} 
    1, & \text{si } (x,y) \in D \\
    0, & \text{si } (x,y) \notin D 
    \end{cases}
    \]
    Definimos la \textbf{extensión por cero} de $f$ al rectángulo $R$, denotada como $f^* : R \to \mathbb{R}$, mediante el producto $f^*(x,y) = f(x,y)\chi_D(x,y)$. Explícitamente:
    \[
    f^*(x,y) = \begin{cases} 
    f(x,y), & \text{si } (x,y) \in D \\
    0, & \text{si } (x,y) \in R \setminus D 
    \end{cases}
    \]
\end{definicionbox}

A partir de esta extensión, la definición natural de integral sobre $D$ se traslada a la integral de $f^*$ sobre el rectángulo $R$.

\begin{definicionbox}{Integral sobre una región general}{}
    Sea $D \subset \mathbb{R}^2$ un conjunto acotado y $f: D \to \mathbb{R}$ una función acotada. Se dice que \textbf{$f$ es integrable sobre $D$} si la extensión por cero $f^*$ es integrable sobre un rectángulo cerrado $R$ que contenga a $D$. En tal caso, la \textbf{integral doble de $f$ sobre $D$} se define como:
    \[
    \iint_D f(x,y) \, dA = \iint_R f^*(x,y) \, dA = \iint_R f(x,y)\chi_D(x,y) \, d(x,y)
    \]
\end{definicionbox}

\begin{observacionbox}{Independencia del Rectángulo}
    Es imperativo demostrar que esta definición \textit{no depende} de la elección del rectángulo $R$. Supongamos que $R_1$ y $R_2$ son dos rectángulos tales que $D \subset R_1$ y $D \subset R_2$. Consideremos $R_3 = R_1 \cap R_2$. Al evaluar las sumas superiores e inferiores sobre particiones, se demuestra rigurosamente que $\iint_{R_1} f^* dA = \iint_{R_3} f^* dA = \iint_{R_2} f^* dA$, garantizando la buena definición del concepto.
\end{observacionbox}

\subsection{Condiciones de Integrabilidad}

Para que la función extendida $f^*$ sea integrable en $R$, las discontinuidades de $f^*$ deben tener medida cero (por el Teorema de Lebesgue). Notemos que el conjunto de discontinuidades de $f^*$ está contenido en la unión de dos conjuntos: el conjunto de discontinuidades de $f$ dentro del interior de $D$, y la frontera de $D$ (denotada $\partial D$).

\begin{teoremabox}{Integrabilidad sobre regiones con frontera de contenido nulo}{}
    Sea $D \subset \mathbb{R}^2$ un conjunto acotado cuya frontera $\partial D$ tiene contenido nulo (y por tanto, medida nula). Si $f: D \to \mathbb{R}$ es una función acotada que es continua en $D$ (o continua salvo en un conjunto de medida nula), entonces $f$ es integrable sobre $D$.
\end{teoremabox}

\textbf{Demostración:}
Sea $R$ un rectángulo cerrado tal que $D \subset R$. La función $f^*$ (extensión por cero) puede presentar discontinuidades únicamente en:
1. Puntos de $D$ donde $f$ es discontinua (conjunto de medida nula, por hipótesis).
2. Puntos de la frontera $\partial D$.
Puesto que $\partial D$ tiene contenido nulo, es un conjunto de medida nula. Por lo tanto, el conjunto total de discontinuidades de $f^*$ en $R$ es la unión de dos conjuntos de medida nula, resultando en un conjunto de medida nula. Como $f$ es acotada en $D$, $f^*$ es acotada en $R$. Por el criterio de Lebesgue, $f^*$ es integrable en $R$, concluyendo la demostración. $\blacksquare$

\begin{errorbox}{Fronteras Patológicas}
    Es un error común asumir que toda región acotada $D$ permite la integración de funciones continuas. Si $\partial D$ tiene medida de Lebesgue positiva (como en el caso de ciertos conjuntos de tipo Cantor de medida positiva extendidos al plano), la función constante $f(x,y) = 1$ no será integrable en sentido de Riemann sobre $D$, puesto que $\chi_D$ sería discontinua en un conjunto de medida no nula. La hipótesis del contenido nulo de la frontera es estructural.
\end{errorbox}

\section{Área de una Región Plana}

La medida bidimensional en la teoría de integración de Riemann es el \textit{área} o \textit{medida de Jordan}.

\begin{definicionbox}{Área de una región}{}
    Sea $D \subset \mathbb{R}^2$ un conjunto acotado. Si la función constante $f(x,y) = 1$ es integrable sobre $D$, se dice que $D$ es \textbf{medible en el sentido de Jordan}. El \textbf{área} de $D$, denotada por $A(D)$ o $\mu(D)$, se define como:
    \[
    A(D) = \iint_D 1 \, dA = \iint_D dA
    \]
\end{definicionbox}

\section{Teorema de Fubini en Regiones Generales}

El teorema de Fubini para rectángulos puede adaptarse a regiones de formas especiales, denominadas regiones proyectables o de Tipo I y Tipo II. Estas regiones forman la base del cálculo práctico de integrales múltiples mediante integrales iteradas.

\subsection{Regiones Tipo I y Tipo II}

\begin{definicionbox}{Región Tipo I}{}
    Una región plana $D$ es de \textbf{Tipo I} (proyectable sobre el eje $x$) si puede describirse como:
    \[
    D = \{ (x,y) \in \mathbb{R}^2 \mid a \leq x \leq b, \; \phi_1(x) \leq y \leq \phi_2(x) \}
    \]
    donde $\phi_1, \phi_2 : [a,b] \to \mathbb{R}$ son funciones continuas.
\end{definicionbox}

\begin{definicionbox}{Región Tipo II}{}
    Una región plana $D$ es de \textbf{Tipo II} (proyectable sobre el eje $y$) si puede describirse como:
    \[
    D = \{ (x,y) \in \mathbb{R}^2 \mid c \leq y \leq d, \; \psi_1(y) \leq x \leq \psi_2(y) \}
    \]
    donde $\psi_1, \psi_2 : [c,d] \to \mathbb{R}$ son funciones continuas.
\end{definicionbox}

\begin{teoremabox}{Teorema de Fubini para Regiones Proyectables}{}
    Sea $f: D \to \mathbb{R}$ una función continua sobre una región $D$.
    \begin{itemize}
        \item[(i)] Si $D$ es una región de \textbf{Tipo I}, entonces:
        \[
        \iint_D f(x,y) \, dA = \int_a^b \left( \int_{\phi_1(x)}^{\phi_2(x)} f(x,y) \, dy \right) dx
        \]
        \item[(ii)] Si $D$ es una región de \textbf{Tipo II}, entonces:
        \[
        \iint_D f(x,y) \, dA = \int_c^d \left( \int_{\psi_1(y)}^{\psi_2(y)} f(x,y) \, dx \right) dy
        \]
    \end{itemize}
\end{teoremabox}

\textbf{Demostración (Caso Tipo I):}
Como $D$ es compacta (cerrada y acotada) y $f$ es continua, $f$ es integrable sobre $D$. Además, la frontera de $D$ consta de las gráficas de las funciones continuas $\phi_1$ y $\phi_2$, junto con los segmentos verticales $x=a$ y $x=b$. Por teoremas previos, la gráfica de una función continua en un intervalo cerrado tiene contenido nulo. Por ende, $\partial D$ tiene contenido nulo.
Sea $R = [a,b] \times [c,d]$ un rectángulo que contiene a $D$, escogiendo $c = \min_{x \in [a,b]} \phi_1(x)$ y $d = \max_{x \in [a,b]} \phi_2(x)$.
Definamos la extensión $f^*(x,y) = f(x,y)$ si $(x,y) \in D$, y $f^*(x,y) = 0$ en $R \setminus D$.
Por definición:
\[
\iint_D f(x,y) \, dA = \iint_R f^*(x,y) \, dA
\]
Aplicando el Teorema de Fubini clásico sobre el rectángulo $R$:
\[
\iint_R f^*(x,y) \, dA = \int_a^b \left( \int_c^d f^*(x,y) \, dy \right) dx
\]
Fijemos $x \in [a,b]$. La integral interior es sobre la variable $y$ en el intervalo $[c,d]$. Observamos la definición de $f^*(x,y)$:
\[
f^*(x,y) = \begin{cases} 
0, & c \leq y < \phi_1(x) \\
f(x,y), & \phi_1(x) \leq y \leq \phi_2(x) \\
0, & \phi_2(x) < y \leq d 
\end{cases}
\]
Por las propiedades de aditividad del dominio para la integral de Riemann en una variable:
\begin{align*}
\int_c^d f^*(x,y) \, dy &= \int_c^{\phi_1(x)} 0 \, dy + \int_{\phi_1(x)}^{\phi_2(x)} f(x,y) \, dy + \int_{\phi_2(x)}^d 0 \, dy \\
&= \int_{\phi_1(x)}^{\phi_2(x)} f(x,y) \, dy
\end{align*}
Sustituyendo esto en la integral exterior:
\[
\iint_D f(x,y) \, dA = \int_a^b \left( \int_{\phi_1(x)}^{\phi_2(x)} f(x,y) \, dy \right) dx
\]
Lo cual concluye la demostración rigurosa. El caso Tipo II es completamente análogo por simetría de variables. $\blacksquare$

\section{Cambio del Orden de Integración}

Cuando una región es proyectable en ambos ejes simultáneamente (es de Tipo I y Tipo II al mismo tiempo), la integral doble puede evaluarse mediante iteradas en cualquier orden: $dx dy$ o $dy dx$. A menudo, el cálculo analítico de una integral iterada puede ser imposible en un orden determinado debido a la falta de antiderivadas expresables en funciones elementales, volviéndose trivial tras intercambiar el orden.

Para cambiar el orden:
1. \textbf{Esbozar la región $D$} a partir de los límites de integración iterados.
2. \textbf{Redescribir $D$} intercambiando los roles de variable dependiente e independiente (de Tipo I a Tipo II o viceversa).
3. \textbf{Plantear la nueva integral} usando la nueva descripción.

\begin{ejemplobox}{Cambio de orden: $e^{y^2}$}{}
    Evalúe la integral iterada:
    \[
    I = \int_0^1 \int_x^1 e^{y^2} \, dy dx
    \]
\end{ejemplobox}

\textbf{Solución:}
La integral interior es $\int e^{y^2} dy$. Esta antiderivada no es una función elemental (está relacionada con la función error de Gauss), lo que bloquea el cálculo directo.
Identifiquemos la región de integración $D$. Los límites dictan que $D$ es una región Tipo I:
\[
0 \leq x \leq 1 \quad \text{y} \quad x \leq y \leq 1
\]
Esto es, $D = \{ (x,y) \in \mathbb{R}^2 \mid 0 \leq x \leq 1, \; x \leq y \leq 1 \}$.
Esta región es el triángulo delimitado por las rectas $x=0$ (eje $y$), $y=1$ y $y=x$.

\begin{center}
\begin{tikzpicture}
\begin{axis}[
    axis lines=middle,
    xlabel={$x$},
    ylabel={$y$},
    xmin=-0.2, xmax=1.5,
    ymin=-0.2, ymax=1.5,
    axis equal,
    enlargelimits,
    domain=0:1,
    samples=100
]
\addplot [fill=blue!20, draw=none] coordinates {(0,0) (1,1) (0,1)} \closedcycle;
\addplot [thick, black] {x};
\addplot [thick, black] coordinates {(0,1) (1,1)};
\addplot [thick, black] coordinates {(0,0) (0,1)};
\node at (axis cs: 0.3, 0.7) {$D$};
\node[right] at (axis cs: 1, 1) {$y = x$};
\node[above] at (axis cs: 0.5, 1) {$y = 1$};
\end{axis}
\end{tikzpicture}
\end{center}

Ahora, describimos $D$ como una región de Tipo II, proyectándola sobre el eje $y$. En el eje $y$, la región varía desde $y=0$ hasta $y=1$. Para un $y$ fijo en ese intervalo, la coordenada $x$ varía desde la recta vertical $x=0$ hasta la recta oblicua $x=y$. Por tanto:
\[
D = \{ (x,y) \in \mathbb{R}^2 \mid 0 \leq y \leq 1, \; 0 \leq x \leq y \}
\]
Reescribiendo la integral con este nuevo orden ($dx$ seguido de $dy$):
\[
I = \int_0^1 \int_0^y e^{y^2} \, dx dy
\]
Ahora evaluamos. La integral interior respecto a $x$ trata a $y$ como constante:
\[
\int_0^y e^{y^2} \, dx = e^{y^2} \left[ x \right]_0^y = y e^{y^2}
\]
Sustituyendo en la integral exterior:
\[
I = \int_0^1 y e^{y^2} \, dy
\]
Esta integral se resuelve trivialmente mediante la sustitución $u = y^2$, de donde $du = 2y \, dy \implies y \, dy = \frac{du}{2}$. Los límites cambian a $0^2=0$ y $1^2=1$:
\[
I = \int_0^1 \frac{1}{2} e^u \, du = \frac{1}{2} \left[ e^u \right]_0^1 = \frac{1}{2} (e - 1)
\]
El cambio de orden transformó un problema irresoluble en uno elemental. $\blacksquare$

\begin{ejerciciobox}{Problema Tipo UNI - Valor absoluto}{}
    Calcule el valor de la integral:
    \[
    \iint_D |x - y| \, dA
    \]
    donde $D = \{ (x,y) \in \mathbb{R}^2 \mid 0 \leq x \leq 2, \; 0 \leq y \leq 2 \}$.
\end{ejerciciobox}

\textbf{Resolución rigurosa paso a paso:}
El integrando contiene un valor absoluto. La estrategia fundamental en Cálculo Avanzado para tratar valores absolutos dentro de integrales múltiples es particionar el dominio $D$ en subregiones donde el argumento del valor absoluto no cambie de signo.
Sabemos que:
\[
|x - y| = \begin{cases} 
x - y, & \text{si } x \geq y \\
y - x, & \text{si } y \geq x 
\end{cases}
\]
Definimos las subregiones $D_1$ y $D_2$ dividiendo el cuadrado $D = [0,2] \times [0,2]$ por la recta $y=x$:
\begin{align*}
D_1 &= \{ (x,y) \in D \mid y \leq x \} \\
D_2 &= \{ (x,y) \in D \mid y \geq x \}
\end{align*}
Geométricamente, la recta $y=x$ es la diagonal principal del cuadrado $D$. $D_1$ es el triángulo inferior y $D_2$ es el triángulo superior. Por aditividad del dominio sobre conjuntos con intersección de medida nula (el segmento sobre $y=x$):
\[
\iint_D |x - y| \, dA = \iint_{D_1} (x - y) \, dA + \iint_{D_2} (y - x) \, dA
\]

\textbf{Evaluación sobre $D_1$:}
Describimos $D_1$ como región Tipo I: $0 \leq x \leq 2$, y para $x$ fijo, $0 \leq y \leq x$.
\begin{align*}
I_1 &= \iint_{D_1} (x - y) \, dA = \int_0^2 \int_0^x (x - y) \, dy dx \\
    &= \int_0^2 \left[ xy - \frac{y^2}{2} \right]_{y=0}^{y=x} dx \\
    &= \int_0^2 \left( x(x) - \frac{x^2}{2} - (0 - 0) \right) dx \\
    &= \int_0^2 \left( x^2 - \frac{x^2}{2} \right) dx \\
    &= \int_0^2 \frac{x^2}{2} \, dx \\
    &= \frac{1}{2} \left[ \frac{x^3}{3} \right]_0^2 = \frac{1}{2} \left( \frac{8}{3} \right) = \frac{4}{3}
\end{align*}

\textbf{Evaluación sobre $D_2$:}
Por simetría pura del problema, si intercambiamos las variables $x$ e $y$, el dominio $D$ permanece invariante (es un cuadrado $0\le x \le 2, 0\le y \le 2$), el subdominio $D_2$ se mapea sobre $D_1$, y el integrando $(y-x)$ se transforma en $(x-y)$. Es inmediato predecir que $I_2 = I_1 = \frac{4}{3}$. No obstante, procedemos con el cálculo explícito por completitud:
Describimos $D_2$ como Tipo I: $0 \leq x \leq 2$, y para $x$ fijo, $x \leq y \leq 2$.
\begin{align*}
I_2 &= \iint_{D_2} (y - x) \, dA = \int_0^2 \int_x^2 (y - x) \, dy dx \\
    &= \int_0^2 \left[ \frac{y^2}{2} - xy \right]_{y=x}^{y=2} dx \\
    &= \int_0^2 \left( \left( \frac{2^2}{2} - 2x \right) - \left( \frac{x^2}{2} - x^2 \right) \right) dx \\
    &= \int_0^2 \left( 2 - 2x - \left( -\frac{x^2}{2} \right) \right) dx \\
    &= \int_0^2 \left( \frac{x^2}{2} - 2x + 2 \right) dx \\
    &= \left[ \frac{x^3}{6} - x^2 + 2x \right]_0^2 \\
    &= \left( \frac{8}{6} - 4 + 4 \right) - (0) = \frac{4}{3}
\end{align*}

Finalmente, sumamos los aportes de las dos regiones:
\[
\iint_D |x - y| \, dA = I_1 + I_2 = \frac{4}{3} + \frac{4}{3} = \frac{8}{3}
\]
$\blacksquare$

\begin{ejerciciobox}{Cálculo Teórico - Identidad de Doble Integral}{}
    Sea $F$ la antiderivada de una función continua $f: [0,1] \to \mathbb{R}$, esto es, $F(x) = \int_0^x f(t) dt$. Demuestre usando integrales dobles que:
    \[
    \int_0^1 \int_0^x F(y) \, dy dx = \int_0^1 \frac{(1-x)^2}{2} f(x) \, dx
    \]
\end{ejerciciobox}

\textbf{Demostración:}
Denotemos el lado izquierdo como $J$:
\[
J = \int_0^1 \int_0^x F(y) \, dy dx
\]
La región de integración es $D = \{ (x,y) \mid 0 \leq x \leq 1, \; 0 \leq y \leq x \}$.
Aplicamos el Teorema de Fubini para cambiar el orden de integración, escribiendo $D$ como región de Tipo II:
\[
D = \{ (x,y) \mid 0 \leq y \leq 1, \; y \leq x \leq 1 \}
\]
Entonces:
\[
J = \int_0^1 \int_y^1 F(y) \, dx dy = \int_0^1 F(y) \left( \int_y^1 dx \right) dy = \int_0^1 F(y)(1 - y) \, dy
\]
Aplicamos ahora integración por partes en la integral de una variable resultante:
Sea $u = F(y)$ y $dv = (1-y) dy$.
Por el Teorema Fundamental del Cálculo (puesto que $f$ es continua), $du = f(y) dy$.
Elegimos la antiderivada de $dv$ como $v = -\frac{(1-y)^2}{2}$. (Nota: elegimos esta constante de integración específica para simplificar la evaluación en los límites).
Aplicando la fórmula $\int_a^b u dv = \left[uv\right]_a^b - \int_a^b v du$:
\begin{align*}
J &= \left[ F(y) \left( -\frac{(1-y)^2}{2} \right) \right]_{y=0}^{y=1} - \int_0^1 \left( -\frac{(1-y)^2}{2} \right) f(y) \, dy \\
  &= \left( F(1)(0) - F(0)\left(-\frac{1}{2}\right) \right) + \int_0^1 \frac{(1-y)^2}{2} f(y) \, dy
\end{align*}
Dado que $F(x) = \int_0^x f(t)dt$, entonces $F(0) = \int_0^0 f(t)dt = 0$. El término evaluado se anula completamente:
\[
J = 0 + \int_0^1 \frac{(1-y)^2}{2} f(y) \, dy
\]
Dado que la variable de integración es muda, podemos renombrar $y$ como $x$:
\[
J = \int_0^1 \frac{(1-x)^2}{2} f(x) \, dx
\]
Queda demostrada la identidad teórica. $\blacksquare$

\vspace{0.5cm}
Con estas herramientas de cambio de orden y reducción a integrales iteradas en dominios generales, el estudiante está preparado para abordar el estudio de los cambios de variable (como coordenadas polares) mediante el Teorema de Transformación (Jacobiano), que será objeto de la Parte III.
\part*{Parte III: Fubini en sólidos y cálculo de volúmenes}

\section{Teorema de Fubini para Integrales Triples}

\begin{definicionbox}{Sólidos Simples o Regiones de Integración}
Una región sólida $E \subset \mathbb{R}^3$ se dice que es de \textbf{Tipo 1} si se encuentra entre las gráficas de dos funciones continuas de $x$ e $y$, es decir,
\[ E = \{(x,y,z) \in \mathbb{R}^3 : (x,y) \in D, \quad u_1(x,y) \leq z \leq u_2(x,y)\} \]
donde $D$ es la proyección de $E$ sobre el plano $xy$. De forma análoga:
\begin{itemize}
    \item \textbf{Tipo 2:} La región se proyecta en el plano $yz$. $E = \{(x,y,z) : (y,z) \in D, \quad u_1(y,z) \leq x \leq u_2(y,z)\}$.
    \item \textbf{Tipo 3:} La región se proyecta en el plano $xz$. $E = \{(x,y,z) : (x,z) \in D, \quad u_1(x,z) \leq y \leq u_2(x,z)\}$.
\end{itemize}
\end{definicionbox}

\begin{teoremabox}{Teorema de Fubini para Integrales Triples}
Si $f$ es continua en una región sólida $E$ de Tipo 1, entonces
\[ \iiint_E f(x,y,z) \,dV = \iint_D \left[ \int_{u_1(x,y)}^{u_2(x,y)} f(x,y,z) \,dz \right] dA \]
donde $D$ es la proyección de $E$ sobre el plano $xy$. La integral doble sobre $D$ se puede evaluar iteradamente dependiendo de si $D$ es una región plana de Tipo I o Tipo II, lo que da lugar a los órdenes de integración $dz\,dy\,dx$ y $dz\,dx\,dy$.
\end{teoremabox}

\begin{observacionbox}{Los seis órdenes de integración}
Dado un sólido general, existen $3! = 6$ posibles órdenes de integración: $dz\,dy\,dx$, $dz\,dx\,dy$, $dx\,dy\,dz$, $dx\,dz\,dy$, $dy\,dx\,dz$, y $dy\,dz\,dx$. El cambio en el orden de integración en integrales triples requiere proyectar el sólido $E$ sobre uno de los tres planos coordenados y establecer las superficies superior e inferior (o derecha/izquierda, anterior/posterior) que delimitan al sólido a partir de esa proyección.
\end{observacionbox}

\begin{center}
\begin{tikzpicture}[scale=1.5]
    % Axes
    \draw[->] (0,0,0) -- (2.5,0,0) node[right]{$y$};
    \draw[->] (0,0,0) -- (0,2.5,0) node[above]{$z$};
    \draw[->] (0,0,0) -- (0,0,3) node[below left]{$x$};
    
    % Projection D on xy-plane (z=0)
    \draw[fill=gray!20, opacity=0.8] (1,0,1) .. controls (2,0,0.5) and (2,0,2) .. (0.5,0,2.5) .. controls (0,0,2) and (0.5,0,1.5) .. (1,0,1);
    \node at (1.2,0,1.6) {$D$};
    
    % Solid E
    \draw[fill=blue!10, opacity=0.6] (1,1,1) .. controls (2,1.2,0.5) and (2,1.5,2) .. (0.5,1.1,2.5) .. controls (0,1.3,2) and (0.5,1.4,1.5) .. (1,1,1);
    \draw[fill=blue!30, opacity=0.6] (1,2,1) .. controls (2,2.2,0.5) and (2,2.5,2) .. (0.5,2.1,2.5) .. controls (0,2.3,2) and (0.5,2.4,1.5) .. (1,2,1);
    
    % Vertical lines
    \draw[dashed] (1,0,1) -- (1,1,1);
    \draw[dashed] (2,0,1.25) -- (2,1.35,1.25);
    \draw[dashed] (0.5,0,2.5) -- (0.5,1.1,2.5);
    
    \node at (1.2,2.3,1.6) {$z=u_2(x,y)$};
    \node at (1.2,0.8,1.6) {$z=u_1(x,y)$};
    \node at (2.2, 1.5, 1) {$E$};
\end{tikzpicture}
\end{center}

\begin{errorbox}{Error al cambiar el orden de integración}
Un error sumamente frecuente al cambiar el orden de integración en integrales triples es intercambiar directamente los límites de las integrales iteradas sin reevaluar las proyecciones. Es estrictamente necesario graficar (o imaginar) el sólido $E$, determinar su nueva proyección bidimensional y redefinir las funciones que lo acotan en la nueva dirección geométrica.
\end{errorbox}

\section{Cálculo de Volúmenes usando Coordenadas Cartesianas}
El volumen de un sólido $E$ se obtiene integrando la función constante $f(x,y,z) = 1$:
\[ V(E) = \iiint_E \,dV \]

\begin{ejemplobox}{Volumen del sólido de Steinmetz}
\textbf{Problema:} Determine el volumen del sólido $E$ limitado por la intersección de dos cilindros circulares rectos de radio $r$ cuyos ejes se intersecan perpendicularmente.

\textbf{Solución:} 
Sin pérdida de generalidad, supongamos que los cilindros tienen como ejes el eje $x$ y el eje $y$. Sus ecuaciones son:
\[ x^2 + z^2 = r^2 \quad \text{y} \quad y^2 + z^2 = r^2 \]
El sólido $E$ exhibe simetría en los $8$ octantes. Calcularemos el volumen en el primer octante y lo multiplicaremos por $8$. En el primer octante, $x, y, z \geq 0$.
De las ecuaciones, vemos que el sólido está limitado superiormente por el techo que proviene de la intersección geométrica. En realidad, para un valor fijo de $z$, las restricciones para $x$ e $y$ son:
\[ 0 \leq x \leq \sqrt{r^2 - z^2}, \quad 0 \leq y \leq \sqrt{r^2 - z^2} \]
Podemos proyectar sobre el plano $xy$, pero resulta más sencillo integrar con el orden $dx\,dy\,dz$, tomando secciones transversales a una altura $z$. La sección a altura $z$ es un cuadrado de lado $2\sqrt{r^2-z^2}$.
El volumen de un octante es:
\[ V_1 = \int_0^r \int_0^{\sqrt{r^2-z^2}} \int_0^{\sqrt{r^2-z^2}} dx\,dy\,dz \]
Resolviendo la integral más interna respecto a $x$:
\[ \int_0^{\sqrt{r^2-z^2}} 1 \,dx = \sqrt{r^2-z^2} \]
Luego, integrando respecto a $y$:
\[ \int_0^{\sqrt{r^2-z^2}} \sqrt{r^2-z^2} \,dy = \sqrt{r^2-z^2} \left[ y \right]_0^{\sqrt{r^2-z^2}} = (\sqrt{r^2-z^2})^2 = r^2 - z^2 \]
Finalmente, respecto a $z$:
\[ V_1 = \int_0^r (r^2 - z^2) \,dz = \left[ r^2 z - \frac{z^3}{3} \right]_0^r = r^3 - \frac{r^3}{3} = \frac{2r^3}{3} \]
El volumen total es $8$ veces este valor:
\[ V = 8 \left(\frac{2r^3}{3}\right) = \frac{16r^3}{3} \]
Esta es una demostración clásica exigida en el curso.
\end{ejemplobox}

\begin{ejemplobox}{Tetraedro y sólido entre paraboloides}
\textbf{Tetraedro:} Considere el tetraedro limitado por los planos coordenados y el plano $x+y+z=1$. Proyectando sobre $z=0$, obtenemos el triángulo $D$ limitado por $x=0$, $y=0$, $x+y=1$. El volumen es:
\[ V = \int_0^1 \int_0^{1-x} \int_0^{1-x-y} dz\,dy\,dx = \int_0^1 \int_0^{1-x} (1-x-y) \,dy\,dx = \int_0^1 \frac{(1-x)^2}{2} \,dx = \frac{1}{6} \]

\textbf{Sólido entre paraboloides:} Considere $E$ acotado por $z = x^2+y^2$ y $z = 4 - 3x^2 - 3y^2$. La intersección ocurre cuando $x^2+y^2 = 4 - 3x^2 - 3y^2 \implies 4(x^2+y^2) = 4 \implies x^2+y^2 = 1$. La proyección $D$ es el disco unidad.
\[ V = \iint_D \left( \int_{x^2+y^2}^{4-3x^2-3y^2} dz \right) dA = \iint_D (4 - 4x^2 - 4y^2) dA \]
\end{ejemplobox}

\section{Aplicaciones: Masa, Centro de Masa y Momentos de Inercia}

\begin{definicionbox}{Distribuciones de Masa}
Si una región sólida $E$ tiene una densidad de masa volumétrica $\rho(x,y,z)$ (masa por unidad de volumen), se definen las siguientes cantidades físicas:
\begin{itemize}
    \item \textbf{Masa total:} $m = \iiint_E \rho(x,y,z) \,dV$.
    \item \textbf{Momentos estáticos:} Respecto a los planos coordenados:
    \[ M_{yz} = \iiint_E x \rho(x,y,z) \,dV, \quad M_{xz} = \iiint_E y \rho(x,y,z) \,dV, \quad M_{xy} = \iiint_E z \rho(x,y,z) \,dV \]
    \item \textbf{Centro de masa:} El punto $(\bar{x}, \bar{y}, \bar{z})$ tal que $\bar{x} = \frac{M_{yz}}{m}$, $\bar{y} = \frac{M_{xz}}{m}$, $\bar{z} = \frac{M_{xy}}{m}$.
    \item \textbf{Momentos de inercia:} Respecto a los ejes coordenados (miden la inercia rotacional):
    \[ I_x = \iiint_E (y^2+z^2) \rho(x,y,z) \,dV, \quad I_y = \iiint_E (x^2+z^2) \rho(x,y,z) \,dV, \quad I_z = \iiint_E (x^2+y^2) \rho(x,y,z) \,dV \]
\end{itemize}
\end{definicionbox}

\begin{ejerciciobox}{Masa de un adorno sólido (Problema tipo PC pasada)}
\textbf{Problema:} Se desea fabricar un adorno macizo cuya forma es la del sólido $E$ limitado inferiormente por el plano $z=0$ y superiormente por el paraboloide $z = 5 - x^2 - y^2$. La densidad volumétrica del material está dada por $\rho(x,y,z) = |x|$. Calcule la masa total del adorno utilizando integrales iteradas de Fubini en coordenadas cartesianas estrictamente.

\textbf{Solución:}
La masa $m$ está dada por:
\[ m = \iiint_E |x| \,dV \]
El sólido $E$ intersecta al plano $z=0$ en $0 = 5 - x^2 - y^2 \implies x^2 + y^2 = 5$. La proyección de $E$ sobre el plano $xy$ es el disco $D = \{(x,y) : x^2 + y^2 \leq 5\}$.
Dado que la densidad es $\rho(x,y,z) = |x|$, es una función par respecto a $x$. El dominio de integración es simétrico respecto al plano $yz$ (es decir, el dominio en el plano $xy$ es simétrico respecto al eje $y$).
Podemos calcular la masa en la mitad derecha ($x \geq 0$) y multiplicarla por 2:
\[ m = 2 \iiint_{E \cap \{x \geq 0\}} x \,dV \]
El orden de integración cartesiano será $dz\,dy\,dx$. Proyectamos el medio disco $x \geq 0$ sobre el plano $xy$. Los límites para $x$ van de $0$ a $\sqrt{5}$. Para un $x$ fijo, $y$ varía de $-\sqrt{5-x^2}$ a $\sqrt{5-x^2}$. Para $(x,y)$ fijos, $z$ varía desde $0$ hasta el techo del paraboloide $5-x^2-y^2$.
\[ m = 2 \int_{0}^{\sqrt{5}} \int_{-\sqrt{5-x^2}}^{\sqrt{5-x^2}} \int_{0}^{5-x^2-y^2} x \,dz\,dy\,dx \]
Evaluando la integral interna respecto a $z$:
\[ \int_{0}^{5-x^2-y^2} x \,dz = x \left[ z \right]_0^{5-x^2-y^2} = x(5-x^2-y^2) \]
Sustituyendo en la integral doble:
\[ m = 2 \int_{0}^{\sqrt{5}} \int_{-\sqrt{5-x^2}}^{\sqrt{5-x^2}} (5x - x^3 - xy^2) \,dy\,dx \]
Integrando respecto a $y$:
\begin{align*}
\int_{-\sqrt{5-x^2}}^{\sqrt{5-x^2}} (5x - x^3 - xy^2) \,dy &= \left[ (5x-x^3)y - x \frac{y^3}{3} \right]_{-\sqrt{5-x^2}}^{\sqrt{5-x^2}} \\
&= 2 \left( (5x-x^3)\sqrt{5-x^2} - \frac{x}{3}(\sqrt{5-x^2})^3 \right)
\end{align*}
Notemos que $5x-x^3 = x(5-x^2)$, entonces el primer término es $x(5-x^2)^{3/2}$.
El segundo término es $\frac{x}{3}(5-x^2)^{3/2}$.
La resta algebraica resulta en:
\[ 2 \left( x(5-x^2)^{3/2} - \frac{1}{3}x(5-x^2)^{3/2} \right) = 2 \left( \frac{2}{3}x(5-x^2)^{3/2} \right) = \frac{4}{3}x(5-x^2)^{3/2} \]
Ahora sustituimos esto en la integral externa respecto a $x$:
\[ m = 2 \int_{0}^{\sqrt{5}} \frac{4}{3}x(5-x^2)^{3/2} \,dx = \frac{8}{3} \int_{0}^{\sqrt{5}} x(5-x^2)^{3/2} \,dx \]
Hacemos el cambio de variable $u = 5-x^2$, de lo cual se deduce que $du = -2x \,dx \implies x \,dx = -\frac{1}{2} \,du$.
Los nuevos límites de integración son: si $x=0$, $u=5$; si $x=\sqrt{5}$, $u=0$.
\[ m = \frac{8}{3} \int_{5}^{0} (u^{3/2}) \left(-\frac{1}{2} du\right) = \frac{4}{3} \int_{0}^{5} u^{3/2} \,du \]
Evaluamos la integral polinómica final:
\[ m = \frac{4}{3} \left[ \frac{2}{5} u^{5/2} \right]_0^5 = \frac{8}{15} (5)^{5/2} = \frac{8}{15} (25\sqrt{5}) = \frac{200\sqrt{5}}{15} = \frac{40\sqrt{5}}{3} \]
Por lo tanto, la masa total del adorno sólido es de $\frac{40\sqrt{5}}{3}$ unidades de masa. 
(Se observan todos los pasos de la integración iterada directa).
\end{ejerciciobox}
\part*{Parte IV: Sistemas de Coordenadas Polares, Cilíndricas y Esféricas}
\addcontentsline{toc}{part}{Parte IV: Sistemas de Coordenadas Polares, Cilíndricas y Esféricas}

El estudio de la integración múltiple a menudo involucra regiones cuya descripción analítica en coordenadas cartesianas resulta excesivamente complicada o cuyas integrales asociadas son inmanejables. En esta sesión abordamos rigurosamente los sistemas de coordenadas polares (para el plano $\mathbb{R}^2$), y cilíndricas y esféricas (para el espacio $\mathbb{R}^3$). La elección del sistema de coordenadas adecuado permite transformar dominios geométricamente complejos, como discos, cilindros o esferas, en regiones más simples (idealmente rectangulares o paralelepípedos), facilitando sustancialmente el cálculo integral.

\section{Coordenadas Polares en $\mathbb{R}^2$}

\begin{definicionbox}{Sistema de Coordenadas Polares}
Dado un punto fijo $O$ llamado \textbf{polo} y un rayo inicial llamado \textbf{eje polar} en el plano, cada punto $P \in \mathbb{R}^2$ puede representarse mediante el par ordenado $(r, \theta)$, donde:
\begin{itemize}
    \item $r$ es la distancia dirigida desde el polo $O$ hasta el punto $P$.
    \item $\theta$ es el ángulo dirigido desde el eje polar hasta el segmento $OP$.
\end{itemize}
Convencionalmente, el polo se ubica en el origen $(0,0)$ del sistema cartesiano y el eje polar coincide con el semieje positivo de las abscisas. La relación entre las coordenadas cartesianas $(x,y)$ y las polares $(r,\theta)$ está dada por:
\[ x = r \cos \theta, \quad y = r \operatorname{sen} \theta \]
donde habitualmente se restringen los parámetros al dominio principal $r \geq 0$ y $0 \leq \theta < 2\pi$, o alternativamente $-\pi < \theta \leq \pi$, para garantizar la unicidad de representación, salvo en el polo, donde $r=0$ y $\theta$ es arbitrario.
Adicionalmente, se satisface que $x^2 + y^2 = r^2$ y, si $x \neq 0$, $\tan \theta = \frac{y}{x}$.
\end{definicionbox}

\begin{errorbox}{Ambigüedad en el ángulo}
Un error común al calcular $\theta$ mediante $\arctan(y/x)$ es no considerar el cuadrante en el que se ubica el punto $(x,y)$. La función arco tangente tiene imagen en $(-\pi/2, \pi/2)$, por lo que, si $x < 0$, es necesario sumar o restar $\pi$ al valor obtenido.
\end{errorbox}

\subsection{Esbozo de Regiones en Coordenadas Polares}

El análisis de curvas polares típicas es fundamental. Entre ellas destacan:
\begin{itemize}
    \item \textbf{Círculos centrados en el origen}: $r = a$ (constante).
    \item \textbf{Círculos desplazados}: $r = 2a \cos \theta$ (círculo de radio $|a|$ centrado en $(a,0)$) y $r = 2a \operatorname{sen} \theta$ (centro en $(0,a)$).
    \item \textbf{Cardioides}: $r = a(1 \pm \cos \theta)$ o $r = a(1 \pm \operatorname{sen} \theta)$, donde $a > 0$.
    \item \textbf{Rosas polares}: $r = a \cos(n\theta)$ o $r = a \operatorname{sen}(n\theta)$. Si $n$ es impar, tiene $n$ pétalos; si $n$ es par, tiene $2n$ pétalos.
    \item \textbf{Lemniscata de Bernoulli}: $r^2 = a^2 \cos(2\theta)$ o $r^2 = a^2 \operatorname{sen}(2\theta)$.
\end{itemize}

\begin{ejemplobox}{Intersección de Curvas Polares}
Halle y dibuje la región plana $\mathcal{D}$ interior a la curva $r = 3 \cos \theta$ y exterior a la curva $r = 1 + \cos \theta$.

\textbf{Resolución:}
La curva $r = 3 \cos \theta$ corresponde a un círculo de radio $3/2$ con centro en $(3/2, 0)$. La curva $r = 1 + \cos \theta$ es un cardioide simétrico respecto al eje polar.
Para hallar los puntos de intersección, igualamos ambas ecuaciones:
\[ 3 \cos \theta = 1 + \cos \theta \implies 2 \cos \theta = 1 \implies \cos \theta = \frac{1}{2} \]
En el intervalo principal $(-\pi, \pi]$, las soluciones son $\theta = -\frac{\pi}{3}$ y $\theta = \frac{\pi}{3}$.
La región se describe por:
\[ \mathcal{D} = \left\{ (r, \theta) : -\frac{\pi}{3} \leq \theta \leq \frac{\pi}{3}, \quad 1 + \cos \theta \leq r \leq 3 \cos \theta \right\} \]
\begin{center}
\begin{tikzpicture}[scale=1]
\draw[->] (-1.2,0) -- (4.2,0) node[right] {$x$};
\draw[->] (0,-2.5) -- (0,2.5) node[above] {$y$};
\draw[domain=-180:180, samples=150, smooth, variable=\t, blue, thick] plot ({3*cos(\t)*cos(\t)}, {3*cos(\t)*sin(\t)});
\draw[domain=-180:180, samples=150, smooth, variable=\t, red, thick] plot ({(1+cos(\t))*cos(\t)}, {(1+cos(\t))*sin(\t)});
\fill[blue, opacity=0.2] plot[domain=-60:60, samples=80, variable=\t] ({3*cos(\t)*cos(\t)}, {3*cos(\t)*sin(\t)}) -- plot[domain=60:-60, samples=80, variable=\t] ({(1+cos(\t))*cos(\t)}, {(1+cos(\t))*sin(\t)}) -- cycle;
\node at (2,0.5) {$\mathcal{D}$};
\end{tikzpicture}
\end{center}
\end{ejemplobox}

\section{Coordenadas Cilíndricas en $\mathbb{R}^3$}

\begin{definicionbox}{Sistema de Coordenadas Cilíndricas}
Las coordenadas cilíndricas $(r, \theta, z)$ extienden el sistema polar al espacio tridimensional mediante la adición de la coordenada cartesiana $z$. Las ecuaciones de transformación desde cartesianas a cilíndricas son:
\[ x = r \cos \theta, \quad y = r \operatorname{sen} \theta, \quad z = z \]
El dominio estándar para los parámetros es $r \geq 0$, $0 \leq \theta < 2\pi$ y $z \in \mathbb{R}$.
Además, $x^2 + y^2 = r^2$.
\end{definicionbox}

Las coordenadas cilíndricas resultan particularmente convenientes cuando el dominio de interés posee un eje de simetría (usualmente el eje $z$).

\textbf{Superficies notables en cilíndricas:}
\begin{itemize}
    \item \textbf{Cilindro circular recto} centrado en el eje $z$: $r = a$.
    \item \textbf{Planos que contienen al eje $z$}: $\theta = \alpha$ (un semiplano si $r \geq 0$).
    \item \textbf{Cono circular recto}: $z = c \sqrt{x^2+y^2} \implies z = cr$.
    \item \textbf{Paraboloide de revolución}: $z = c(x^2+y^2) \implies z = c r^2$.
    \item \textbf{Esfera}: $x^2+y^2+z^2 = a^2 \implies r^2 + z^2 = a^2$.
\end{itemize}

\begin{ejerciciobox}{Descripción de un Sólido}
Describa analíticamente en coordenadas cilíndricas el sólido $\mathcal{E}$ acotado inferiormente por el paraboloide $z = x^2 + y^2$ y superiormente por el plano $z = 4$.

\textbf{Resolución:}
\begin{enumerate}
    \item Sustituimos las ecuaciones a cilíndricas:
          El paraboloide $z = x^2 + y^2$ se transforma en $z = r^2$.
          El plano superior se mantiene como $z = 4$.
    \item Determinamos la intersección para proyectar sobre el plano $xy$:
          $r^2 = 4 \implies r = 2$ (dado que $r \geq 0$).
    \item La proyección $D$ en el plano polar es el disco $0 \leq r \leq 2$, con $\theta \in [0, 2\pi]$.
    \item Para un par $(r, \theta)$ fijo en $D$, la coordenada $z$ varía desde la superficie del paraboloide hasta el plano horizontal.
\end{enumerate}
Por lo tanto, la descripción del sólido es:
\[ \mathcal{E} = \{ (r, \theta, z) \in \mathbb{R}^3 : 0 \leq \theta \leq 2\pi, \; 0 \leq r \leq 2, \; r^2 \leq z \leq 4 \} \]
\end{ejerciciobox}

\section{Coordenadas Esféricas en $\mathbb{R}^3$}

\begin{definicionbox}{Sistema de Coordenadas Esféricas}
El sistema de coordenadas esféricas ubica un punto $P(x,y,z) \in \mathbb{R}^3$ utilizando la tupla $(\rho, \theta, \phi)$, donde:
\begin{itemize}
    \item $\rho$ es la distancia euclidiana del origen $O$ a $P$ ($\rho \geq 0$).
    \item $\theta$ es el ángulo azimutal, medido en el plano $xy$ desde el eje $x$ positivo hacia el eje $y$ positivo, como en polares ($0 \leq \theta < 2\pi$).
    \item $\phi$ es el ángulo de elevación o colatitud, medido desde el eje $z$ positivo hacia el vector $\vec{OP}$ ($0 \leq \phi \leq \pi$).
\end{itemize}
Las ecuaciones de transformación son:
\[ x = \rho \operatorname{sen} \phi \cos \theta \]
\[ y = \rho \operatorname{sen} \phi \operatorname{sen} \theta \]
\[ z = \rho \cos \phi \]
Se verifica inmediatamente que $\rho^2 = x^2 + y^2 + z^2$ y que $r = \rho \operatorname{sen} \phi$, vinculando las cilíndricas con las esféricas.
\end{definicionbox}

\begin{teoremabox}{Traducción de Desigualdades y Superficies}
En coordenadas esféricas, ciertas superficies clásicas adquieren expresiones sumamente simples:
\begin{enumerate}
    \item \textbf{Esfera centrada en el origen}: $\rho = a$ (con $a > 0$).
    \item \textbf{Semicono superior} de eje $z$: $\phi = c$ (con $0 < c < \pi/2$). Si $c \in (\pi/2, \pi)$, es un semicono inferior.
    \item \textbf{Semiplano} con bisagra en el eje $z$: $\theta = \alpha$.
    \item \textbf{Esfera desplazada tangente al origen}: Por ejemplo, una esfera de radio $a$ con centro en $(0,0,a)$ tiene ecuación cartesiana $x^2 + y^2 + (z-a)^2 = a^2$, que se expande a $x^2 + y^2 + z^2 = 2az$. Al sustituir esféricas obtenemos $\rho^2 = 2a (\rho \cos \phi)$, luego $\rho = 2a \cos \phi$, restringido a $0 \leq \phi \leq \pi/2$.
\end{enumerate}
\end{teoremabox}

\begin{ejemplobox}{Sólido en Forma de ``Cono de Helado''}
Considere el sólido $\mathcal{E}$ delimitado inferiormente por el cono $z = \sqrt{x^2+y^2}$ y superiormente por la esfera $x^2+y^2+z^2 = z$. Describa $\mathcal{E}$ en coordenadas esféricas.

\textbf{Resolución:}
\begin{enumerate}
    \item Transformamos la ecuación de la esfera: $\rho^2 = \rho \cos \phi \implies \rho = \cos \phi$.
    \item Transformamos la ecuación del cono: $z = \sqrt{x^2+y^2}$ implica, utilizando cilíndricas intermedias, $z = r$. En esféricas: $\rho \cos \phi = \rho \operatorname{sen} \phi \implies \tan \phi = 1 \implies \phi = \frac{\pi}{4}$.
    \item Intersección: La esfera y el cono se intersecan a lo largo de una circunferencia donde $\phi = \pi/4$.
    \item El sólido se encuentra "dentro" del semicono (es decir, el ángulo $\phi$ varía desde el eje $z$, $\phi=0$, hasta la pared del cono, $\phi=\pi/4$) y "debajo" del techo esférico (la distancia $\rho$ varía de $0$ a $\cos \phi$). El ángulo azimutal $\theta$ da una vuelta completa.
\end{enumerate}
En notación conjuntista:
\[ \mathcal{E} = \left\{ (\rho, \theta, \phi) : 0 \leq \theta \leq 2\pi, \; 0 \leq \phi \leq \frac{\pi}{4}, \; 0 \leq \rho \leq \cos \phi \right\} \]
\begin{center}
% Esbozo en 3D
\begin{tikzpicture}
\begin{axis}[
    view={35}{30},
    axis lines=center,
    xlabel=$x$, ylabel=$y$, zlabel=$z$,
    xmin=-0.6, xmax=0.6,
    ymin=-0.6, ymax=0.6,
    zmin=0, zmax=1.1,
    ticklabel style={font=\tiny},
    enlargelimits=true
]
% Cono
\addplot3[surf, domain=0:2*pi, y domain=0:0.5, colormap/cool, opacity=0.5, faceted color=black!60]
    ({y*cos(deg(x))}, {y*sin(deg(x))}, {y});
% Casquete esférico superior
\addplot3[surf, domain=0:2*pi, y domain=0:45, colormap/hot, opacity=0.5, faceted color=black!60]
    ({cos(y)*sin(y)*cos(deg(x))}, {cos(y)*sin(y)*sin(deg(x))}, {cos(y)*cos(y)});
\end{axis}
\end{tikzpicture}
\end{center}
\end{ejemplobox}

\begin{observacionbox}
Al cambiar de un sistema a otro en el cálculo de integrales múltiples, la transformación introduce un factor de escala en el elemento de volumen debido al Jacobiano de la transformación. Adelantamos que:
\begin{itemize}
    \item En cilíndricas: $dV = r \, dr \, d\theta \, dz$
    \item En esféricas: $dV = \rho^2 \operatorname{sen} \phi \, d\rho \, d\phi \, d\theta$
\end{itemize}
\end{observacionbox}
\part*{Parte V: Teorema General de Cambio de Variable}
\addcontentsline{toc}{part}{Parte V: Teorema General de Cambio de Variable}

\section*{Sesión 29: El Teorema de Cambio de Variable en Múltiples Dimensiones}

El cálculo de integrales múltiples sobre regiones complicadas en $\mathbb{R}^n$ suele requerir parametrizaciones complejas que dificultan la evaluación analítica. El teorema de cambio de variable proporciona una herramienta fundamental para transformar regiones de integración complejas en regiones más simples (como rectángulos o cajas), a costa de introducir un factor de distorsión geométrica conocido como el valor absoluto del determinante jacobiano. Este teorema generaliza la regla de sustitución para integrales de una variable.

\subsection*{El Determinante Jacobiano y su Interpretación Geométrica}

Antes de enunciar el teorema principal, es imprescindible definir formalmente la matriz jacobiana y su determinante, y entender su rol geométrico.

\begin{definicionbox}{Matriz Jacobiana y Determinante Jacobiano}
Sea $T: U \subset \mathbb{R}^n \to \mathbb{R}^n$ un campo vectorial de clase $C^1$ en un conjunto abierto $U$, dado por $T(u_1, u_2, \dots, u_n) = (x_1, x_2, \dots, x_n)$, donde $x_i = x_i(u_1, u_2, \dots, u_n)$ para $i=1, \dots, n$. La matriz jacobiana de $T$ en el punto $\mathbf{u} = (u_1, \dots, u_n) \in U$ se denota por $DT(\mathbf{u})$ o $J_T(\mathbf{u})$ y se define como la matriz de derivadas parciales:
\begin{equation*}
J_T(\mathbf{u}) = \begin{pmatrix}
\frac{\partial x_1}{\partial u_1} & \frac{\partial x_1}{\partial u_2} & \cdots & \frac{\partial x_1}{\partial u_n} \\
\frac{\partial x_2}{\partial u_1} & \frac{\partial x_2}{\partial u_2} & \cdots & \frac{\partial x_2}{\partial u_n} \\
\vdots & \vdots & \ddots & \vdots \\
\frac{\partial x_n}{\partial u_1} & \frac{\partial x_n}{\partial u_2} & \cdots & \frac{\partial x_n}{\partial u_n}
\end{pmatrix}
\end{equation*}
El determinante de esta matriz se conoce como el \textbf{Jacobiano} de $T$ y se denota comúnmente por:
\begin{equation*}
\det J_T(\mathbf{u}) = \frac{\partial(x_1, x_2, \dots, x_n)}{\partial(u_1, u_2, \dots, u_n)}
\end{equation*}
\end{definicionbox}

\begin{observacionbox}{Justificación Geométrica del Jacobiano}
El valor absoluto del determinante jacobiano, $|\det J_T(\mathbf{u})|$, mide la tasa local de expansión o contracción del volumen (o área en $\mathbb{R}^2$) bajo la transformación $T$ en las proximidades del punto $\mathbf{u}$.

Para entender esto, consideremos una transformación lineal $L: \mathbb{R}^2 \to \mathbb{R}^2$ dada por $L(u,v) = (au+bv, cu+dv)$. La matriz jacobiana de $L$ es constante:
\begin{equation*}
J_L = \begin{pmatrix} a & b \\ c & d \end{pmatrix}
\end{equation*}
Si tomamos un rectángulo fundamental en el plano $uv$ definido por los vectores $\mathbf{e}_1 = (1,0)$ y $\mathbf{e}_2 = (0,1)$, su área es 1. Bajo $L$, estos vectores se transforman en $\mathbf{v}_1 = L(\mathbf{e}_1) = (a,c)$ y $\mathbf{v}_2 = L(\mathbf{e}_2) = (b,d)$. El área del paralelogramo generado por $\mathbf{v}_1$ y $\mathbf{v}_2$ en el plano $xy$ es, por definición, el valor absoluto del determinante de la matriz formada por estos vectores:
\begin{equation*}
\text{Área} = \left| \det \begin{pmatrix} a & b \\ c & d \end{pmatrix} \right| = |ad - bc| = |\det J_L|
\end{equation*}
Dado que cualquier transformación $T$ de clase $C^1$ es localmente aproximable por su derivada (una transformación lineal), la imagen de un pequeño rectángulo $\Delta u \times \Delta v$ en el plano $uv$ se aproxima mediante un paralelogramo en el plano $xy$ con área aproximadamente igual a $|\det J_T(u,v)| \Delta u \Delta v$. Al integrar, este factor corrige la distorsión del área.
\end{observacionbox}

Para calcular transformaciones complejas, es frecuente que se conozca la transformación inversa de manera más natural. En estos casos, la siguiente propiedad es indispensable.

\begin{teoremabox}{Propiedad del Jacobiano Inverso}
Sea $T: U \subset \mathbb{R}^n \to V \subset \mathbb{R}^n$ un difeomorfismo de clase $C^1$ entre conjuntos abiertos, esto es, $T$ es biyectiva y tanto $T$ como $T^{-1}$ son continuamente diferenciables. Si denotamos las variables en $U$ por $\mathbf{u}$ y en $V$ por $\mathbf{x}$, entonces para todo $\mathbf{u} \in U$:
\begin{equation*}
\det J_{T^{-1}}(T(\mathbf{u})) = \frac{1}{\det J_T(\mathbf{u})}
\end{equation*}
Escrito en notación diferencial para $n=2$, si $(x,y) = T(u,v)$ y $(u,v) = T^{-1}(x,y)$, entonces:
\begin{equation*}
\frac{\partial(u,v)}{\partial(x,y)} = \frac{1}{\frac{\partial(x,y)}{\partial(u,v)}}
\end{equation*}
\end{teoremabox}

\subsection*{El Teorema General de Cambio de Variable}

\begin{teoremabox}{Teorema de Cambio de Variable para Integrales Múltiples}
Sean $U$ y $V$ subconjuntos abiertos de $\mathbb{R}^n$. Sea $T: U \to V$ una transformación de clase $C^1$. Supóngase que:
\begin{enumerate}
    \item $T$ es inyectiva en $U \setminus N$, donde $N$ es un conjunto de contenido $n$-dimensional nulo (por ejemplo, la frontera de una región).
    \item $T(U) = V$, es decir, $T$ es sobreyectiva sobre $V$.
    \item $\det J_T(\mathbf{u}) \neq 0$ para todo $\mathbf{u} \in U \setminus N$.
\end{enumerate}
Sea $f: V \to \mathbb{R}$ una función integrable. Entonces $f \circ T$ es integrable en $U$, y se cumple que:
\begin{equation*}
\int \dots \int_V f(x_1, \dots, x_n) \, dx_1 \dots dx_n = \int \dots \int_U f(T(u_1, \dots, u_n)) \left| \frac{\partial(x_1, \dots, x_n)}{\partial(u_1, \dots, u_n)} \right| \, du_1 \dots du_n
\end{equation*}
\end{teoremabox}

\begin{errorbox}{Errores Comunes en el Cambio de Variable}
1. \textbf{Olvido del valor absoluto:} El factor de corrección es el valor \textit{absoluto} del determinante jacobiano. Ignorar el signo puede llevar a áreas o volúmenes negativos.
2. \textbf{Falta de inyectividad:} Aplicar el teorema en regiones donde $T$ no es inyectiva (salvo en conjuntos de medida nula) duplica áreas. Por ejemplo, en polares, integrar $\theta$ de $0$ a $4\pi$ sobrecontará el volumen por un factor de 2.
3. \textbf{Confusión entre directa e inversa:} A veces se calcula $\frac{\partial(u,v)}{\partial(x,y)}$ y se multiplica directamente, en lugar de dividir o usar la transformación inversa.
\end{errorbox}

\subsection*{Aplicaciones en Integrales Dobles}

Una de las aplicaciones más elegantes del cambio de variable es la resolución de la integral gaussiana.

\begin{ejerciciobox}{Coordenadas Polares y la Integral de Gauss}
Demuestre de manera rigurosa que:
\begin{equation*}
I = \int_{-\infty}^{\infty} e^{-x^2} \, dx = \sqrt{\pi}
\end{equation*}
\textbf{Solución:}
La integral definida $I$ es impropia y la función $e^{-x^2}$ no tiene una antiderivada elemental. El truco clásico es elevar $I$ al cuadrado. Dado que $x$ es una variable muda, podemos escribir:
\begin{equation*}
I^2 = \left( \int_{-\infty}^{\infty} e^{-x^2} \, dx \right) \left( \int_{-\infty}^{\infty} e^{-y^2} \, dy \right)
\end{equation*}
Utilizando el Teorema de Fubini para integrales impropias (justificado por ser la función integranda positiva, aplicando el Teorema de Tonelli), podemos combinar las integrales:
\begin{equation*}
I^2 = \int_{-\infty}^{\infty} \int_{-\infty}^{\infty} e^{-x^2} e^{-y^2} \, dx \, dy = \iint_{\mathbb{R}^2} e^{-(x^2+y^2)} \, dx \, dy
\end{equation*}
Ahora, aplicaremos un cambio de variable a coordenadas polares. Sea la transformación $T: (0, \infty) \times [0, 2\pi) \to \mathbb{R}^2 \setminus \{(0,0)\}$ dada por:
\begin{equation*}
x = r \cos \theta, \quad y = r \sin \theta
\end{equation*}
Calculamos el jacobiano de esta transformación:
\begin{equation*}
J_T(r, \theta) = \begin{pmatrix} \frac{\partial x}{\partial r} & \frac{\partial x}{\partial \theta} \\ \frac{\partial y}{\partial r} & \frac{\partial y}{\partial \theta} \end{pmatrix} = \begin{pmatrix} \cos \theta & -r \sin \theta \\ \sin \theta & r \cos \theta \end{pmatrix}
\end{equation*}
\begin{equation*}
\det J_T(r, \theta) = (\cos \theta)(r \cos \theta) - (-r \sin \theta)(\sin \theta) = r \cos^2 \theta + r \sin^2 \theta = r
\end{equation*}
Como $r > 0$, el valor absoluto del jacobiano es $|r| = r$. La región $\mathbb{R}^2$ se describe en polares por el límite cuando $R \to \infty$ del disco $D_R$ de radio $R$. El conjunto donde $T$ falla ser inyectiva o sobreyectiva es el origen, de medida nula.
Sustituyendo en la integral:
\begin{equation*}
I^2 = \lim_{R \to \infty} \iint_{D_R} e^{-(x^2+y^2)} \, dx \, dy = \lim_{R \to \infty} \int_{0}^{2\pi} \int_{0}^{R} e^{-r^2} r \, dr \, d\theta
\end{equation*}
Resolvemos primero la integral interior respecto a $r$. Realizamos la sustitución $u = r^2$, lo que implica $du = 2r \, dr \implies \frac{du}{2} = r \, dr$. Los límites cambian: si $r=0 \implies u=0$; si $r=R \implies u=R^2$.
\begin{equation*}
\int_{0}^{R} e^{-r^2} r \, dr = \int_{0}^{R^2} e^{-u} \frac{du}{2} = \frac{1}{2} \left[ -e^{-u} \right]_{0}^{R^2} = \frac{1}{2} ( -e^{-R^2} - (-1) ) = \frac{1}{2} (1 - e^{-R^2})
\end{equation*}
Sustituyendo este resultado en la integral exterior:
\begin{equation*}
I^2 = \lim_{R \to \infty} \int_{0}^{2\pi} \frac{1}{2} (1 - e^{-R^2}) \, d\theta = \lim_{R \to \infty} \left[ \frac{1}{2} (1 - e^{-R^2}) \theta \right]_{0}^{2\pi} = \lim_{R \to \infty} \pi (1 - e^{-R^2})
\end{equation*}
Tomando el límite cuando $R \to \infty$, notamos que $e^{-R^2} \to 0$:
\begin{equation*}
I^2 = \pi(1 - 0) = \pi
\end{equation*}
Dado que la función $e^{-x^2}$ es estrictamente positiva para todo $x$, su integral $I$ debe ser positiva. Por lo tanto, extraemos la raíz cuadrada positiva:
\begin{equation*}
I = \sqrt{\pi}
\end{equation*}
Esto concluye la demostración rigurosa y sin saltos de la integral de Gauss.
\end{ejerciciobox}

En ocasiones, las fronteras de la región de integración sugieren directamente un cambio de variable no lineal.

\begin{ejerciciobox}{Región acotada por hipérbolas y rectas}
Calcule el área de la región $D$ en el primer cuadrante acotada por las hipérbolas $xy = 1$, $xy = 4$ y las rectas $y = x$, $y = 3x$.

\textbf{Solución:}
El área de la región $D$ está dada por $A = \iint_D 1 \, dx \, dy$. Las curvas que limitan $D$ sugieren el siguiente cambio de variables. Definimos:
\begin{equation*}
u = xy \quad \text{y} \quad v = \frac{y}{x}
\end{equation*}
Con esta elección, las fronteras de $D$ en el plano $xy$ se transforman en fronteras rectangulares en el plano $uv$.
Específicamente, $xy=1 \implies u=1$; $xy=4 \implies u=4$; $y=x \implies y/x=1 \implies v=1$; $y=3x \implies y/x=3 \implies v=3$.
Así, la nueva región de integración $D^*$ es el rectángulo $D^* = \{(u,v) \mid 1 \le u \le 4, \, 1 \le v \le 3\}$.

En lugar de despejar $x$ e $y$ en términos de $u$ y $v$ (lo cual es posible pero tedioso), calcularemos el jacobiano inverso $\frac{\partial(u,v)}{\partial(x,y)}$ y usaremos la propiedad del jacobiano recíproco.
Calculamos las derivadas parciales:
\begin{align*}
\frac{\partial u}{\partial x} &= y, & \frac{\partial u}{\partial y} &= x \\
\frac{\partial v}{\partial x} &= -\frac{y}{x^2}, & \frac{\partial v}{\partial y} &= \frac{1}{x}
\end{align*}
El jacobiano inverso es:
\begin{equation*}
\frac{\partial(u,v)}{\partial(x,y)} = \det \begin{pmatrix} y & x \\ -\frac{y}{x^2} & \frac{1}{x} \end{pmatrix} = (y)\left(\frac{1}{x}\right) - (x)\left(-\frac{y}{x^2}\right) = \frac{y}{x} + \frac{y}{x} = 2\frac{y}{x} = 2v
\end{equation*}
Por la propiedad del jacobiano inverso:
\begin{equation*}
\left| \frac{\partial(x,y)}{\partial(u,v)} \right| = \frac{1}{\left| \frac{\partial(u,v)}{\partial(x,y)} \right|} = \frac{1}{|2v|}
\end{equation*}
Dado que en la región $D^*$, $1 \le v \le 3$, sabemos que $v > 0$, por lo que $|2v| = 2v$. El determinante jacobiano es $\frac{1}{2v}$.
Planteamos la integral del área:
\begin{equation*}
A = \iint_D 1 \, dx \, dy = \iint_{D^*} 1 \left| \frac{\partial(x,y)}{\partial(u,v)} \right| \, du \, dv = \int_{1}^{3} \int_{1}^{4} \frac{1}{2v} \, du \, dv
\end{equation*}
Integramos primero con respecto a $u$:
\begin{equation*}
\int_{1}^{4} \frac{1}{2v} \, du = \frac{1}{2v} \Big[ u \Big]_{1}^{4} = \frac{1}{2v} (4 - 1) = \frac{3}{2v}
\end{equation*}
Ahora integramos con respecto a $v$:
\begin{equation*}
A = \int_{1}^{3} \frac{3}{2v} \, dv = \frac{3}{2} \Big[ \ln|v| \Big]_{1}^{3} = \frac{3}{2} (\ln 3 - \ln 1) = \frac{3}{2} \ln 3
\end{equation*}
El área de la región es $\frac{3}{2} \ln 3$. Se aprecia cómo el cambio de variable simplifica enormemente la integral.
\end{ejerciciobox}


\begin{ejemplobox}{Transformaciones Lineales: Un paralelogramo}
Para evaluar la integral $\iint_R (x+y)^2 e^{x-y} \, dx \, dy$ donde $R$ es el paralelogramo con vértices en $(0,0)$, $(1,1)$, $(2,0)$ y $(1,-1)$.
Un bosquejo rápido de la región muestra que está acotada por las rectas $x-y=0$, $x-y=2$, $x+y=0$, $x+y=2$.
Definimos $u = x+y$ y $v = x-y$. La región en el plano $uv$ es el rectángulo $0 \le u \le 2$ y $0 \le v \le 2$.
El jacobiano inverso es:
\begin{equation*}
\frac{\partial(u,v)}{\partial(x,y)} = \det \begin{pmatrix} 1 & 1 \\ 1 & -1 \end{pmatrix} = -1 - 1 = -2
\end{equation*}
Entonces, el jacobiano directo es $\left| \frac{\partial(x,y)}{\partial(u,v)} \right| = \frac{1}{|-2|} = \frac{1}{2}$.
La integral se convierte en:
\begin{equation*}
\iint_{R^*} u^2 e^v \frac{1}{2} \, du \, dv = \frac{1}{2} \left( \int_{0}^{2} u^2 \, du \right) \left( \int_{0}^{2} e^v \, dv \right) = \frac{1}{2} \left[ \frac{u^3}{3} \right]_0^2 \left[ e^v \right]_0^2 = \frac{1}{2} \left(\frac{8}{3}\right) (e^2 - 1) = \frac{4}{3}(e^2 - 1)
\end{equation*}
\end{ejemplobox}


\subsection*{Cambio de Variable en Integrales Triples}

El principio se extiende a cualquier dimensión. En $\mathbb{R}^3$, el factor de corrección es el determinante de una matriz $3 \times 3$. Consideremos un caso fundamental en la geometría espacial.

\begin{ejerciciobox}{Volumen del Elipsoide Sólido}
Demuestre que el volumen del sólido $E$ encerrado por el elipsoide
\begin{equation*}
\frac{x^2}{a^2} + \frac{y^2}{b^2} + \frac{z^2}{c^2} \le 1
\end{equation*}
donde $a, b, c > 0$, es $V = \frac{4}{3}\pi abc$.

\textbf{Solución:}
El volumen de la región tridimensional $E$ viene dado por la integral triple:
\begin{equation*}
V = \iiint_E 1 \, dx \, dy \, dz
\end{equation*}
La región es compleja de integrar directamente en coordenadas cartesianas, aunque posible, produce integrales muy extensas. Proponemos un cambio de variables que estira (o comprime) los ejes para transformar el elipsoide en una esfera de radio 1. Definimos las nuevas variables $u, v, w$:
\begin{equation*}
x = au, \quad y = bv, \quad z = cw
\end{equation*}
Bajo esta transformación $T(u,v,w) = (au, bv, cw)$, la inecuación que define $E$ se transforma en:
\begin{equation*}
\frac{(au)^2}{a^2} + \frac{(bv)^2}{b^2} + \frac{(cw)^2}{c^2} \le 1 \implies u^2 + v^2 + w^2 \le 1
\end{equation*}
Esta última inecuación define la bola unitaria cerrada $B$ en $\mathbb{R}^3$ con respecto a las variables $u,v,w$.
Calculamos el determinante jacobiano de la transformación:
\begin{equation*}
J_T(u,v,w) = \begin{pmatrix}
\frac{\partial x}{\partial u} & \frac{\partial x}{\partial v} & \frac{\partial x}{\partial w} \\
\frac{\partial y}{\partial u} & \frac{\partial y}{\partial v} & \frac{\partial y}{\partial w} \\
\frac{\partial z}{\partial u} & \frac{\partial z}{\partial v} & \frac{\partial z}{\partial w}
\end{pmatrix} = \begin{pmatrix}
a & 0 & 0 \\
0 & b & 0 \\
0 & 0 & c
\end{pmatrix}
\end{equation*}
El determinante es el producto de la diagonal de esta matriz diagonal:
\begin{equation*}
\det J_T(u,v,w) = a \cdot b \cdot c = abc
\end{equation*}
Dado que $a, b, c > 0$, el valor absoluto es $|abc| = abc$. Aplicando el teorema de cambio de variable:
\begin{equation*}
V = \iiint_E 1 \, dx \, dy \, dz = \iiint_B 1 \cdot \left| \frac{\partial(x,y,z)}{\partial(u,v,w)} \right| \, du \, dv \, dw = \iiint_B abc \, du \, dv \, dw
\end{equation*}
Dado que $abc$ es una constante, sale de la integral:
\begin{equation*}
V = abc \iiint_B 1 \, du \, dv \, dw
\end{equation*}
La integral triple $\iiint_B 1 \, du \, dv \, dw$ representa geométricamente el volumen de la bola unitaria $B$ de radio 1 en $\mathbb{R}^3$. Sabemos por geometría básica (o por coordenadas esféricas) que el volumen de una esfera de radio $R$ es $\frac{4}{3}\pi R^3$. Por lo tanto, para $R=1$:
\begin{equation*}
\text{Volumen}(B) = \frac{4}{3}\pi (1)^3 = \frac{4}{3}\pi
\end{equation*}
Si deseamos demostrar esto también rigurosamente usando esféricas ($u = \rho \sin\phi \cos\theta, v = \rho \sin\phi \sin\theta, w = \rho \cos\phi$ con jacobiano $\rho^2 \sin\phi$):
\begin{equation*}
\iiint_B 1 \, du \, dv \, dw = \int_0^{2\pi} \int_0^\pi \int_0^1 \rho^2 \sin\phi \, d\rho \, d\phi \, d\theta
\end{equation*}
\begin{equation*}
= \left( \int_0^{2\pi} d\theta \right) \left( \int_0^\pi \sin\phi \, d\phi \right) \left( \int_0^1 \rho^2 \, d\rho \right)
\end{equation*}
\begin{equation*}
= \Big[ \theta \Big]_0^{2\pi} \cdot \Big[ -\cos\phi \Big]_0^\pi \cdot \left[ \frac{\rho^3}{3} \right]_0^1 = (2\pi) \cdot (-\cos\pi + \cos 0) \cdot \left(\frac{1}{3}\right) = 2\pi \cdot (1+1) \cdot \frac{1}{3} = \frac{4}{3}\pi
\end{equation*}
Reemplazando en la expresión del volumen:
\begin{equation*}
V = abc \left( \frac{4}{3}\pi \right) = \frac{4}{3}\pi abc
\end{equation*}
Quedando demostrado rigurosamente y paso a paso que el volumen de un elipsoide con semiejes $a, b$ y $c$ es $\frac{4}{3}\pi abc$.
\end{ejerciciobox}

\begin{observacionbox}{Composición de Cambios de Variable}
El ejercicio anterior ilustra indirectamente una propiedad fundamental del cambio de variables: se pueden componer transformaciones. Si tenemos una región muy complicada, podríamos aplicar un cambio $T_1$ a unas variables intermedias, y luego un cambio $T_2$ a unas variables finales. El jacobiano total será el producto de los jacobianos, acorde a la regla de la cadena para la matriz jacobiana $J_{T_2 \circ T_1} = J_{T_2} J_{T_1}$.
En el caso del elipsoide, implícitamente hicimos la transformación lineal $x=au, y=bv, z=cw$ seguida de una transformación a coordenadas esféricas tridimensionales.
\end{observacionbox}

La comprensión e internalización del Teorema de Cambio de Variable resulta capital no solo para el análisis matemático, sino que subyace a toda la geometría diferencial moderna a través del concepto de integración sobre variedades orientables mediante cartas y particiones de la unidad.
\part*{Parte VI: Integrales en cilíndricas, esféricas y en $\mathbb{R}^n$}

El estudio de las integrales triples se enriquece profundamente al adoptar sistemas de coordenadas curvilíneos que se adapten a las simetrías intrínsecas de la región de integración. En esta sección abordaremos las coordenadas cilíndricas y esféricas, así como la generalización del concepto de integral a espacios euclidianos de dimensión arbitraria, $\mathbb{R}^n$.

\section{Coordenadas Cilíndricas y Esféricas: Fundamentos Teóricos}

\subsection{Coordenadas Cilíndricas}

Las coordenadas cilíndricas $(r, \theta, z)$ constituyen una extensión natural de las coordenadas polares en $\mathbb{R}^2$ incorporando la altura $z$ como tercera dimensión.

\begin{definicionbox}{Transformación a Coordenadas Cilíndricas}{}
    Se define la transformación $T: D \subset \mathbb{R}^3_{r,\theta,z} \to \mathbb{R}^3_{x,y,z}$ dada por:
    \[
    T(r, \theta, z) = (r \cos\theta, r \sin\theta, z)
    \]
    donde usualmente $r \geq 0$, $0 \leq \theta < 2\pi$ y $z \in \mathbb{R}$.
\end{definicionbox}

Para realizar el cambio de variable en integrales triples, es imperativo calcular el determinante de la matriz Jacobiana de $T$.

\begin{teoremabox}{Jacobiano de las Coordenadas Cilíndricas}{}
    Sea $T(r, \theta, z)$ la transformación a coordenadas cilíndricas. El Jacobiano de la transformación es $J_T = r$.
\end{teoremabox}

\textit{Demostración.} Calculamos la matriz Jacobiana:
\[
D T(r,\theta,z) = \begin{pmatrix}
\frac{\partial x}{\partial r} & \frac{\partial x}{\partial \theta} & \frac{\partial x}{\partial z} \\
\frac{\partial y}{\partial r} & \frac{\partial y}{\partial \theta} & \frac{\partial y}{\partial z} \\
\frac{\partial z}{\partial r} & \frac{\partial z}{\partial \theta} & \frac{\partial z}{\partial z}
\end{pmatrix} = \begin{pmatrix}
\cos\theta & -r \sin\theta & 0 \\
\sin\theta & r \cos\theta & 0 \\
0 & 0 & 1
\end{pmatrix}
\]
El determinante se evalúa expandiendo por la tercera fila:
\[
\det(D T) = 1 \cdot \left[ (\cos\theta)(r \cos\theta) - (-r \sin\theta)(\sin\theta) \right] = r \cos^2\theta + r \sin^2\theta = r.
\]
Puesto que $r \geq 0$, el valor absoluto del Jacobiano es $|J_T| = r$. Q.E.D.

\subsection{Coordenadas Esféricas}

Las coordenadas esféricas $(\rho, \phi, \theta)$ resultan óptimas para regiones con simetría puntual, donde $\rho$ es la distancia al origen, $\phi$ es el ángulo polar (respecto al eje $z$ positivo) y $\theta$ es el ángulo acimutal.

\begin{definicionbox}{Transformación a Coordenadas Esféricas}{}
    Se define la transformación $T: D \subset \mathbb{R}^3_{\rho,\phi,\theta} \to \mathbb{R}^3_{x,y,z}$ mediante:
    \begin{align*}
        x &= \rho \sin\phi \cos\theta \\
        y &= \rho \sin\phi \sin\theta \\
        z &= \rho \cos\phi
    \end{align*}
    donde $\rho \geq 0$, $0 \leq \phi \leq \pi$ y $0 \leq \theta < 2\pi$.
\end{definicionbox}

\begin{teoremabox}{Jacobiano de las Coordenadas Esféricas}{}
    Sea $T(\rho, \phi, \theta)$ la transformación a esféricas. El Jacobiano de la transformación es $J_T = \rho^2 \sin\phi$.
\end{teoremabox}

\textit{Demostración.} La matriz Jacobiana es:
\[
D T = \begin{pmatrix}
\sin\phi \cos\theta & \rho \cos\phi \cos\theta & -\rho \sin\phi \sin\theta \\
\sin\phi \sin\theta & \rho \cos\phi \sin\theta & \rho \sin\phi \cos\theta \\
\cos\phi & -\rho \sin\phi & 0
\end{pmatrix}
\]
Calculamos el determinante expandiendo por la tercera fila:
\begin{align*}
\det(D T) &= \cos\phi \begin{vmatrix} \rho \cos\phi \cos\theta & -\rho \sin\phi \sin\theta \\ \rho \cos\phi \sin\theta & \rho \sin\phi \cos\theta \end{vmatrix} - (-\rho \sin\phi) \begin{vmatrix} \sin\phi \cos\theta & -\rho \sin\phi \sin\theta \\ \sin\phi \sin\theta & \rho \sin\phi \cos\theta \end{vmatrix} \\
&= \cos\phi \left[ \rho^2 \sin\phi \cos\phi \cos^2\theta + \rho^2 \sin\phi \cos\phi \sin^2\theta \right] + \rho \sin\phi \left[ \rho \sin^2\phi \cos^2\theta + \rho \sin^2\phi \sin^2\theta \right] \\
&= \cos\phi \left[ \rho^2 \sin\phi \cos\phi \right] + \rho \sin\phi \left[ \rho \sin^2\phi \right] \\
&= \rho^2 \sin\phi \cos^2\phi + \rho^2 \sin^3\phi = \rho^2 \sin\phi (\cos^2\phi + \sin^2\phi) = \rho^2 \sin\phi.
\end{align*}
Como $\phi \in [0, \pi]$, $\sin\phi \geq 0$, por lo que el valor absoluto del Jacobiano coincide con $\rho^2 \sin\phi$. Q.E.D.

\begin{errorbox}{Precaución con los ángulos en esféricas}{}
    Una fuente frecuente de error en exámenes es confundir la convención física y la convención matemática. En este curso (convención matemática estándar), $\phi$ es el ángulo que forma el radio vector con el eje $z$ positivo. ¡Verifica siempre la posición de las variables antes de integrar! Además, ¡nunca olvides multiplicar por el Jacobiano $\rho^2 \sin\phi$!
\end{errorbox}


\section{Problemas de Integración en $\mathbb{R}^3$}

A continuación, analizaremos diversos problemas tipo PC y examen.

\begin{ejemplobox}{Sólido entre paraboloide y cono}{}
    Halle el volumen de la región $E$ limitada inferiormente por el paraboloide $z = x^2 + y^2$ y superiormente por el cono $z = 2 - \sqrt{x^2+y^2}$.
\end{ejemplobox}

\textit{Solución.} Observemos la simetría de la región alrededor del eje $z$, la cual sugiere que el uso de coordenadas cilíndricas es apropiado.
En coordenadas cilíndricas, las superficies son:
1. Paraboloide (límite inferior): $z = r^2$.
2. Cono (límite superior): $z = 2 - r$.

Intersección de ambas superficies para hallar la proyección sobre el plano $xy$:
\[
r^2 = 2 - r \implies r^2 + r - 2 = 0 \implies (r+2)(r-1) = 0.
\]
Dado que $r \geq 0$, la única solución válida es $r=1$. Por lo tanto, la proyección de la región en el plano $xy$ es el disco $D = \{(r,\theta): 0 \leq r \leq 1, \, 0 \leq \theta \leq 2\pi\}$.
Los límites para $z$ en función de $r$ serán desde el paraboloide hasta el cono: $r^2 \leq z \leq 2-r$.
El volumen se calcula como:
\begin{align*}
V(E) &= \iiint_E 1 \, dV = \int_0^{2\pi} \int_0^1 \int_{r^2}^{2-r} 1 \, (r \, dz \, dr \, d\theta) \\
&= \int_0^{2\pi} d\theta \int_0^1 r [z]_{r^2}^{2-r} \, dr \\
&= 2\pi \int_0^1 r (2 - r - r^2) \, dr \\
&= 2\pi \int_0^1 (2r - r^2 - r^3) \, dr \\
&= 2\pi \left[ r^2 - \frac{r^3}{3} - \frac{r^4}{4} \right]_0^1 \\
&= 2\pi \left( 1 - \frac{1}{3} - \frac{1}{4} \right) = 2\pi \left( \frac{12 - 4 - 3}{12} \right) = 2\pi \left( \frac{5}{12} \right) = \frac{5\pi}{6}.
\end{align*}


\begin{ejerciciobox}{Integral sobre una esfera desplazada (Ref: PC5 2024-II)}{}
    Calcule $\iiint_{E} z \, dV$, donde $E$ es la región interior a la esfera $x^2 + y^2 + z^2 \leq 2z$.
\end{ejerciciobox}

\textit{Solución.} Completamos cuadrados para identificar la región:
\[
x^2 + y^2 + (z-1)^2 \leq 1
\]
Se trata de una esfera de radio $1$ centrada en $(0,0,1)$. Como la región es esférica y pasa por el origen, es útil emplear coordenadas esféricas.
La ecuación de la frontera en esféricas se obtiene sustituyendo:
\[
\rho^2 = 2(\rho \cos\phi) \implies \rho = 2\cos\phi.
\]
Para que $\rho \geq 0$, es necesario que $\cos\phi \geq 0$, lo que implica $0 \leq \phi \leq \pi/2$. El ángulo $\theta$ recorre todo su dominio: $0 \leq \theta \leq 2\pi$. Por consiguiente, la región $E$ se describe por:
\[
E = \{(\rho, \phi, \theta) \mid 0 \leq \theta \leq 2\pi, \, 0 \leq \phi \leq \pi/2, \, 0 \leq \rho \leq 2\cos\phi\}.
\]
La función integrando es $z = \rho \cos\phi$. Incorporando el Jacobiano de la transformación ($\rho^2 \sin\phi$), la integral adopta la forma:
\begin{align*}
\iiint_{E} z \, dV &= \int_0^{2\pi} \int_0^{\pi/2} \int_0^{2\cos\phi} (\rho \cos\phi) (\rho^2 \sin\phi) \, d\rho \, d\phi \, d\theta \\
&= \int_0^{2\pi} d\theta \int_0^{\pi/2} \cos\phi \sin\phi \left[ \int_0^{2\cos\phi} \rho^3 \, d\rho \right] d\phi \\
&= 2\pi \int_0^{\pi/2} \cos\phi \sin\phi \left[ \frac{\rho^4}{4} \right]_0^{2\cos\phi} \, d\phi \\
&= 2\pi \int_0^{\pi/2} \cos\phi \sin\phi \left( 4\cos^4\phi \right) \, d\phi \\
&= 8\pi \int_0^{\pi/2} \cos^5\phi \sin\phi \, d\phi.
\end{align*}
Para evaluar la última integral, aplicamos la sustitución $u = \cos\phi \implies du = -\sin\phi \, d\phi$. Los límites de integración cambian de $\phi=0 \to u=1$ y $\phi=\pi/2 \to u=0$:
\begin{align*}
\iiint_{E} z \, dV &= 8\pi \int_1^0 u^5 (-du) = 8\pi \int_0^1 u^5 \, du \\
&= 8\pi \left[ \frac{u^6}{6} \right]_0^1 = 8\pi \left( \frac{1}{6} \right) = \frac{4\pi}{3}.
\end{align*}

\begin{observacionbox}{Coordenadas Esféricas Generalizadas}{}
    Para integrar sobre regiones elipsoidales de la forma $\frac{x^2}{a^2} + \frac{y^2}{b^2} + \frac{z^2}{c^2} \leq 1$, se introducen las \emph{coordenadas esféricas generalizadas}:
    $x = a \rho \sin\phi \cos\theta, \quad y = b \rho \sin\phi \sin\theta, \quad z = c \rho \cos\phi$.
    El Jacobiano en este caso es el producto del Jacobiano de las esféricas usuales por la dilatación asimétrica: $|J| = a b c \, \rho^2 \sin\phi$.
\end{observacionbox}

\section{Integrales Múltiples en $\mathbb{R}^n$}

La teoría desarrollada para dimensiones bajas se puede extender de forma natural a funciones $f: \Omega \subset \mathbb{R}^n \to \mathbb{R}$, en donde la integral $n$-múltiple mide el "hipervolumen" u otras magnitudes ponderadas en $\mathbb{R}^n$. Por el Teorema de Fubini generalizado, si la región es elemental y la función es integrable, la integral múltiple se evalúa como $n$ integrales iteradas simples.

\subsection{Volumen del $n$-simplex estándar}

El $n$-simplex estándar se define como:
\[
\Delta^n = \{ (x_1, \ldots, x_n) \in \mathbb{R}^n \mid x_i \geq 0, \, x_1 + \dots + x_n \leq 1 \}
\]

\begin{ejemplobox}{Volumen de $\Delta^n$}{}
    Demuestre de manera rigurosa que el $n$-volumen de $\Delta^n$ es $V_n = \frac{1}{n!}$.
\end{ejemplobox}

\textit{Demostración.} Usaremos inducción matemática sobre la dimensión $n$.
Caso base: Para $n=1$, $\Delta^1 = \{ x_1 \mid 0 \leq x_1 \leq 1 \}$. Claramente $V_1 = \int_0^1 dx_1 = 1 = \frac{1}{1!}$. Se cumple.
Supongamos que para $n-1$, el volumen es $V_{n-1}(t) = \frac{t^{n-1}}{(n-1)!}$ para el simplex modificado $\Delta^{n-1}(t) = \{ x_i \geq 0, \sum_{i=1}^{n-1} x_i \leq t \}$.
Para el paso inductivo, planteamos el volumen del simplex de dimensión $n$:
\[
V_n = \int_{\Delta^n} dV = \int_0^1 \left( \int_{\Delta^{n-1}(1-x_n)} d x_1 \cdots d x_{n-1} \right) dx_n
\]
Donde la integral interna corresponde exactamente al volumen de un $(n-1)$-simplex a escala, cuyo límite máximo es $1-x_n$. Por la hipótesis inductiva, esta integral vale $\frac{(1-x_n)^{n-1}}{(n-1)!}$. Sustituyendo:
\[
V_n = \int_0^1 \frac{(1-x_n)^{n-1}}{(n-1)!} dx_n
\]
Haciendo $u = 1-x_n \implies du = -dx_n$, límites $1 \to 0$:
\[
V_n = \int_1^0 \frac{u^{n-1}}{(n-1)!} (-du) = \int_0^1 \frac{u^{n-1}}{(n-1)!} du = \left[ \frac{u^n}{n(n-1)!} \right]_0^1 = \frac{1}{n!}.
\]
Esto demuestra rigurosamente que $V(\Delta^n) = \frac{1}{n!}$. Q.E.D.

\subsection{Volumen de la bola $n$-dimensional: El caso de $\mathbb{R}^4$}

El volumen de la bola de radio $R$ en $\mathbb{R}^n$, denotada por $B_n(R) = \{ \mathbf{x} \in \mathbb{R}^n \mid \|\mathbf{x}\|^2 \leq R^2 \}$, obedece a fórmulas muy interesantes dependiendo de la dimensión. Para dimensiones altas, se utilizan funciones Especiales ($\Gamma(x)$). 

\begin{ejerciciobox}{Volumen de la Hiperbola en $\mathbb{R}^4$}{}
    Calcule explícitamente el hipervolumen de la bola tetradimensional de radio $R$, $B_4(R) = \{ (x,y,z,w) \mid x^2+y^2+z^2+w^2 \leq R^2 \}$.
\end{ejerciciobox}

\textit{Solución.} El volumen $V_4(R)$ se define como la integral sobre la región tetradimensional:
\[
V_4(R) = \int\limits_{B_4(R)} dx \, dy \, dz \, dw
\]
Podemos seccionar la bola en $\mathbb{R}^4$ fijando dos coordenadas, digamos $(z,w)$. Sea $D_R$ el disco bi-dimensional $z^2+w^2 \leq R^2$. Para cada punto fijo $(z,w) \in D_R$, las variables $x, y$ deben satisfacer:
\[
x^2+y^2 \leq R^2 - z^2 - w^2
\]
Esto define un círculo en el plano $xy$ con radio al cuadrado $r^2 = R^2 - (z^2+w^2)$. Por lo tanto, el área de esta sección en el plano $xy$ (fijados $z, w$) es simplemente el área del disco correspondiente:
\[
A_{xy}(z,w) = \pi (R^2 - z^2 - w^2).
\]
Luego, integramos esta área $A_{xy}$ sobre todos los valores posibles de $(z,w)$, es decir, sobre $D_R$:
\[
V_4(R) = \iint_{D_R} \pi (R^2 - z^2 - w^2) \, dz \, dw.
\]
Para calcular esta integral doble, pasamos a coordenadas polares en el plano $zw$, digamos $z = \rho \cos\alpha$ y $w = \rho \sin\alpha$. El disco $D_R$ se mapea a $0 \leq \rho \leq R$ y $0 \leq \alpha \leq 2\pi$. El Jacobiano es $\rho$. Sustituyendo:
\begin{align*}
V_4(R) &= \int_0^{2\pi} \int_0^R \pi (R^2 - \rho^2) \rho \, d\rho \, d\alpha \\
&= \pi \left( \int_0^{2\pi} d\alpha \right) \left( \int_0^R (R^2 \rho - \rho^3) \, d\rho \right) \\
&= \pi (2\pi) \left[ R^2 \frac{\rho^2}{2} - \frac{\rho^4}{4} \right]_0^R \\
&= 2\pi^2 \left( \frac{R^4}{2} - \frac{R^4}{4} \right) \\
&= 2\pi^2 \left( \frac{R^4}{4} \right) = \frac{\pi^2 R^4}{2}.
\end{align*}

Este resultado demuestra que el hipervolumen de la bola de radio $R$ en $4$ dimensiones es $V_4(R) = \frac{\pi^2}{2} R^4$. En general, la fórmula se puede expresar usando la función Gamma. Para cualquier entero $n \geq 1$:
\[
V_n(R) = \frac{\pi^{n/2}}{\Gamma(\frac{n}{2} + 1)} R^n
\]
lo que coincide para $n=4$ dado que $\Gamma(3) = 2$.
\newpage
\part*{Parte VI: Banco Completo de 90 Ejercicios Graduados (Nivel 0 al Nivel 5)}
\addcontentsline{toc}{part}{Parte VI: Banco Completo de 90 Ejercicios Graduados (Nivel 0 al Nivel 5)}
En este capítulo se presenta la resolución íntegra y exhaustiva de los 90 ejercicios pedagógicos que componen la plataforma interactiva. Cada subtema contiene exactamente 18 ejercicios, estructurados en seis niveles de maestría analítica (3 problemas por nivel). No se omite ningún paso de integración, factorización o verificación de hipótesis topológicas.

\section{Tema 1: Regiones Generales en $\mathbb{R}^2$, Jordan y Fubini}

\subsection{Nivel 0: Calentamiento Conceptual}

\begin{nivelbox}{1.0.1}{Identificación de tipo de región en el plano}
\textbf{Enunciado:}\\
Sea $\Omega \subset \mathbb{R}^2$ el triángulo delimitado por las rectas $x = 0$, $y = 0$ y $x + y = 2$.
(a) Describa $\Omega$ como una región Tipo 1 ($y$ acotada por funciones de $x$).
(b) Describa $\Omega$ como una región Tipo 2 ($x$ acotada por funciones de $y$).

\begin{observacionbox}{Pista Metodológica}
Para Tipo 1, proyecte sobre el eje $X$ hallando el intervalo $[a,b]$ y luego determine los pisos y techos $y_1(x) \le y \le y_2(x)$. Para Tipo 2, proyecte sobre el eje $Y$.
\end{observacionbox}

\begin{ejerciciobox}{Resolución Analítica Paso a Paso}
\smallskip\noindent\textbf{1.} \textbf{Descripción Tipo 1:}
La proyección sobre el eje $X$ va desde $x = 0$ hasta $x = 2$. Para cada $x \in [0, 2]$, la recta vertical entra por el eje $X$ ($y = 0$) y sale por la recta $y = 2 - x$.
\[\Omega = \left\{ (x,y) \in \mathbb{R}^2 : 0 \le x \le 2, \ 0 \le y \le 2 - x \right\}.\]

\smallskip\noindent\textbf{2.} \textbf{Descripción Tipo 2:}
La proyección sobre el eje $Y$ va desde $y = 0$ hasta $y = 2$. Para cada $y \in [0, 2]$, el segmento horizontal entra por el eje $Y$ ($x = 0$) y sale por $x = 2 - y$.
\[\Omega = \left\{ (x,y) \in \mathbb{R}^2 : 0 \le y \le 2, \ 0 \le x \le 2 - y \right\}.\]
\end{ejerciciobox}

\begin{errorbox}{¡Cuidado con esta trampa de examen!}
Poner límites variables en ambas integrales a la vez. La integral exterior DEBE tener siempre límites constantes.
\end{errorbox}
\end{nivelbox}
\bigskip

\begin{nivelbox}{1.0.2}{Área de una región mediante integral doble de la función constante 1}
\textbf{Enunciado:}\\
Calcule el área de la región parabólica $\Omega = \left\{ (x,y) \in \mathbb{R}^2 : 0 \le x \le 1, \ x^2 \le y \le x \right\}$ usando $\operatorname{Area}(\Omega) = \iint_\Omega 1 \, dA$.

\begin{observacionbox}{Pista Metodológica}
Escriba la integral iterada con límites en $y$ desde $x^2$ hasta $x$, e integre luego respecto a $x$ de $0$ a $1$.
\end{observacionbox}

\begin{ejerciciobox}{Resolución Analítica Paso a Paso}
\smallskip\noindent\textbf{1.} \textbf{Planteamiento iterado:}
\[\operatorname{Area}(\Omega) = \int_0^1 \left( \int_{x^2}^x 1 \, dy \right) dx.\]

\smallskip\noindent\textbf{2.} \textbf{Integración interior:}
\[\int_{x^2}^x 1 \, dy = [y]_{x^2}^x = x - x^2.\]

\smallskip\noindent\textbf{3.} \textbf{Integración exterior:}
\[\int_0^1 (x - x^2) \, dx = \left[ \frac{x^2}{2} - \frac{x^3}{3} \right]_0^1 = \frac{1}{2} - \frac{1}{3} = \frac{1}{6}.\]
\end{ejerciciobox}

\begin{errorbox}{¡Cuidado con esta trampa de examen!}
Confundir cuál función está por encima. En $[0,1]$, $x \ge x^2$, por lo que el techo es $y = x$ y el piso $y = x^2$.
\end{errorbox}
\end{nivelbox}
\bigskip

\begin{nivelbox}{1.0.3}{Propiedad de separación para funciones producto en rectángulos}
\textbf{Enunciado:}\\
Demuestre y calcule $\iint_R x e^y \, dA$ sobre el rectángulo $R = [0, 2] \times [0, 1]$ separando el producto de integrales unidimensionales.

\begin{observacionbox}{Pista Metodológica}
Si el dominio es un rectángulo $[a,b] \times [c,d]$ y $f(x,y) = g(x)h(y)$, entonces $\iint_R f \, dA = \left(\int_a^b g(x)dx\right)\left(\int_c^d h(y)dy\right)$.
\end{observacionbox}

\begin{ejerciciobox}{Resolución Analítica Paso a Paso}
\smallskip\noindent\textbf{1.} \textbf{Factorización por Teorema de Fubini:}
\[\iint_R x e^y \, dA = \left( \int_0^2 x \, dx \right) \left( \int_0^1 e^y \, dy \right).\]

\smallskip\noindent\textbf{2.} \textbf{Cálculo directo:}
\[\int_0^2 x \, dx = \left[ \frac{x^2}{2} \right]_0^2 = 2.\]
\[\int_0^1 e^y \, dy = [e^y]_0^1 = e - 1.\]

\smallskip\noindent\textbf{3.} \textbf{Resultado:}
\[I = 2(e - 1) = 2e - 2.\]
\end{ejerciciobox}

\begin{errorbox}{¡Cuidado con esta trampa de examen!}
Intentar aplicar esta separación cuando los límites de integración son variables o la región no es un rectángulo cartesiano.
\end{errorbox}
\end{nivelbox}
\bigskip

\subsection{Nivel 1: Procedimental Guiado}

\begin{nivelbox}{1.1.1}{Integral doble con polinomio sobre región triangular}
\textbf{Enunciado:}\\
Calcule $\iint_\Omega (2x + 3y) \, dA$, donde $\Omega$ es el triángulo con vértices $(0,0)$, $(1,0)$ y $(1,2)$.

\begin{observacionbox}{Pista Metodológica}
La hipotenusa une $(0,0)$ con $(1,2)$, por lo que su ecuación es $y = 2x$. Como región Tipo 1, $0 \le x \le 1$ y $0 \le y \le 2x$.
\end{observacionbox}

\begin{ejerciciobox}{Resolución Analítica Paso a Paso}
\smallskip\noindent\textbf{1.} \textbf{Iterada Tipo 1:}
\[I = \int_0^1 \int_0^{2x} (2x + 3y) \, dy \, dx.\]

\smallskip\noindent\textbf{2.} \textbf{Integral interior:}
\[\int_0^{2x} (2x + 3y) \, dy = \left[ 2xy + \frac{3y^2}{2} \right]_0^{2x} = 2x(2x) + \frac{3(4x^2)}{2} = 4x^2 + 6x^2 = 10x^2.\]

\smallskip\noindent\textbf{3.} \textbf{Integral exterior:}
\[I = \int_0^1 10x^2 \, dx = \left[ \frac{10x^3}{3} \right]_0^1 = \frac{10}{3}.\]
\end{ejerciciobox}

\begin{errorbox}{¡Cuidado con esta trampa de examen!}
Olvidar multiplicar el término $2x$ por la longitud del intervalo $2x$ al integrar respecto a $y$.
\end{errorbox}
\end{nivelbox}
\bigskip

\begin{nivelbox}{1.1.2}{Inversión elemental de orden de integración}
\textbf{Enunciado:}\\
Invierta el orden de integración de $I = \int_0^1 \int_x^1 \cos(y^2) \, dy \, dx$ y halle su valor exacto.

\begin{observacionbox}{Pista Metodológica}
La función $\cos(y^2)$ no tiene primitiva elemental en términos de funciones estándar. Al cambiar a Tipo 2 ($dx dy$), aparecerá un factor $y$ que facilita la sustitución $u = y^2$.
\end{observacionbox}

\begin{ejerciciobox}{Resolución Analítica Paso a Paso}
\smallskip\noindent\textbf{1.} \textbf{Identificación de la región:}
\[\Omega = \left\{ (x,y) : 0 \le x \le 1, \ x \le y \le 1 \right\}.\]
En Tipo 2: $y$ varía en $[0, 1]$ y para cada $y$, $x$ entra por $x = 0$ y sale por $x = y$.

\smallskip\noindent\textbf{2.} \textbf{Integral invertida:}
\[I = \int_0^1 \int_0^y \cos(y^2) \, dx \, dy = \int_0^1 \cos(y^2) [x]_0^y \, dy = \int_0^1 y \cos(y^2) \, dy.\]

\smallskip\noindent\textbf{3.} \textbf{Sustitución simple ($u = y^2$, $du = 2y dy$):}
\[I = \frac{1}{2} [\sin(y^2)]_0^1 = \frac{\sin(1)}{2}.\]
\end{ejerciciobox}

\begin{errorbox}{¡Cuidado con esta trampa de examen!}
Cambiar el orden escribiendo simplemente $\int_x^1 \int_0^1 \dots$ dejando la variable $x$ en la integral exterior.
\end{errorbox}
\end{nivelbox}
\bigskip

\begin{nivelbox}{1.1.3}{Evaluación con función valor absoluto en integrales dobles}
\textbf{Enunciado:}\\
Calcule $I = \iint_R x |y| \, dA$, donde $R = [0, 1] \times [-1, 2]$ (Problema de la PC5 2025-II).

\begin{observacionbox}{Pista Metodológica}
Separe la región en dos rectángulos $R_1 = [0, 1] \times [-1, 0]$ donde $|y| = -y$, y $R_2 = [0, 1] \times [0, 2]$ donde $|y| = y$.
\end{observacionbox}

\begin{ejerciciobox}{Resolución Analítica Paso a Paso}
\smallskip\noindent\textbf{1.} \textbf{Descomposición por aditividad del dominio:}
\[I = \left( \int_0^1 x \, dx \right) \left( \int_{-1}^2 |y| \, dy \right).\]

\smallskip\noindent\textbf{2.} \textbf{Integral en $x$:}
\[\int_0^1 x \, dx = \frac{1}{2}.\]

\smallskip\noindent\textbf{3.} \textbf{Integral en $y$:}
\[\int_{-1}^2 |y| \, dy = \int_{-1}^0 (-y) \, dy + \int_0^2 y \, dy = \left[ -\frac{y^2}{2} \right]_{-1}^0 + \left[ \frac{y^2}{2} \right]_0^2 = \frac{1}{2} + 2 = \frac{5}{2}.\]

\smallskip\noindent\textbf{4.} \textbf{Producto final:}
\[I = \frac{1}{2} \cdot \frac{5}{2} = \frac{5}{4}.\]
\end{ejerciciobox}

\begin{errorbox}{¡Cuidado con esta trampa de examen!}
Integrar $|y|$ como si fuera $y$, obteniendo $\left[y^2/2\right]_{-1}^2 = 2 - 1/2 = 3/2$, lo cual es un error grave de signo.
\end{errorbox}
\end{nivelbox}
\bigskip

\subsection{Nivel 2: Integración de Conceptos}

\begin{nivelbox}{1.2.1}{Integral doble sobre región acotada por exponencial y rectas}
\textbf{Enunciado:}\\
Calcule $I = \iint_R e^{x+y} \, dA$, donde $R$ está limitada por las curvas $y = e^x$, $x = 0$ e $y = 2$ (PC5 2025-II, Pregunta 2a).

\begin{observacionbox}{Pista Metodológica}
Es mucho más sencillo describirla como región Tipo 2: $y$ va desde $1$ (cuando $x=0$) hasta $2$, y $x$ entra por $x = 0$ y sale por $x = \ln y$.
\end{observacionbox}

\begin{ejerciciobox}{Resolución Analítica Paso a Paso}
\smallskip\noindent\textbf{1.} \textbf{Intersección de frontera:}
Para $x = 0$, $y = e^0 = 1$. Para $y = 2$, $e^x = 2 \implies x = \ln 2$.
La región en Tipo 2 es: $1 \le y \le 2$ y $0 \le x \le \ln y$.

\smallskip\noindent\textbf{2.} \textbf{Planteamiento de la integral:}
\[I = \int_1^2 e^y \left( \int_0^{\ln y} e^x \, dx \right) dy.\]

\smallskip\noindent\textbf{3.} \textbf{Integral interior:}
\[\int_0^{\ln y} e^x \, dx = [e^x]_0^{\ln y} = y - 1.\]

\smallskip\noindent\textbf{4.} \textbf{Integral exterior por partes:}
\[I = \int_1^2 (y - 1) e^y \, dy.\]
Haciendo $u = y-1, du = dy, dv = e^y dy, v = e^y$:
\[I = \left[ (y - 1)e^y \right]_1^2 - \int_1^2 e^y \, dy = e^2 - [e^y]_1^2 = e^2 - (e^2 - e) = e.\]
\end{ejerciciobox}

\begin{errorbox}{¡Cuidado con esta trampa de examen!}
Tomar los límites en $y$ desde $0$ hasta $2$. La curva $y = e^x$ interseca a $x = 0$ en $y = 1$, nunca en $y = 0$.
\end{errorbox}
\end{nivelbox}
\bigskip

\begin{nivelbox}{1.2.2}{Inversión con función exponencial cúbica $e^{x^3}$}
\textbf{Enunciado:}\\
Evalúe la integral iterada $I = \int_0^4 \int_{\sqrt{y}}^2 e^{x^3} \, dx \, dy$ invirtiendo el orden de integración.

\begin{observacionbox}{Pista Metodológica}
Grafique $x = \sqrt{y} \iff y = x^2$ para $y \in [0, 4]$. La proyección en $X$ es $[0, 2]$.
\end{observacionbox}

\begin{ejerciciobox}{Resolución Analítica Paso a Paso}
\smallskip\noindent\textbf{1.} \textbf{Región original:}
\[\Omega = \left\{ (x,y) : 0 \le y \le 4, \ \sqrt{y} \le x \le 2 \right\}.\]

\smallskip\noindent\textbf{2.} \textbf{Descripción Tipo 1:}
\[0 \le x \le 2, \qquad 0 \le y \le x^2.\]

\smallskip\noindent\textbf{3.} \textbf{Evaluación de la integral:}
\[I = \int_0^2 \int_0^{x^2} e^{x^3} \, dy \, dx = \int_0^2 e^{x^3} [y]_0^{x^2} \, dx = \int_0^2 x^2 e^{x^3} \, dx.\]
Sustitución $u = x^3, du = 3x^2 dx$:
\[I = \frac{1}{3} \left[ e^{x^3} \right]_0^2 = \frac{e^8 - 1}{3}.\]
\end{ejerciciobox}

\begin{errorbox}{¡Cuidado con esta trampa de examen!}
Olvidar el factor $1/3$ proveniente de la derivada de $x^3$.
\end{errorbox}
\end{nivelbox}
\bigskip

\begin{nivelbox}{1.2.3}{Cálculo de volumen bajo un paraboloide hiperbólico}
\textbf{Enunciado:}\\
Halle el volumen del sólido delimitado superiormente por $z = x y$ e inferiormente por el plano $z = 0$, sobre la región $\Omega$ acotada por $y = x$ e $y = x^2$.

\begin{observacionbox}{Pista Metodológica}
Los puntos de corte son $(0,0)$ y $(1,1)$. Como región Tipo 1, $0 \le x \le 1$ y $x^2 \le y \le x$.
\end{observacionbox}

\begin{ejerciciobox}{Resolución Analítica Paso a Paso}
\smallskip\noindent\textbf{1.} \textbf{Integral de volumen:}
\[V = \int_0^1 \int_{x^2}^x x y \, dy \, dx.\]

\smallskip\noindent\textbf{2.} \textbf{Integral interior:}
\[\int_{x^2}^x x y \, dy = x \left[ \frac{y^2}{2} \right]_{x^2}^x = \frac{x}{2} (x^2 - x^4) = \frac{x^3 - x^5}{2}.\]

\smallskip\noindent\textbf{3.} \textbf{Integral exterior:}
\[V = \frac{1}{2} \int_0^1 (x^3 - x^5) \, dx = \frac{1}{2} \left[ \frac{x^4}{4} - \frac{x^6}{6} \right]_0^1 = \frac{1}{2} \left( \frac{1}{4} - \frac{1}{6} \right) = \frac{1}{2} \left( \frac{1}{12} \right) = \frac{1}{24}.\]
\end{ejerciciobox}

\begin{errorbox}{¡Cuidado con esta trampa de examen!}
Restar los polinomios al revés, obteniendo un volumen negativo.
\end{errorbox}
\end{nivelbox}
\bigskip

\subsection{Nivel 3: Problemas Tipo Examen UNI}

\begin{nivelbox}{1.3.1}{Inversión general de Fubini con función genérica continua (PC5 2024-II)}
\textbf{Enunciado:}\\
Demuestre la identidad de la PC5 2024-II (Pregunta 1a):
\[\int_0^a \left[ \int_0^y e^{m(a-x)} f(x) \, dx \right] dy = \int_0^a (a - x) e^{m(a-x)} f(x) \, dx,\]
donde $a, m$ son constantes reales con $a > 0$ y $f$ es continua en $[0, a]$.

\begin{observacionbox}{Pista Metodológica}
Dibuje el triángulo $0 \le y \le a$ y $0 \le x \le y$. Invierta el orden de integración a $0 \le x \le a$ y $x \le y \le a$.
\end{observacionbox}

\begin{ejerciciobox}{Resolución Analítica Paso a Paso}
\smallskip\noindent\textbf{1.} \textbf{Identificación de la región:}
\[R = \left\{ (x,y) \in \mathbb{R}^2 : 0 \le y \le a, \ 0 \le x \le y \right\}.\]

\smallskip\noindent\textbf{2.} \textbf{Inversión del orden:}
La proyección sobre $X$ es $x \in [0, a]$. Para cada $x$ fijado, la variable vertical $y$ va desde la recta $y = x$ hasta el techo horizontal $y = a$:
\[R = \left\{ (x,y) \in \mathbb{R}^2 : 0 \le x \le a, \ x \le y \le a \right\}.\]

\smallskip\noindent\textbf{3.} \textbf{Cálculo de la integral:}
\[I = \int_0^a \left[ \int_x^a e^{m(a-x)} f(x) \, dy \right] dx.\]
Como el integrando no depende de $y$:
\[\int_x^a e^{m(a-x)} f(x) \, dy = e^{m(a-x)} f(x) \int_x^a 1 \, dy = e^{m(a-x)} f(x) (a - x).\]

\smallskip\noindent\textbf{4.} \textbf{Conclusión:}
\[I = \int_0^a (a - x) e^{m(a-x)} f(x) \, dx. \qquad \blacksquare\]
\end{ejerciciobox}

\begin{errorbox}{¡Cuidado con esta trampa de examen!}
Intentar integrar por partes antes de invertir el orden. El cambio de orden resuelve el ejercicio en 2 líneas.
\end{errorbox}
\end{nivelbox}
\bigskip

\begin{nivelbox}{1.3.2}{Teorema Fundamental del Cálculo en derivadas cruzadas (PC5 2024-II)}
\textbf{Enunciado:}\\
Sea $f: A \subset \mathbb{R}^2 \to \mathbb{R}$ de clase $C^2$ en un abierto $A$. Si $R = [a,b] \times [c,d] \subset A$, pruebe mediante el Teorema de Fubini que:
\[\iint_R \frac{\partial^2 f}{\partial x \partial y} \, dA = f(b,d) - f(a,d) - f(b,c) + f(a,c).\]

\begin{observacionbox}{Pista Metodológica}
Escriba la integral iterada $\int_c^d \left( \int_a^b \frac{\partial}{\partial x}\left(\frac{\partial f}{\partial y}\right) dx \right) dy$ y aplique el TFC unidimensional.
\end{observacionbox}

\begin{ejerciciobox}{Resolución Analítica Paso a Paso}
\smallskip\noindent\textbf{1.} \textbf{Aplicación de Fubini:}
Por continuidad de $\frac{\partial^2 f}{\partial x \partial y}$:
\[\iint_R \frac{\partial^2 f}{\partial x \partial y} \, dA = \int_c^d \left[ \int_a^b \frac{\partial}{\partial x} \left( \frac{\partial f}{\partial y}(x,y) \right) dx \right] dy.\]

\smallskip\noindent\textbf{2.} \textbf{TFC en la variable $x$:}
\[\int_a^b \frac{\partial}{\partial x} \left( \frac{\partial f}{\partial y}(x,y) \right) dx = \frac{\partial f}{\partial y}(b, y) - \frac{\partial f}{\partial y}(a, y).\]

\smallskip\noindent\textbf{3.} \textbf{TFC en la variable $y$:}
\[\int_c^d \left[ \frac{\partial f}{\partial y}(b, y) - \frac{\partial f}{\partial y}(a, y) \right] dy = [f(b, y)]_c^d - [f(a, y)]_c^d.\]
\[= (f(b,d) - f(b,c)) - (f(a,d) - f(a,c)) = f(b,d) - f(a,d) - f(b,c) + f(a,c). \quad \blacksquare\]
\end{ejerciciobox}

\begin{errorbox}{¡Cuidado con esta trampa de examen!}
Confundir los signos de los cuatro términos. Note el patrón alternado $+ - - +$ característico de la medida del rectángulo.
\end{errorbox}
\end{nivelbox}
\bigskip

\begin{nivelbox}{1.3.3}{Derivación de la función acumuladora doble}
\textbf{Enunciado:}\\
Sea $F(x,y) = \int_a^x \int_c^y f(u,v) \, dv \, du$ con $f$ continua en un entorno de $(x,y)$. Demuestre rigurosamente que $\frac{\partial^2 F}{\partial y \partial x}(x,y) = f(x,y)$.

\begin{observacionbox}{Pista Metodológica}
Defina $G(u, y) = \int_c^y f(u,v) dv$. Derive primero respecto a $x$ usando Leibniz y luego respecto a $y$.
\end{observacionbox}

\begin{ejerciciobox}{Resolución Analítica Paso a Paso}
\smallskip\noindent\textbf{1.} \textbf{Primera derivada respecto a $x$:}
\[F(x,y) = \int_a^x G(u, y) \, du, \qquad \text{donde } G(u,y) = \int_c^y f(u,v) \, dv.\]
Por el Teorema Fundamental del Cálculo aplicado a la variable $x$:
\[\frac{\partial F}{\partial x}(x,y) = G(x, y) = \int_c^y f(x, v) \, dv.\]

\smallskip\noindent\textbf{2.} \textbf{Segunda derivada respecto a $y$:}
Ahora derivamos respecto a $y$. Nuevamente por el TFC unidimensional:
\[\frac{\partial^2 F}{\partial y \partial x}(x,y) = \frac{\partial}{\partial y} \left( \int_c^y f(x, v) \, dv \right) = f(x, y). \qquad \blacksquare\]
\end{ejerciciobox}

\begin{errorbox}{¡Cuidado con esta trampa de examen!}
Asumir que se puede conmutar el orden de derivadas sin justificar que $f(x,y)$ es continua.
\end{errorbox}
\end{nivelbox}
\bigskip

\subsection{Nivel 4: Desafíos Teóricos y Contraejemplos}

\begin{nivelbox}{1.4.1}{Demostración de identidad integral por simetría cuadrática (PC5 2025-II)}
\textbf{Enunciado:}\\
Sea $f: [a,b] \to \mathbb{R}$ continua. Demuestre formalmente que:
\[\int_a^b \int_x^b f(x)f(y) \, dy \, dx = \frac{1}{2} \left( \int_a^b f(x) \, dx \right)^2.\]

\begin{observacionbox}{Pista Metodológica}
Considere el cuadrado $Q = [a,b] \times [a,b]$ y divídalo por la diagonal $y = x$ en dos triángulos simétricos $T_1$ y $T_2$.
\end{observacionbox}

\begin{ejerciciobox}{Resolución Analítica Paso a Paso}
\smallskip\noindent\textbf{1.} \textbf{Integral sobre el cuadrado completo:}
\[I_{\text{cuadrado}} = \int_a^b \int_a^b f(x)f(y) \, dy \, dx = \left( \int_a^b f(x) \, dx \right) \left( \int_a^b f(y) \, dy \right) = \left( \int_a^b f(x) \, dx \right)^2.\]

\smallskip\noindent\textbf{2.} \textbf{Descomposición por la diagonal $y = x$:}
El cuadrado $Q$ es la unión casi disjunta de:
\[T_1 = \{ (x,y) : a \le x \le b, \ x \le y \le b \}, \qquad T_2 = \{ (x,y) : a \le y \le b, \ y \le x \le b \}.\]
Por aditividad de la integral: $I_{\text{cuadrado}} = \iint_{T_1} f(x)f(y) \, dA + \iint_{T_2} f(x)f(y) \, dA$.

\smallskip\noindent\textbf{3.} \textbf{Igualdad por simetría de variables mudas:}
En $\iint_{T_2} f(x)f(y) \, dA$, intercambiando los nombres de las variables $x \leftrightarrow y$:
\[\iint_{T_2} f(x)f(y) \, dx dy = \iint_{T_1} f(y)f(x) \, dy dx = \iint_{T_1} f(x)f(y) \, dA.\]

\smallskip\noindent\textbf{4.} \textbf{Conclusión:}
\[2 \iint_{T_1} f(x)f(y) \, dA = \left( \int_a^b f(x) \, dx \right)^2 \implies \int_a^b \int_x^b f(x)f(y) \, dy \, dx = \frac{1}{2} \left( \int_a^b f(x) \, dx \right)^2. \quad \blacksquare\]
\end{ejerciciobox}

\begin{errorbox}{¡Cuidado con esta trampa de examen!}
Intentar calcular primitivas algebraicas explícitas de una función arbitraria $f$, en lugar de explotar la simetría geométrica del dominio.
\end{errorbox}
\end{nivelbox}
\bigskip

\begin{nivelbox}{1.4.2}{Discontinuidades en un conjunto de contenido nulo}
\textbf{Enunciado:}\\
Sea $f: [0,1] \times [0,1] \to \mathbb{R}$ definida por $f(x,y) = 1$ si $y > x^2$, y $f(x,y) = 0$ si $y \le x^2$. Demuestre que $f$ es integrable según Riemann y calcule su integral doble.

\begin{observacionbox}{Pista Metodológica}
El conjunto de discontinuidades de $f$ es exactamente la parábola $y = x^2$. Demuestre que una curva suave en $\mathbb{R}^2$ tiene contenido de Jordan nulo.
\end{observacionbox}

\begin{ejerciciobox}{Resolución Analítica Paso a Paso}
\smallskip\noindent\textbf{1.} \textbf{Conjunto de discontinuidades:}
$D_f = \{ (x, x^2) \in \mathbb{R}^2 : x \in [0, 1] \}$.

\smallskip\noindent\textbf{2.} \textbf{Prueba de contenido nulo:}
Como $g(x) = x^2$ es continua en $[0,1]$, es uniformemente continua: para todo $\varepsilon > 0$, existe $\delta > 0$ tal que $|x_1 - x_2| < \delta \implies |x_1^2 - x_2^2| < \varepsilon$.
Tomando una partición con $n > 1/\delta$ rectángulos, podemos cubrir la gráfica con rectángulos cuya suma total de áreas es menor que $\varepsilon$. Por tanto, $\operatorname{cont}(D_f) = 0$.

\smallskip\noindent\textbf{3.} \textbf{Criterio de Lebesgue-Jordan:}
Una función acotada en un rectángulo cuyas discontinuidades forman un conjunto de contenido nulo es integrable Riemann.

\smallskip\noindent\textbf{4.} \textbf{Cálculo de la integral por Fubini:}
\[\iint_{[0,1]^2} f \, dA = \int_0^1 \left( \int_{x^2}^1 1 \, dy \right) dx = \int_0^1 (1 - x^2) \, dx = \left[ x - \frac{x^3}{3} \right]_0^1 = \frac{2}{3}.\]
\end{ejerciciobox}

\begin{errorbox}{¡Cuidado con esta trampa de examen!}
Creer que cualquier función discontinua no es integrable Riemann. La integrabilidad se preserva si las discontinuidades tienen contenido nulo.
\end{errorbox}
\end{nivelbox}
\bigskip

\begin{nivelbox}{1.4.3}{Contraejemplo de Fubini para función no integrable}
\textbf{Enunciado:}\\
Considere la función $f(x,y) = \frac{x^2 - y^2}{(x^2 + y^2)^2}$ en $(0,1] \times (0,1]$. Calcule las dos integrales iteradas $\int_0^1 \left(\int_0^1 f(x,y) dy\right) dx$ y $\int_0^1 \left(\int_0^1 f(x,y) dx\right) dy$. Explique por qué sus valores difieren y por qué esto no contradice el Teorema de Fubini.

\begin{observacionbox}{Pista Metodológica}
Note que $\frac{\partial}{\partial y}\left(\frac{y}{x^2+y^2}\right) = \frac{x^2-y^2}{(x^2+y^2)^2}$.
\end{observacionbox}

\begin{ejerciciobox}{Resolución Analítica Paso a Paso}
\smallskip\noindent\textbf{1.} \textbf{Primera iterada:}
\[\int_0^1 \frac{x^2 - y^2}{(x^2 + y^2)^2} \, dy = \left[ \frac{y}{x^2 + y^2} \right]_{y=0}^{y=1} = \frac{1}{x^2 + 1}.\]
Integrando respecto a $x$:
\[\int_0^1 \frac{1}{x^2 + 1} \, dx = [\arctan(x)]_0^1 = \frac{\pi}{4}.\]

\smallskip\noindent\textbf{2.} \textbf{Segunda iterada:}
Por antisimetría $f(y,x) = -f(x,y)$:
\[\int_0^1 \left( \int_0^1 \frac{x^2 - y^2}{(x^2 + y^2)^2} \, dx \right) dy = -\frac{\pi}{4}.\]

\smallskip\noindent\textbf{3.} \textbf{Justificación:}
Las integrales iteradas dan $\frac{\pi}{4} \neq -\frac{\pi}{4}$. Esto NO contradice el Teorema de Fubini porque la función $f(x,y)$ \textbf{no está acotada} cerca del origen $(0,0)$ y por tanto la integral doble de Riemann no existe (no es integrable en sentido ordinario).
\end{ejerciciobox}

\begin{errorbox}{¡Cuidado con esta trampa de examen!}
Concluir que Fubini es falso, olvidando verificar la hipótesis crítica de que la función debe ser acotada e integrable.
\end{errorbox}
\end{nivelbox}
\bigskip

\subsection{Nivel 5: Dominio Avanzado y Rigor Topológico}

\begin{nivelbox}{1.5.1}{Teorema de Darboux: Suma inferior menor o igual a suma superior arbitraria}
\textbf{Enunciado:}\\
Sea $f: B \subset \mathbb{R}^n \to \mathbb{R}$ acotada en un bloque $B$. Demuestre que para cualesquiera dos particiones $P$ y $Q$ de $B$, se cumple rigurosamente que $s(f, P) \le S(f, Q)$ (Pregunta 1c de la PC5 2024-II).

\begin{observacionbox}{Pista Metodológica}
Utilice la partición refinamiento común $P + Q = P \cup Q$ y las propiedades de monotonicidad de sumas inferiores y superiores frente al refinamiento.
\end{observacionbox}

\begin{ejerciciobox}{Resolución Analítica Paso a Paso}
\smallskip\noindent\textbf{1.} \textbf{Refinamiento común:}
Sea $P^* = P \cup Q$ el refinamiento común de las particiones $P$ y $Q$. Es decir, $P \subset P^*$ y $Q \subset P^*$.

\smallskip\noindent\textbf{2.} \textbf{Monotonicidad de las sumas de Darboux:}
- Al añadir puntos a una partición, la suma inferior no decrece: $s(f, P) \le s(f, P^*)$.
- Al añadir puntos a una partición, la suma superior no aumenta: $S(f, P^*) \le S(f, Q)$.

\smallskip\noindent\textbf{3.} \textbf{Relación básica sobre una misma partición:}
Para cualquier partición fijada, en particular $P^*$, el ínfimo en cada sub-bloque es menor o igual al supremo: $m_i \le M_i$, luego:
\[s(f, P^*) \le S(f, P^*).\]

\smallskip\noindent\textbf{4.} \textbf{Cadena de desigualdades:}
\[s(f, P) \le s(f, P^*) \le S(f, P^*) \le S(f, Q).\]
Por transitividad en $\mathbb{R}$: $s(f, P) \le S(f, Q)$. $\blacksquare$
\end{ejerciciobox}

\begin{errorbox}{¡Cuidado con esta trampa de examen!}
Intentar comparar directamente los rectángulos de $P$ con los de $Q$ sin introducir la partición común $P \cup Q$.
\end{errorbox}
\end{nivelbox}
\bigskip

\begin{nivelbox}{1.5.2}{Frontera con contenido nulo y medida de Jordan}
\textbf{Enunciado:}\\
Demuestre formalmente que un conjunto acotado $\Omega \subset \mathbb{R}^2$ es medible según Jordan si y solo si su frontera topológica $\partial\Omega$ tiene contenido de Jordan nulo.

\begin{observacionbox}{Pista Metodológica}
Considere la función característica $\chi_\Omega$ en un bloque contenedor $B \supset \Omega$. Analice dónde $\chi_\Omega$ es discontinua.
\end{observacionbox}

\begin{ejerciciobox}{Resolución Analítica Paso a Paso}
\smallskip\noindent\textbf{1.} \textbf{Caracterización de discontinuidades:}
Sea $\chi_\Omega: B \to \mathbb{R}$ la función característica de $\Omega$.
- Si $x \in \operatorname{int}(\Omega)$, existe una vecindad $V \subset \Omega$ donde $\chi_\Omega(y) = 1$ para todo $y \in V$, por lo que $\chi_\Omega$ es continua en $x$.
- Si $x \in \operatorname{ext}(\Omega) = B \setminus \overline{\Omega}$, existe una vecindad donde $\chi_\Omega(y) = 0$, luego es continua en $x$.
- Si $x \in \partial\Omega$, en toda vecindad de $x$ existen puntos de $\Omega$ (donde $\chi_\Omega = 1$) y puntos fuera de $\Omega$ (donde $\chi_\Omega = 0$). Por tanto, la oscilación es $\omega(x) = 1 > 0$, lo que implica que $\chi_\Omega$ es discontinua en todo punto de $\partial\Omega$.

\smallskip\noindent\textbf{2.} \textbf{Criterio de Riemann-Lebesgue:}
Por definición, $\Omega$ es Jordan-medible $\iff \chi_\Omega$ es integrable Riemann en $B$.
Una función acotada es integrable Riemann $\iff$ su conjunto de discontinuidades tiene contenido nulo.
Como $D_{\chi_\Omega} = \partial\Omega$, esto ocurre si y solo si $\operatorname{cont}(\partial\Omega) = 0$. $\blacksquare$
\end{ejerciciobox}

\begin{errorbox}{¡Cuidado con esta trampa de examen!}
Confundir contenido nulo (un número finito de cubrimientos abiertos) con medida de Lebesgue cero (cubrimientos numerables). Para integrales de Riemann se exige contenido nulo.
\end{errorbox}
\end{nivelbox}
\bigskip

\begin{nivelbox}{1.5.3}{La función de Dirichlet bidimensional no es integrable}
\textbf{Enunciado:}\\
Sea $B = [0,1] \times [0,1]$ y defina $f: B \to \mathbb{R}$ como $f(x,y) = 1$ si $(x,y) \in \mathbb{Q}^2$, y $f(x,y) = 0$ si $(x,y) \notin \mathbb{Q}^2$. Demuestre rigurosamente que $f$ no es integrable según Riemann y que las funciones integrables no son necesariamente continuas (PC5 2025-II y 2024-II).

\begin{observacionbox}{Pista Metodológica}
Utilice la densidad de $\mathbb{Q}^2$ y $(\mathbb{R} \setminus \mathbb{Q})^2$ para evaluar los ínfimos y supremos en cualquier sub-rectángulo.
\end{observacionbox}

\begin{ejerciciobox}{Resolución Analítica Paso a Paso}
\smallskip\noindent\textbf{1.} \textbf{Densidad en sub-rectángulos:}
Para cualquier partición $P$ de $B$, cada sub-rectángulo $B_{ij}$ es un producto de intervalos no degenerados $[x_{i-1}, x_i] \times [y_{j-1}, y_j]$.
Por densidad de racionales en $\mathbb{R}$, existe $(u,v) \in B_{ij} \cap \mathbb{Q}^2$, donde $f(u,v) = 1$. Por ende, $M_{ij} = \sup_{(x,y) \in B_{ij}} f(x,y) = 1$.
Por densidad de irracionales, existe $(s,t) \in B_{ij} \setminus \mathbb{Q}^2$, donde $f(s,t) = 0$. Por ende, $m_{ij} = \inf_{(x,y) \in B_{ij}} f(x,y) = 0$.

\smallskip\noindent\textbf{2.} \textbf{Sumas de Darboux:}
\[S(f, P) = \sum_{i,j} M_{ij} \operatorname{vol}(B_{ij}) = 1 \cdot \sum_{i,j} \operatorname{vol}(B_{ij}) = \operatorname{vol}(B) = 1.\]
\[s(f, P) = \sum_{i,j} m_{ij} \operatorname{vol}(B_{ij}) = 0 \cdot \sum_{i,j} \operatorname{vol}(B_{ij}) = 0.\]

\smallskip\noindent\textbf{3.} \textbf{Conclusión de no integrabilidad:}
La integral superior es $\overline{\int} f = \inf_P S(f, P) = 1$, mientras que la inferior es $\underline{\int} f = \sup_P s(f, P) = 0$.
Como $\underline{\int} f \neq \overline{\int} f$, la función de Dirichlet no es integrable Riemann.

\smallskip\noindent\textbf{4.} \textbf{Contraejemplo a 'integrable $\implies$ continua':}
Una función escalonada simple $g(x,y) = 1$ en $[0, 1/2] \times [0,1]$ y $0$ en $(1/2, 1] \times [0,1]$ tiene discontinuidades en la recta $x = 1/2$ (contenido nulo), por lo que es integrable Riemann pero claramente no es continua. $\blacksquare$
\end{ejerciciobox}

\begin{errorbox}{¡Cuidado con esta trampa de examen!}
Intentar evaluar la integral calculando límites puntuales en vez de usar la definición formal de sumas superior e inferior de Darboux.
\end{errorbox}
\end{nivelbox}
\bigskip


\section{Tema 2: Teorema de Cambio de Variables y Coordenadas Polares}

\subsection{Nivel 0: Calentamiento Conceptual}

\begin{nivelbox}{2.0.1}{Jacobiano de la transformación a coordenadas polares}
\textbf{Enunciado:}\\
Considere la transformación a coordenadas polares $T(r, \theta) = (r\cos\theta, r\sin\theta)$.
(a) Calcule la matriz Jacobiana $DT(r, \theta)$.
(b) Demuestre que su determinante Jacobiano es $\det DT(r, \theta) = r$.

\begin{observacionbox}{Pista Metodológica}
Escriba $x = r\cos\theta$ e $y = r\sin\theta$. Calcule las 4 derivadas parciales y el determinante $2 \times 2$.
\end{observacionbox}

\begin{ejerciciobox}{Resolución Analítica Paso a Paso}
\smallskip\noindent\textbf{1.} \textbf{Derivadas parciales:}
\[\frac{\partial x}{\partial r} = \cos\theta, \quad \frac{\partial x}{\partial \theta} = -r\sin\theta.\]
\[\frac{\partial y}{\partial r} = \sin\theta, \quad \frac{\partial y}{\partial \theta} = r\cos\theta.\]

\smallskip\noindent\textbf{2.} \textbf{Matriz Jacobiana:}
\[DT(r,\theta) = \begin{pmatrix} \cos\theta & -r\sin\theta \\ \sin\theta & r\cos\theta \end{pmatrix}.\]

\smallskip\noindent\textbf{3.} \textbf{Determinante:}
\[\det DT = (\cos\theta)(r\cos\theta) - (-r\sin\theta)(\sin\theta) = r\cos^2\theta + r\sin^2\theta = r(\cos^2\theta + \sin^2\theta) = r.\]
Por tanto, el diferencial de área es $dA = r \, dr \, d\theta$.
\end{ejerciciobox}

\begin{errorbox}{¡Cuidado con esta trampa de examen!}
Olvidar el signo negativo en $\frac{\partial}{\partial \theta}(r\cos\theta) = -r\sin\theta$.
\end{errorbox}
\end{nivelbox}
\bigskip

\begin{nivelbox}{2.0.2}{Descripción de un sector circular en coordenadas polares}
\textbf{Enunciado:}\\
Describa en coordenadas polares el disco de radio 2 centrado en el origen en el primer cuadrante $D = \left\{ (x,y) \in \mathbb{R}^2 : x^2 + y^2 \le 4, \ x \ge 0, \ y \ge 0 \right\}$.

\begin{observacionbox}{Pista Metodológica}
Sustituya $x^2 + y^2 = r^2$. El primer cuadrante corresponde a $\theta \in [0, \pi/2]$.
\end{observacionbox}

\begin{ejerciciobox}{Resolución Analítica Paso a Paso}
\smallskip\noindent\textbf{1.} \textbf{Radio:} $r^2 \le 4 \implies 0 \le r \le 2$.
\smallskip\noindent\textbf{2.} \textbf{Ángulo:} El primer cuadrante va desde el semieje positivo $X$ ($\theta = 0$) hasta el semieje positivo $Y$ ($\theta = \pi/2$).
\smallskip\noindent\textbf{3.} \textbf{Región en polares:}
\[D' = [0, 2] \times [0, \pi/2] = \left\{ (r,\theta) : 0 \le r \le 2, \ 0 \le \theta \le \frac{\pi}{2} \right\}.\]
\end{ejerciciobox}

\begin{errorbox}{¡Cuidado con esta trampa de examen!}
Tomar el ángulo hasta $\pi$ o $2\pi$, cubriendo más cuadrantes de los solicitados.
\end{errorbox}
\end{nivelbox}
\bigskip

\begin{nivelbox}{2.0.3}{Integral directa en polares de una función con simetría radial}
\textbf{Enunciado:}\\
Evalúe $\iint_D (x^2 + y^2) \, dA$, donde $D$ es el disco $x^2 + y^2 \le 1$.

\begin{observacionbox}{Pista Metodológica}
Reemplace $x^2 + y^2 = r^2$ y $dA = r \, dr \, d\theta$.
\end{observacionbox}

\begin{ejerciciobox}{Resolución Analítica Paso a Paso}
\smallskip\noindent\textbf{1.} \textbf{Conversión a polares:}
\[\iint_D (x^2 + y^2) \, dA = \int_0^{2\pi} \int_0^1 (r^2) \cdot r \, dr \, d\theta = \int_0^{2\pi} \int_0^1 r^3 \, dr \, d\theta.\]

\smallskip\noindent\textbf{2.} \textbf{Separación de integrales:}
\[\left( \int_0^{2\pi} 1 \, d\theta \right) \left( \int_0^1 r^3 \, dr \right) = (2\pi) \left[ \frac{r^4}{4} \right]_0^1 = 2\pi \cdot \frac{1}{4} = \frac{\pi}{2}.\]
\end{ejerciciobox}

\begin{errorbox}{¡Cuidado con esta trampa de examen!}
Olvidar el factor jacobiano $r$, integrando $r^2$ en lugar de $r^3$.
\end{errorbox}
\end{nivelbox}
\bigskip

\subsection{Nivel 1: Procedimental Guiado}

\begin{nivelbox}{2.1.1}{Integral de Gauss bidimensional en el primer cuadrante}
\textbf{Enunciado:}\\
Calcule $I = \iint_Q e^{-(x^2 + y^2)} \, dA$, donde $Q$ es el primer cuadrante $x \ge 0, y \ge 0$. Concluya el valor de $\int_0^\infty e^{-x^2} dx$.

\begin{observacionbox}{Pista Metodológica}
En polares, $r \in [0, \infty)$ y $\theta \in [0, \pi/2]$.
\end{observacionbox}

\begin{ejerciciobox}{Resolución Analítica Paso a Paso}
\smallskip\noindent\textbf{1.} \textbf{Integral en polares:}
\[I = \int_0^{\pi/2} \int_0^\infty e^{-r^2} r \, dr \, d\theta.\]

\smallskip\noindent\textbf{2.} \textbf{Integración en $r$:}
\[\int_0^\infty e^{-r^2} r \, dr = \left[ -\frac{e^{-r^2}}{2} \right]_0^\infty = 0 - \left(-\frac{1}{2}\right) = \frac{1}{2}.\]

\smallskip\noindent\textbf{3.} \textbf{Integración en $\theta$:}
\[I = \frac{\pi}{2} \cdot \frac{1}{2} = \frac{\pi}{4}.\]

\smallskip\noindent\textbf{4.} \textbf{Deducción de la integral de Gauss:}
Como $Q = [0, \infty) \times [0, \infty)$, por Fubini $I = \left( \int_0^\infty e^{-x^2} dx \right)^2 = \frac{\pi}{4} \implies \int_0^\infty e^{-x^2} dx = \frac{\sqrt{\pi}}{2}$.
\end{ejerciciobox}

\begin{errorbox}{¡Cuidado con esta trampa de examen!}
Integrar respecto a $\theta$ en $[0, 2\pi]$ cuando solo se pedía el primer cuadrante.
\end{errorbox}
\end{nivelbox}
\bigskip

\begin{nivelbox}{2.1.2}{Área de un anillo elíptico mediante coordenadas elípticas}
\textbf{Enunciado:}\\
Sea $\Omega$ la región delimitada por las elipses $\frac{x^2}{a^2} + \frac{y^2}{b^2} = 1$ y $\frac{x^2}{a^2} + \frac{y^2}{b^2} = 4$. Calcule su área mediante el cambio de variables $x = a r \cos\theta, y = b r \sin\theta$.

\begin{observacionbox}{Pista Metodológica}
Calcule primero el determinante Jacobiano del cambio de coordenadas elípticas.
\end{observacionbox}

\begin{ejerciciobox}{Resolución Analítica Paso a Paso}
\smallskip\noindent\textbf{1.} \textbf{Jacobiano de la transformación:}
\[DT = \begin{pmatrix} a\cos\theta & -ar\sin\theta \\ b\sin\theta & br\cos\theta \end{pmatrix} \implies \det DT = ab r.\]

\smallskip\noindent\textbf{2.} \textbf{Límites de integración:}
Las elipses corresponden a $r = 1$ y $r = 2$, con $\theta \in [0, 2\pi]$.

\smallskip\noindent\textbf{3.} \textbf{Cálculo del área:}
\[\operatorname{Area}(\Omega) = \int_0^{2\pi} \int_1^2 ab r \, dr \, d\theta = 2\pi ab \left[ \frac{r^2}{2} \right]_1^2 = 2\pi ab \left( 2 - \frac{1}{2} \right) = 3\pi ab.\]
\end{ejerciciobox}

\begin{errorbox}{¡Cuidado con esta trampa de examen!}
Olvidar el producto de los semiejes $ab$ en el determinante jacobiano.
\end{errorbox}
\end{nivelbox}
\bigskip

\begin{nivelbox}{2.1.3}{Cambio lineal en un paralelogramo}
\textbf{Enunciado:}\\
Calcule $\iint_R (x + y) \, dA$, donde $R$ es el paralelogramo con vértices $(0,0), (2,1), (3,3), (1,2)$, acotado por las rectas $x - 2y = 0, x - 2y = -3, y - x = 0, y - x = 1$.

\begin{observacionbox}{Pista Metodológica}
Defina $u = x - 2y$ y $v = y - x$. Calcule $\frac{\partial(u,v)}{\partial(x,y)}$ e inviértalo para hallar $\det J$.
\end{observacionbox}

\begin{ejerciciobox}{Resolución Analítica Paso a Paso}
\smallskip\noindent\textbf{1.} \textbf{Jacobiano directo:}
\[\frac{\partial(u,v)}{\partial(x,y)} = \det \begin{pmatrix} 1 & -2 \\ -1 & 1 \end{pmatrix} = 1 - 2 = -1.\]
Por ende, $|\det J| = \left| \frac{1}{-1} \right| = 1$.

\smallskip\noindent\textbf{2.} \textbf{Expresión del integrando:}
Sumando las ecuaciones: $u + v = -y \implies y = -(u+v)$. Luego $x = u + 2y = u - 2u - 2v = -u - 2v$.
Entonces $x + y = -2u - 3v$.

\smallskip\noindent\textbf{3.} \textbf{Límites:} $u \in [-3, 0]$ y $v \in [0, 1]$.

\smallskip\noindent\textbf{4.} \textbf{Integración:}
\[\int_0^1 \int_{-3}^0 (-2u - 3v) \cdot 1 \, du \, dv = \int_0^1 \left[ -u^2 - 3uv \right]_{-3}^0 dv = \int_0^1 (9 - 9v) \, dv = \left[ 9v - \frac{9v^2}{2} \right]_0^1 = \frac{9}{2}.\]
\end{ejerciciobox}

\begin{errorbox}{¡Cuidado con esta trampa de examen!}
Usar el Jacobiano con signo negativo, lo que alteraría el signo de la medida del área.
\end{errorbox}
\end{nivelbox}
\bigskip

\subsection{Nivel 2: Integración de Conceptos}

\begin{nivelbox}{2.2.1}{Círculo desplazado en coordenadas polares}
\textbf{Enunciado:}\\
Calcule $\iint_D x \, dA$, donde $D$ es el círculo $(x - 1)^2 + y^2 \le 1$.

\begin{observacionbox}{Pista Metodológica}
En coordenadas polares, $x^2 - 2x + y^2 \le 0 \iff r^2 \le 2r\cos\theta \iff r \le 2\cos\theta$. El disco está en $-\pi/2 \le \theta \le \pi/2$.
\end{observacionbox}

\begin{ejerciciobox}{Resolución Analítica Paso a Paso}
\smallskip\noindent\textbf{1.} \textbf{Frontera en polares:}
$r = 2\cos\theta$. Para que $r \ge 0$, requerimos $\cos\theta \ge 0 \implies -\pi/2 \le \theta \le \pi/2$.

\smallskip\noindent\textbf{2.} \textbf{Integral en polares ($x = r\cos\theta, dA = r dr d\theta$):}
\[I = \int_{-\pi/2}^{\pi/2} \int_0^{2\cos\theta} (r\cos\theta) r \, dr \, d\theta = \int_{-\pi/2}^{\pi/2} \cos\theta \left[ \frac{r^3}{3} \right]_0^{2\cos\theta} d\theta.\]

\smallskip\noindent\textbf{3.} \textbf{Sustitución trigonométrica:}
\[I = \frac{8}{3} \int_{-\pi/2}^{\pi/2} \cos^4\theta \, d\theta = \frac{16}{3} \int_0^{\pi/2} \cos^4\theta \, d\theta.\]
Usando la fórmula de Wallis $\int_0^{\pi/2} \cos^4\theta d\theta = \frac{3}{4} \cdot \frac{1}{2} \cdot \frac{\pi}{2} = \frac{3\pi}{16}$:
\[I = \frac{16}{3} \cdot \frac{3\pi}{16} = \pi.\]
\end{ejerciciobox}

\begin{errorbox}{¡Cuidado con esta trampa de examen!}
Integrar de $0$ a $2\pi$, lo que recorrería el círculo dos veces o generaría radios negativos sin justificar.
\end{errorbox}
\end{nivelbox}
\bigskip

\begin{nivelbox}{2.2.2}{Región acotada por hipérbolas y rectas (coordenadas hiperbólicas)}
\textbf{Enunciado:}\\
Evalúe $\iint_R xy \, dA$, donde $R$ es la región en el primer cuadrante limitada por $xy = 1, xy = 3, y = x, y = 4x$.

\begin{observacionbox}{Pista Metodológica}
Efectúe el cambio $u = xy$ y $v = y/x$.
\end{observacionbox}

\begin{ejerciciobox}{Resolución Analítica Paso a Paso}
\smallskip\noindent\textbf{1.} \textbf{Jacobiano:}
\[\frac{\partial(u,v)}{\partial(x,y)} = \det \begin{pmatrix} y & x \\ -y/x^2 & 1/x \end{pmatrix} = \frac{y}{x} - \left(-\frac{y}{x}\right) = \frac{2y}{x} = 2v.\]
Por ende, $|\det J| = \left| \frac{\partial(x,y)}{\partial(u,v)} \right| = \frac{1}{2v}$.

\smallskip\noindent\textbf{2.} \textbf{Límites de integración:}
$u \in [1, 3]$ y $v \in [1, 4]$.

\smallskip\noindent\textbf{3.} \textbf{Cálculo de la integral:}
\[\iint_R xy \, dA = \int_1^4 \int_1^3 u \cdot \frac{1}{2v} \, du \, dv = \frac{1}{2} \left( \int_1^3 u \, du \right) \left( \int_1^4 \frac{1}{v} \, dv \right).\]
\[= \frac{1}{2} \left[ \frac{u^2}{2} \right]_1^3 [\ln v]_1^4 = \frac{1}{2} (4) (\ln 4) = 2 \ln 4 = 4 \ln 2.\]
\end{ejerciciobox}

\begin{errorbox}{¡Cuidado con esta trampa de examen!}
Olvidar invertir el determinante al despejar $dA = \left|\frac{\partial(x,y)}{\partial(u,v)}\right| du dv$.
\end{errorbox}
\end{nivelbox}
\bigskip

\begin{nivelbox}{2.2.3}{Integral de Fresnel / Dirichlet con exponente cuadrático}
\textbf{Enunciado:}\\
Evalúe la integral impropia $I = \iint_{\mathbb{R}^2} \frac{1}{(1 + x^2 + y^2)^2} \, dA$ pasando a coordenadas polares.

\begin{observacionbox}{Pista Metodológica}
El dominio cubre todo $\mathbb{R}^2$, luego $r \in [0, \infty)$ y $\theta \in [0, 2\pi]$.
\end{observacionbox}

\begin{ejerciciobox}{Resolución Analítica Paso a Paso}
\smallskip\noindent\textbf{1.} \textbf{Conversión a polares:}
\[I = \int_0^{2\pi} \int_0^\infty \frac{r}{(1 + r^2)^2} \, dr \, d\theta.\]

\smallskip\noindent\textbf{2.} \textbf{Integral en $r$ (sustitución $w = 1 + r^2, dw = 2r dr$):}
\[\int_0^\infty \frac{r}{(1 + r^2)^2} \, dr = \left[ -\frac{1}{2(1 + r^2)} \right]_0^\infty = 0 - \left( -\frac{1}{2} \right) = \frac{1}{2}.\]

\smallskip\noindent\textbf{3.} \textbf{Integral angular:}
\[I = 2\pi \cdot \frac{1}{2} = \pi.\]
\end{ejerciciobox}

\begin{errorbox}{¡Cuidado con esta trampa de examen!}
Dejar el límite superior de $r$ en 1 pensando que se trata del disco unitario.
\end{errorbox}
\end{nivelbox}
\bigskip

\subsection{Nivel 3: Problemas Tipo Examen UNI}

\begin{nivelbox}{2.3.1}{Volumen de la guía del embrague Ferrari (PC5 2024-II, Pregunta 3)}
\textbf{Enunciado:}\\
Se quiere construir una guía del embrague de tracción de un auto Ferrari que tiene la forma del sólido limitado por el paraboloide $z = 4 - x^2 - y^2$, encima del plano $XY$, intersecado por el cilindro hueco $(x-1)^2 + y^2 = 1$ con base en el plano $XY$. Calcule el volumen de dicha pieza usando coordenadas polares.

\begin{observacionbox}{Pista Metodológica}
En polares, el cilindro es $r = 2\cos\theta$ para $\theta \in [0, \pi]$ (o $[-\pi/2, \pi/2]$) y el paraboloide es $z = 4 - r^2$.
\end{observacionbox}

\begin{ejerciciobox}{Resolución Analítica Paso a Paso}
\smallskip\noindent\textbf{1.} \textbf{Ecuaciones en coordenadas polares:}
Cilindro: $x^2 - 2x + y^2 = 0 \implies r^2 = 2r\cos\theta \implies r = 2\cos\theta$.
Paraboloide: $z = 4 - (x^2+y^2) = 4 - r^2$.

\smallskip\noindent\textbf{2.} \textbf{Región de integración:}
$R = \left\{ (r,\theta) : -\frac{\pi}{2} \le \theta \le \frac{\pi}{2}, \ 0 \le r \le 2\cos\theta \right\}$. Por simetría respecto al eje $X$, duplicamos la integral para $\theta \in [0, \pi/2]$:
\[V = 2 \int_0^{\pi/2} \int_0^{2\cos\theta} (4 - r^2) r \, dr \, d\theta.\]

\smallskip\noindent\textbf{3.} \textbf{Integración en $r$:}
\[\int_0^{2\cos\theta} (4r - r^3) \, dr = \left[ 2r^2 - \frac{r^4}{4} \right]_0^{2\cos\theta} = 2(4\cos^2\theta) - \frac{16\cos^4\theta}{4} = 8\cos^2\theta - 4\cos^4\theta.\]

\smallskip\noindent\textbf{4.} \textbf{Integración en $\theta$:}
\[V = 2 \int_0^{\pi/2} (8\cos^2\theta - 4\cos^4\theta) \, d\theta.\]
Usando Wallis: $\int_0^{\pi/2} \cos^2\theta d\theta = \frac{\pi}{4}$, $\int_0^{\pi/2} \cos^4\theta d\theta = \frac{3\pi}{16}$.
\[V = 2 \left( 8 \cdot \frac{\pi}{4} - 4 \cdot \frac{3\pi}{16} \right) = 2 \left( 2\pi - \frac{3\pi}{4} \right) = 2 \left( \frac{5\pi}{4} \right) = \frac{5\pi}{2}.\]
\end{ejerciciobox}

\begin{errorbox}{¡Cuidado con esta trampa de examen!}
Integrar $r$ hasta $2$ en vez de la frontera móvil $r = 2\cos\theta$, confundiendo el cilindro desplazado con el cilindro centrado.
\end{errorbox}
\end{nivelbox}
\bigskip

\begin{nivelbox}{2.3.2}{Integral sobre cardioide en coordenadas polares}
\textbf{Enunciado:}\\
Halle el área encerrada por la cardioide $r = 1 + \cos\theta$ y calcule $\iint_\Omega x \, dA$.

\begin{observacionbox}{Pista Metodológica}
Para el área, integre $\frac{1}{2} r^2 d\theta$. Para $\iint x dA$, reemplace $x = r\cos\theta$.
\end{observacionbox}

\begin{ejerciciobox}{Resolución Analítica Paso a Paso}
\smallskip\noindent\textbf{1.} \textbf{Área:}
\[\operatorname{Area} = \int_0^{2\pi} \int_0^{1+\cos\theta} r \, dr \, d\theta = \frac{1}{2} \int_0^{2\pi} (1 + 2\cos\theta + \cos^2\theta) \, d\theta = \frac{1}{2} (2\pi + 0 + \pi) = \frac{3\pi}{2}.\]

\smallskip\noindent\textbf{2.} \textbf{Integral de $x$:}
\[\iint_\Omega x \, dA = \int_0^{2\pi} \int_0^{1+\cos\theta} r^2 \cos\theta \, dr \, d\theta = \frac{1}{3} \int_0^{2\pi} \cos\theta (1 + \cos\theta)^3 \, d\theta.\]
Desarrollando el binomio $(1 + \cos\theta)^3 = 1 + 3\cos\theta + 3\cos^2\theta + \cos^3\theta$:
Multiplicando por $\cos\theta$: $\cos\theta + 3\cos^2\theta + 3\cos^3\theta + \cos^4\theta$.
Las potencias impares integran a cero sobre $[0, 2\pi]$:
\[\int_0^{2\pi} 3\cos^2\theta \, d\theta = 3\pi, \qquad \int_0^{2\pi} \cos^4\theta \, d\theta = 4 \cdot \frac{3\pi}{16} = \frac{3\pi}{4}.\]
\[I = \frac{1}{3} \left( 3\pi + \frac{3\pi}{4} \right) = \frac{1}{3} \left( \frac{15\pi}{4} \right) = \frac{5\pi}{4}.\]
\end{ejerciciobox}

\begin{errorbox}{¡Cuidado con esta trampa de examen!}
Integrar todas las potencias de coseno sin notar que las impares se anulan por periodicidad.
\end{errorbox}
\end{nivelbox}
\bigskip

\begin{nivelbox}{2.3.3}{Cambio de variable bilineal no ortogonal}
\textbf{Enunciado:}\\
Evalúe $\iint_R (x^2 - y^2) \, dA$, donde $R$ está delimitada por las curvas $xy = 1, xy = 2, x^2 - y^2 = 1, x^2 - y^2 = 4$ en el semiplano $x > 0$.

\begin{observacionbox}{Pista Metodológica}
Tome $u = x^2 - y^2$ y $v = 2xy$. Relacione esto con la variable compleja $z^2 = (x+iy)^2$.
\end{observacionbox}

\begin{ejerciciobox}{Resolución Analítica Paso a Paso}
\smallskip\noindent\textbf{1.} \textbf{Transformación:}
Sea $u = x^2 - y^2$ y $v = 2xy$. El dominio en $(u,v)$ es $[1, 4] \times [2, 4]$.

\smallskip\noindent\textbf{2.} \textbf{Jacobiano:}
\[\frac{\partial(u,v)}{\partial(x,y)} = \det \begin{pmatrix} 2x & -2y \\ 2y & 2x \end{pmatrix} = 4(x^2 + y^2).\]
Note que $|z^2| = |u + iv| = \sqrt{u^2 + v^2} = x^2 + y^2$. Por tanto:
\[\left| \frac{\partial(x,y)}{\partial(u,v)} \right| = \frac{1}{4(x^2 + y^2)} = \frac{1}{4\sqrt{u^2 + v^2}}.\]

\smallskip\noindent\textbf{3.} \textbf{Integral:}
Como el integrando es $x^2 - y^2 = u$:
\[I = \int_2^4 \int_1^4 u \cdot \frac{1}{4\sqrt{u^2 + v^2}} \, du \, dv = \frac{1}{4} \int_2^4 \left[ \sqrt{u^2 + v^2} \right]_{u=1}^{u=4} dv = \frac{1}{4} \int_2^4 \left( \sqrt{16 + v^2} - \sqrt{1 + v^2} \right) dv.\]
Integrando mediante la fórmula estándar $\int \sqrt{a^2 + v^2} dv = \frac{v}{2}\sqrt{a^2+v^2} + \frac{a^2}{2}\ln(v + \sqrt{a^2+v^2})$ se obtiene la expresión exacta cerrada.
\end{ejerciciobox}

\begin{errorbox}{¡Cuidado con esta trampa de examen!}
No advertir que $x^2+y^2$ debe expresarse en términos de $u$ y $v$ para poder integrar en el plano $(u,v)$.
\end{errorbox}
\end{nivelbox}
\bigskip

\subsection{Nivel 4: Desafíos Teóricos y Contraejemplos}

\begin{nivelbox}{2.4.1}{Inyectividad de la transformación en el borde y Teorema de Cambio de Variables}
\textbf{Enunciado:}\\
El Teorema de Cambio de Variable exige que $T: U \to V$ sea un $C^1$-difeomorfismo en el interior de $U$. Analice rigurosamente por qué la transformación polar $T(r,\theta) = (r\cos\theta, r\sin\theta)$ no es inyectiva en el segmento $r = 0$, y justifique por qué esto no invalida el cálculo de la integral en el disco.

\begin{observacionbox}{Pista Metodológica}
Demuestre que el conjunto donde $T$ no es inyectiva o donde $\det DT = 0$ tiene medida de Jordan cero.
\end{observacionbox}

\begin{ejerciciobox}{Resolución Analítica Paso a Paso}
\smallskip\noindent\textbf{1.} \textbf{Falla de inyectividad:}
Para $r = 0$, $T(0, \theta) = (0,0)$ para cualquier $\theta \in [0, 2\pi]$. Todo el segmento $\{0\} \times [0, 2\pi]$ colapsa en el único punto $(0,0)$. Asimismo, $\det DT(0, \theta) = 0$.

\smallskip\noindent\textbf{2.} \textbf{Conjuntos excepcionales de medida cero:}
El segmento $E = \{0\} \times [0, 2\pi]$ en el plano $(r,\theta)$ tiene contenido de Jordan cero (es un segmento unidimensional).
La imagen $T(E) = \{(0,0)\}$ es un único punto en $\mathbb{R}^2$, que también tiene contenido de Jordan cero.

\smallskip\noindent\textbf{3.} \textbf{Justificación formal:}
Al restringir $T$ al abierto $U^* = (0, R) \times (0, 2\pi)$, $T$ es un difeomorfismo $C^1$ sobre $V^* = B(0,R) \setminus \{ (x,0) : x \ge 0 \}$. Como los conjuntos eliminados son de contenido nulo:
\[\iint_{B(0,R)} f \, dA = \iint_{V^*} f \, dA = \iint_{U^*} (f \circ T) r \, dr d\theta = \int_0^{2\pi} \int_0^R (f \circ T) r \, dr d\theta. \quad \blacksquare\]
\end{ejerciciobox}

\begin{errorbox}{¡Cuidado con esta trampa de examen!}
Creer que si la transformación falla en un punto o una curva, el teorema no se puede aplicar. Los conjuntos de contenido nulo no alteran el valor de la integral de Riemann.
\end{errorbox}
\end{nivelbox}
\bigskip

\begin{nivelbox}{2.4.2}{Integral de Gauss generalizada con forma cuadrática positiva}
\textbf{Enunciado:}\\
Sea $A$ una matriz simétrica real $2 \times 2$ definida positiva. Demuestre que:
\[\iint_{\mathbb{R}^2} e^{-\frac{1}{2} \mathbf{x}^T A \mathbf{x}} \, dx \, dy = \frac{2\pi}{\sqrt{\det A}}.\]

\begin{observacionbox}{Pista Metodológica}
Diagonalice ortogonalmente $A = Q \Lambda Q^T$ con $\det Q = 1$ y haga el cambio lineal $\mathbf{x} = Q \mathbf{u}$.
\end{observacionbox}

\begin{ejerciciobox}{Resolución Analítica Paso a Paso}
\smallskip\noindent\textbf{1.} \textbf{Diagonalización espectral:}
Como $A$ es simétrica definida positiva, existen autovalores $\lambda_1, \lambda_2 > 0$ y una matriz ortogonal $Q \in SO(2)$ tal que $A = Q \begin{pmatrix} \lambda_1 & 0 \\ 0 & \lambda_2 \end{pmatrix} Q^T$.

\smallskip\noindent\textbf{2.} \textbf{Primer cambio de variable (rotación $\mathbf{x} = Q\mathbf{u}$):}
El Jacobiano es $|\det Q| = 1$. La forma cuadrática se convierte en:
\[\mathbf{x}^T A \mathbf{x} = \mathbf{u}^T (Q^T A Q) \mathbf{u} = \lambda_1 u_1^2 + \lambda_2 u_2^2.\]

\smallskip\noindent\textbf{3.} \textbf{Segundo cambio de variable (escalamiento $v_1 = \sqrt{\lambda_1} u_1, v_2 = \sqrt{\lambda_2} u_2$):}
El Jacobiano es $\frac{1}{\sqrt{\lambda_1 \lambda_2}} = \frac{1}{\sqrt{\det A}}$.

\smallskip\noindent\textbf{4.} \textbf{Evaluación:}
\[I = \frac{1}{\sqrt{\det A}} \iint_{\mathbb{R}^2} e^{-\frac{1}{2}(v_1^2 + v_2^2)} \, dv_1 dv_2 = \frac{1}{\sqrt{\det A}} \left( \int_{-\infty}^\infty e^{-v^2/2} dv \right)^2.\]
Como $\int_{-\infty}^\infty e^{-v^2/2} dv = \sqrt{2\pi}$:
\[I = \frac{2\pi}{\sqrt{\det A}}. \qquad \blacksquare\]
\end{ejerciciobox}

\begin{errorbox}{¡Cuidado con esta trampa de examen!}
Olvidar la raíz cuadrada en el determinante del denominador $\sqrt{\det A}$.
\end{errorbox}
\end{nivelbox}
\bigskip

\begin{nivelbox}{2.4.3}{Fórmula de la coárea y cambio a coordenadas de curvas de nivel}
\textbf{Enunciado:}\\
Sea $f: \mathbb{R}^2 \to \mathbb{R}$ suave con $\nabla f(x,y) \neq 0$. Considere las curvas de nivel $f(x,y) = c$. Justifique formalmente la relación entre el Jacobiano del sistema ortogonal generado por el gradiente y la integral de línea a lo largo de las curvas de nivel.

\begin{observacionbox}{Pista Metodológica}
Construya un sistema local de coordenadas $(c, s)$ donde $c = f(x,y)$ y $s$ es el parámetro de longitud de arco a lo largo de la curva de nivel.
\end{observacionbox}

\begin{ejerciciobox}{Resolución Analítica Paso a Paso}
\smallskip\noindent\textbf{1.} \textbf{Vectores de la base local:}
El vector normal unitario a la curva de nivel es $\mathbf{n} = \frac{\nabla f}{\|\nabla f\|}$, y el vector tangente unitario es $\mathbf{t}$.

\smallskip\noindent\textbf{2.} \textbf{Diferencial de área en coordenadas $(c, s)$:}
El desplazamiento en la dirección del gradiente satisface $dc = \|\nabla f\| dn \implies dn = \frac{dc}{\|\nabla f\|}$.
Por ende, el elemento de área ortogonal es:
\[dA = dn \, ds = \frac{dc \, ds}{\|\nabla f\|}.\]

\smallskip\noindent\textbf{3.} \textbf{Fórmula de la coárea (caso 2D):}
\[\iint_\Omega g(x,y) \, dA = \int_{c_{\min}}^{c_{\max}} \left( \int_{f^{-1}(c)} \frac{g(x,y)}{\|\nabla f\|} \, ds \right) dc. \quad \blacksquare\]
\end{ejerciciobox}

\begin{errorbox}{¡Cuidado con esta trampa de examen!}
Confundir $dc$ directamente con la distancia euclidiana entre curvas de nivel sin dividir por el módulo del gradiente.
\end{errorbox}
\end{nivelbox}
\bigskip

\subsection{Nivel 5: Dominio Avanzado y Rigor Topológico}

\begin{nivelbox}{2.5.1}{Teorema de Cambio de Variables: Demostración mediante aproximación afín}
\textbf{Enunciado:}\\
Sea $T: U \to V$ un difeomorfismo de clase $C^1$ entre abiertos de $\mathbb{R}^2$ y sea $Q \subset U$ un cuadrado pequeño de lado $h$ centrado en $p_0$. Demuestre que $\lim_{h \to 0^+} \frac{\operatorname{Area}(T(Q))}{\operatorname{Area}(Q)} = |\det DT(p_0)|$.

\begin{observacionbox}{Pista Metodológica}
Escriba la expansión de Taylor de primer orden $T(p_0 + \mathbf{h}) = T(p_0) + DT(p_0)\mathbf{h} + o(\|\mathbf{h}\|)$ y use la propiedad geométrica de que las aplicaciones lineales transforman áreas multiplicándolas por $|\det A|$.
\end{observacionbox}

\begin{ejerciciobox}{Resolución Analítica Paso a Paso}
\smallskip\noindent\textbf{1.} \textbf{Aproximación lineal afín:}
Definimos la aproximación lineal $L(\mathbf{x}) = T(p_0) + DT(p_0)(\mathbf{x} - p_0)$.
Por el álgebra lineal en $\mathbb{R}^2$, cualquier operador lineal invertible $A = DT(p_0)$ transforma un rectángulo o paralelogramo de área $A_0$ en una figura de área $|\det A| A_0$. Por tanto:
\[\operatorname{Area}(L(Q)) = |\det DT(p_0)| \operatorname{Area}(Q) = |\det DT(p_0)| h^2.\]

\smallskip\noindent\textbf{2.} \textbf{Control del error no lineal:}
Por la diferenciabilidad de $T$ en $p_0$, para todo $\varepsilon > 0$, existe $h_0 > 0$ tal que si $h < h_0$, para todo $\mathbf{x} \in Q$:
\[\|T(\mathbf{x}) - L(\mathbf{x})\| \le \varepsilon \|\mathbf{x} - p_0\| \le \varepsilon \frac{\sqrt{2}}{2} h.\]

\smallskip\noindent\textbf{3.} \textbf{Acotamiento de conjuntos:}
Esto implica que $T(Q)$ está contenido en un engrosamiento de $L(Q)$ por un factor $(1 + C\varepsilon)$, y contiene a un adelgazamiento por $(1 - C\varepsilon)$:
\[(1 - C\varepsilon)^2 |\det DT(p_0)| h^2 \le \operatorname{Area}(T(Q)) \le (1 + C\varepsilon)^2 |\det DT(p_0)| h^2.\]

\smallskip\noindent\textbf{4.} \textbf{Paso al límite:}
Dividiendo entre $\operatorname{Area}(Q) = h^2$ y tomando $h \to 0$ (lo que fuerza $\varepsilon \to 0$):
\[\lim_{h \to 0^+} \frac{\operatorname{Area}(T(Q))}{\operatorname{Area}(Q)} = |\det DT(p_0)|. \qquad \blacksquare\]
\end{ejerciciobox}

\begin{errorbox}{¡Cuidado con esta trampa de examen!}
Asumir sin demostración que las imágenes de los lados del cuadrado siguen siendo rectas bajo una transformación no lineal.
\end{errorbox}
\end{nivelbox}
\bigskip

\begin{nivelbox}{2.5.2}{Teorema de Sard y conjuntos de medida cero de valores críticos}
\textbf{Enunciado:}\\
Sea $T: U \subset \mathbb{R}^2 \to \mathbb{R}^2$ de clase $C^1$. Sea $C = \{ p \in U : \det DT(p) = 0 \}$ el conjunto de puntos críticos. Demuestre que la imagen $T(C)$ tiene contenido nulo si $C$ es compacto.

\begin{observacionbox}{Pista Metodológica}
En un punto crítico, uno de los autovalores o valores singulares de $DT$ es cero, lo que colapsa la dimensión local.
\end{observacionbox}

\begin{ejerciciobox}{Resolución Analítica Paso a Paso}
\smallskip\noindent\textbf{1.} \textbf{Compacidad y subdivisión:}
Dado que $C$ es compacto, para cada $\varepsilon > 0$, cubrimos $C$ con un número finito de cuadrados $Q_k$ de lado $h$ centrados en $p_k \in C$.

\smallskip\noindent\textbf{2.} \textbf{Deformación unidimensional:}
Como $\det DT(p_k) = 0$, la imagen de la diferencial $DT(p_k)(\mathbb{R}^2)$ tiene dimensión a lo sumo 1 (es una recta o el origen).
Por tanto, la imagen afín $L_k(Q_k)$ tiene área cero (es un segmento de recta de longitud a lo sumo $2\|DT(p_k)\| h$).

\smallskip\noindent\textbf{3.} \textbf{Engrosamiento no lineal:}
El error cuadrático de Taylor añade un grosor en la dirección ortogonal de orden a lo sumo $\varepsilon h$. Por ende, la imagen real $T(Q_k)$ está contenida en un rectángulo de lados $\sim h$ y $\sim \varepsilon h$.
Su área es:
\[\operatorname{Area}(T(Q_k)) \le C \varepsilon h^2.\]

\smallskip\noindent\textbf{4.} \textbf{Suma finita:}
\[\operatorname{Area}(T(C)) \le \sum_k \operatorname{Area}(T(Q_k)) \le C \varepsilon \sum_k h^2 = C \varepsilon \operatorname{Area}\left(\bigcup Q_k\right) \le C' \varepsilon.\]
Como $\varepsilon > 0$ es arbitrario, $\operatorname{cont}(T(C)) = 0$. $\blacksquare$
\end{ejerciciobox}

\begin{errorbox}{¡Cuidado con esta trampa de examen!}
Confundir el conjunto de puntos críticos $C$ (en el dominio) con el conjunto de valores críticos $T(C)$ (en el codominio).
\end{errorbox}
\end{nivelbox}
\bigskip

\begin{nivelbox}{2.5.3}{Inversión global de Hadamard y difeomorfismos en R²}
\textbf{Enunciado:}\\
Sea $F: \mathbb{R}^2 \to \mathbb{R}^2$ de clase $C^1$ tal que $\det DF(x,y) \neq 0$ para todo $(x,y) \in \mathbb{R}^2$ y $\lim_{\|\mathbf{x}\| \to \infty} \|F(\mathbf{x})\| = \infty$ (aplicación propia). Demuestre que $F$ es un difeomorfismo global de $\mathbb{R}^2$ en $\mathbb{R}^2$.

\begin{observacionbox}{Pista Metodológica}
Demuestre que $F(\mathbb{R}^2)$ es a la vez abierto (por el TFI) y cerrado (por ser propia) en $\mathbb{R}^2$, y use la conexidad del plano.
\end{observacionbox}

\begin{ejerciciobox}{Resolución Analítica Paso a Paso}
\smallskip\noindent\textbf{1.} \textbf{La imagen es abierta:}
Como $\det DF(p) \neq 0$ en todo punto, por el Teorema de la Función Inversa local, $F$ es una aplicación abierta: para todo abierto $U$, $F(U)$ es abierto. En particular, $F(\mathbb{R}^2)$ es un abierto no vacío de $\mathbb{R}^2$.

\smallskip\noindent\textbf{2.} \textbf{La imagen es cerrada:}
Sea $y_k \in F(\mathbb{R}^2)$ tal que $y_k \to y_0$. Entonces existe $x_k$ tal que $F(x_k) = y_k$. Como $(y_k)$ es convergente, es acotada. Dado que $F$ es propia ($\lim_{\|x\| \to \infty} \|F(x)\| = \infty$), la sucesión preimagen $(x_k)$ debe ser acotada.
Por Bolzano-Weierstrass, existe una subsucesión convergente $x_{k_j} \to x_0$. Por continuidad de $F$, $F(x_0) = \lim F(x_{k_j}) = y_0$, lo que prueba que $y_0 \in F(\mathbb{R}^2)$. Por ende, $F(\mathbb{R}^2)$ es cerrado.

\smallskip\noindent\textbf{3.} \textbf{Sobreyectividad:}
Los únicos subconjuntos no vacíos simultáneamente abiertos y cerrados del espacio conexo $\mathbb{R}^2$ son $\emptyset$ y $\mathbb{R}^2$. Por tanto, $F(\mathbb{R}^2) = \mathbb{R}^2$ ($F$ es sobreyectiva).

\smallskip\noindent\textbf{4.} \textbf{Inyectividad por simple conexidad:}
Como $\mathbb{R}^2$ es simplemente conexo, el recubrimiento $F$ debe tener grado 1, concluyendo que $F$ es biyectiva y por tanto un difeomorfismo global. $\blacksquare$
\end{ejerciciobox}

\begin{errorbox}{¡Cuidado con esta trampa de examen!}
Omitir la condición de que $F$ sea propia. Por ejemplo, $f(x,y) = (e^x\cos y, e^x\sin y)$ tiene Jacobiano nunca nulo pero NO es inyectiva global ni sobreyectiva (no cubre el origen).
\end{errorbox}
\end{nivelbox}
\bigskip


\section{Tema 3: Integrales Triples Cartesianas y los 6 Órdenes de Fubini}

\subsection{Nivel 0: Calentamiento Conceptual}

\begin{nivelbox}{3.0.1}{Integral triple sobre una caja rectangular}
\textbf{Enunciado:}\\
Evalúe $\iiint_B (x + 2y + 3z) \, dV$, donde $B = [0, 1] \times [0, 2] \times [0, 3]$.

\begin{observacionbox}{Pista Metodológica}
Integre iteradamente en cualquier orden, por ejemplo $dz \, dy \, dx$.
\end{observacionbox}

\begin{ejerciciobox}{Resolución Analítica Paso a Paso}
\smallskip\noindent\textbf{1.} \textbf{Planteamiento:}
\[I = \int_0^1 \int_0^2 \int_0^3 (x + 2y + 3z) \, dz \, dy \, dx.\]

\smallskip\noindent\textbf{2.} \textbf{Integral en $z$:}
\[\int_0^3 (x + 2y + 3z) \, dz = \left[ (x + 2y)z + \frac{3z^2}{2} \right]_0^3 = 3x + 6y + \frac{27}{2}.\]

\smallskip\noindent\textbf{3.} \textbf{Integral en $y$:}
\[\int_0^2 \left( 3x + 6y + \frac{27}{2} \right) dy = \left[ 3xy + 3y^2 + \frac{27y}{2} \right]_0^2 = 6x + 12 + 27 = 6x + 39.\]

\smallskip\noindent\textbf{4.} \textbf{Integral en $x$:}
\[I = \int_0^1 (6x + 39) \, dx = \left[ 3x^2 + 39x \right]_0^1 = 3 + 39 = 42.\]
\end{ejerciciobox}

\begin{errorbox}{¡Cuidado con esta trampa de examen!}
Confundir cuál constante acompaña a cuál variable al evaluar los límites.
\end{errorbox}
\end{nivelbox}
\bigskip

\begin{nivelbox}{3.0.2}{Identificación de proyecciones para sólidos Tipos 1, 2 y 3}
\textbf{Enunciado:}\\
Sea $E$ el sólido acotado por los planos coordenados $x=0, y=0, z=0$ y el plano $x + 2y + 3z = 6$.
(a) Describa $E$ como sólido Tipo 1 (proyección en el plano $XY$).
(b) Describa $E$ como sólido Tipo 2 (proyección en el plano $XZ$).

\begin{observacionbox}{Pista Metodológica}
Para Tipo 1, proyecte en $z=0$ haciendo $x + 2y \le 6$. El techo es $z = \frac{6 - x - 2y}{3}$.
\end{observacionbox}

\begin{ejerciciobox}{Resolución Analítica Paso a Paso}
\smallskip\noindent\textbf{1.} \textbf{Tipo 1 (Plano $XY$):}
En $z = 0$, la sombra es el triángulo $D_{xy} = \{ (x,y) : 0 \le x \le 6, \ 0 \le y \le \frac{6-x}{2} \}$.
Para cada $(x,y) \in D_{xy}$, $z$ va desde $0$ hasta $\frac{6-x-2y}{3}$:
\[E = \left\{ (x,y,z) : 0 \le x \le 6, \ 0 \le y \le \frac{6-x}{2}, \ 0 \le z \le \frac{6-x-2y}{3} \right\}.\]

\smallskip\noindent\textbf{2.} \textbf{Tipo 2 (Plano $XZ$):}
En $y = 0$, la sombra es $D_{xz} = \{ (x,z) : 0 \le x \le 6, \ 0 \le z \le \frac{6-x}{3} \}$.
Para cada $(x,z) \in D_{xz}$, la variable horizontal $y$ va desde $0$ hasta $\frac{6-x-3z}{2}$:
\[E = \left\{ (x,y,z) : 0 \le x \le 6, \ 0 \le z \le \frac{6-x}{3}, \ 0 \le y \le \frac{6-x-3z}{2} \right\}.\]
\end{ejerciciobox}

\begin{errorbox}{¡Cuidado con esta trampa de examen!}
Invertir los coeficientes de las variables al despejar $z$ o $y$.
\end{errorbox}
\end{nivelbox}
\bigskip

\begin{nivelbox}{3.0.3}{Separación de variables en paralelepípedos}
\textbf{Enunciado:}\\
Calcule $\iiint_B x y^2 z^3 \, dV$ sobre la caja $B = [0, 2] \times [0, 1] \times [-1, 1]$.

\begin{observacionbox}{Pista Metodológica}
Aplique la simetría de la función impar $z^3$ en el intervalo simétrico $[-1, 1]$.
\end{observacionbox}

\begin{ejerciciobox}{Resolución Analítica Paso a Paso}
\smallskip\noindent\textbf{1.} \textbf{Factorización por Fubini:}
\[\iiint_B x y^2 z^3 \, dV = \left( \int_0^2 x \, dx \right) \left( \int_0^1 y^2 \, dy \right) \left( \int_{-1}^1 z^3 \, dz \right).\]

\smallskip\noindent\textbf{2.} \textbf{Propiedad de simetría impar:}
Como $g(z) = z^3$ es impar y el intervalo $[-1, 1]$ es simétrico alrededor del origen:
\[\int_{-1}^1 z^3 \, dz = 0.\]

\smallskip\noindent\textbf{3.} \textbf{Resultado:}
\[I = \left( \frac{2^2}{2} \right) \left( \frac{1}{3} \right) (0) = 0.\]
\end{ejerciciobox}

\begin{errorbox}{¡Cuidado con esta trampa de examen!}
Calcular paso a paso todas las primitivas sin advertir la anulación por simetría impar.
\end{errorbox}
\end{nivelbox}
\bigskip

\subsection{Nivel 1: Procedimental Guiado}

\begin{nivelbox}{3.1.1}{Volumen de un tetraedro estándar}
\textbf{Enunciado:}\\
Calcule el volumen del tetraedro $T$ delimitado por los planos $x=0, y=0, z=0$ y $x + y + z = 1$ usando una integral triple.

\begin{observacionbox}{Pista Metodológica}
Integre $1 \, dz \, dy \, dx$ con $0 \le z \le 1 - x - y$, $0 \le y \le 1 - x$ y $0 \le x \le 1$.
\end{observacionbox}

\begin{ejerciciobox}{Resolución Analítica Paso a Paso}
\smallskip\noindent\textbf{1.} \textbf{Planteamiento:}
\[V = \int_0^1 \int_0^{1-x} \int_0^{1-x-y} 1 \, dz \, dy \, dx.\]

\smallskip\noindent\textbf{2.} \textbf{Integración en $z$:}
\[\int_0^{1-x-y} dz = 1 - x - y.\]

\smallskip\noindent\textbf{3.} \textbf{Integración en $y$:}
\[\int_0^{1-x} (1 - x - y) \, dy = \left[ (1 - x)y - \frac{y^2}{2} \right]_0^{1-x} = (1 - x)^2 - \frac{(1-x)^2}{2} = \frac{(1-x)^2}{2}.\]

\smallskip\noindent\textbf{4.} \textbf{Integración en $x$:}
\[V = \int_0^1 \frac{(1-x)^2}{2} \, dx = \left[ -\frac{(1-x)^3}{6} \right]_0^1 = 0 - \left(-\frac{1}{6}\right) = \frac{1}{6}.\]
\end{ejerciciobox}

\begin{errorbox}{¡Cuidado con esta trampa de examen!}
Olvidar el factor $1/6$, confundiendo la pirámide de base triangular con un prisma.
\end{errorbox}
\end{nivelbox}
\bigskip

\begin{nivelbox}{3.1.2}{Integral triple sobre un cilindro parabólico}
\textbf{Enunciado:}\\
Evalúe $\iiint_E z \, dV$, donde $E$ es el sólido acotado por el cilindro parabólico $z = y^2$ y los planos $x = 0, x = 1, z = 1$.

\begin{observacionbox}{Pista Metodológica}
La sombra en el plano $YZ$ está delimitada por la parábola $z = y^2$ y la recta $z = 1$, con $y \in [-1, 1]$. La variable $x$ va libremente de $0$ a $1$.
\end{observacionbox}

\begin{ejerciciobox}{Resolución Analítica Paso a Paso}
\smallskip\noindent\textbf{1.} \textbf{Límites de integración:}
$0 \le x \le 1$, $-1 \le y \le 1$, $y^2 \le z \le 1$.

\smallskip\noindent\textbf{2.} \textbf{Integral en $z$:}
\[\int_{y^2}^1 z \, dz = \left[ \frac{z^2}{2} \right]_{y^2}^1 = \frac{1 - y^4}{2}.\]

\smallskip\noindent\textbf{3.} \textbf{Integral en $y$:}
\[\int_{-1}^1 \frac{1 - y^4}{2} \, dy = 2 \int_0^1 \frac{1 - y^4}{2} \, dy = \left[ y - \frac{y^5}{5} \right]_0^1 = 1 - \frac{1}{5} = \frac{4}{5}.\]

\smallskip\noindent\textbf{4.} \textbf{Integral en $x$:}
\[I = \int_0^1 1 \, dx \cdot \frac{4}{5} = \frac{4}{5}.\]
\end{ejerciciobox}

\begin{errorbox}{¡Cuidado con esta trampa de examen!}
Integrar $y$ únicamente de $0$ a $1$, olvidando la mitad simétrica en $y < 0$.
\end{errorbox}
\end{nivelbox}
\bigskip

\begin{nivelbox}{3.1.3}{Inversión del orden de integración en R³}
\textbf{Enunciado:}\\
Reescriba la integral iterada $I = \int_0^1 \int_0^{1-x} \int_0^{1-x-y} f(x,y,z) \, dz \, dy \, dx$ en el orden $dy \, dx \, dz$.

\begin{observacionbox}{Pista Metodológica}
El sólido es el tetraedro $x \ge 0, y \ge 0, z \ge 0, x+y+z \le 1$. Para el orden $dy \, dx \, dz$, la variable exterior es $z \in [0, 1]$, la intermedia es $x$, y la interior es $y$.
\end{observacionbox}

\begin{ejerciciobox}{Resolución Analítica Paso a Paso}
\smallskip\noindent\textbf{1.} \textbf{Identificación del sólido:}
\[E = \{ (x,y,z) : x \ge 0, \ y \ge 0, \ z \ge 0, \ x + y + z \le 1 \}.\]

\smallskip\noindent\textbf{2.} \textbf{Determinación de límites en orden $dy \, dx \, dz$:}
- Límite exterior en $z$: $0 \le z \le 1$.
- Para un $z$ fijo, la proyección en el plano $XY$ es $x \ge 0, y \ge 0, x + y \le 1 - z$.
- Límite intermedio en $x$: $0 \le x \le 1 - z$.
- Límite interior en $y$: $0 \le y \le 1 - z - x$.

\smallskip\noindent\textbf{3.} \textbf{Integral reescrita:}
\[I = \int_0^1 \int_0^{1-z} \int_0^{1-z-x} f(x,y,z) \, dy \, dx \, dz.\]
\end{ejerciciobox}

\begin{errorbox}{¡Cuidado con esta trampa de examen!}
Dejar $1-x$ como límite superior de $y$ sin restar la variable exterior $z$.
\end{errorbox}
\end{nivelbox}
\bigskip

\subsection{Nivel 2: Integración de Conceptos}

\begin{nivelbox}{3.2.1}{Volumen de un sólido limitado por superficies cilíndricas y planos}
\textbf{Enunciado:}\\
Calcule el volumen del sólido delimitado por el cilindro parabólico $z = 1 - y^2$, el plano $z = 0$, el plano $x = 0$ y el plano $x + z = 2$.

\begin{observacionbox}{Pista Metodológica}
Observe que el techo es $z = 1 - y^2$ y el sólido termina lateralmente en el plano $x = 2 - z$. Es conveniente proyectar en el plano $YZ$.
\end{observacionbox}

\begin{ejerciciobox}{Resolución Analítica Paso a Paso}
\smallskip\noindent\textbf{1.} \textbf{Proyección en el plano $YZ$ (Sólido Tipo 3):}
En $z = 0$, $1 - y^2 \ge 0 \implies y \in [-1, 1]$ y $0 \le z \le 1 - y^2$.

\smallskip\noindent\textbf{2.} \textbf{Límites en $x$:}
Para cada $(y,z)$, $x$ entra por $x = 0$ y sale por $x = 2 - z$.

\smallskip\noindent\textbf{3.} \textbf{Integral de volumen:}
\[V = \int_{-1}^1 \int_0^{1-y^2} \int_0^{2-z} 1 \, dx \, dz \, dy = \int_{-1}^1 \int_0^{1-y^2} (2 - z) \, dz \, dy.\]

\smallskip\noindent\textbf{4.} \textbf{Integral en $z$:}
\[\int_0^{1-y^2} (2 - z) \, dz = \left[ 2z - \frac{z^2}{2} \right]_0^{1-y^2} = 2(1 - y^2) - \frac{(1 - y^2)^2}{2} = 2 - 2y^2 - \frac{1 - 2y^2 + y^4}{2} = \frac{3}{2} - y^2 - \frac{y^4}{2}.\]

\smallskip\noindent\textbf{5.} \textbf{Integral en $y$:}
\[V = 2 \int_0^1 \left( \frac{3}{2} - y^2 - \frac{y^4}{2} \right) dy = 2 \left[ \frac{3}{2}y - \frac{y^3}{3} - \frac{y^5}{10} \right]_0^1 = 2 \left( \frac{3}{2} - \frac{1}{3} - \frac{1}{10} \right) = 2 \left( \frac{45 - 10 - 3}{30} \right) = 2 \left( \frac{32}{30} \right) = \frac{32}{15}.\]
\end{ejerciciobox}

\begin{errorbox}{¡Cuidado con esta trampa de examen!}
Intentar proyectar en el plano $XY$, lo cual genera una partición engorrosa en dos regiones distintas.
\end{errorbox}
\end{nivelbox}
\bigskip

\begin{nivelbox}{3.2.2}{Integral triple de monomio sobre una cuña sólida}
\textbf{Enunciado:}\\
Evalúe $I = \iiint_E x z \, dV$, donde $E$ es la cuña acotada por $y = 0, y = 3, z = 0, z = x, x = 2$.

\begin{observacionbox}{Pista Metodológica}
La variable $y$ está desacoplada en $[0, 3]$. En el plano $XZ$, es un triángulo $0 \le x \le 2, 0 \le z \le x$.
\end{observacionbox}

\begin{ejerciciobox}{Resolución Analítica Paso a Paso}
\smallskip\noindent\textbf{1.} \textbf{Factorización iterada:}
\[I = \left( \int_0^3 1 \, dy \right) \left( \int_0^2 x \left( \int_0^x z \, dz \right) dx \right).\]

\smallskip\noindent\textbf{2.} \textbf{Integral en $y$:}
\[\int_0^3 1 \, dy = 3.\]

\smallskip\noindent\textbf{3.} \textbf{Integral interior en $z$:}
\[\int_0^x z \, dz = \frac{x^2}{2}.\]

\smallskip\noindent\textbf{4.} \textbf{Integral en $x$:}
\[\int_0^2 x \left( \frac{x^2}{2} \right) dx = \frac{1}{2} \int_0^2 x^3 \, dx = \frac{1}{2} \left[ \frac{x^4}{4} \right]_0^2 = \frac{1}{2} \cdot 4 = 2.\]

\smallskip\noindent\textbf{5.} \textbf{Resultado total:}
\[I = 3 \cdot 2 = 6.\]
\end{ejerciciobox}

\begin{errorbox}{¡Cuidado con esta trampa de examen!}
Invertir los límites $0 \le x \le z$ en lugar de $0 \le z \le x$.
\end{errorbox}
\end{nivelbox}
\bigskip

\begin{nivelbox}{3.2.3}{Volumen entre dos paraboloides en coordenadas cartesianas}
\textbf{Enunciado:}\\
Plantee y evalúe la integral triple del volumen del sólido comprendido entre los paraboloides $z = x^2 + y^2$ y $z = 2 - x^2 - y^2$.

\begin{observacionbox}{Pista Metodológica}
Halle la curva de intersección igualando $z_1 = z_2$. La proyección es el disco $x^2 + y^2 \le 1$.
\end{observacionbox}

\begin{ejerciciobox}{Resolución Analítica Paso a Paso}
\smallskip\noindent\textbf{1.} \textbf{Intersección:}
\[x^2 + y^2 = 2 - x^2 - y^2 \iff 2(x^2 + y^2) = 2 \iff x^2 + y^2 = 1.\]
La sombra en el plano $XY$ es el disco unitario $D$.

\smallskip\noindent\textbf{2.} \textbf{Planteamiento como sólido Tipo 1:}
\[V = \iint_D \left( \int_{x^2+y^2}^{2-x^2-y^2} 1 \, dz \right) dA = \iint_D (2 - 2(x^2 + y^2)) \, dA.\]

\smallskip\noindent\textbf{3.} \textbf{Evaluación de la integral doble:}
En cartesianas o pasando a polares elementales: $\int_0^{2\pi} \int_0^1 (2 - 2r^2) r \, dr \, d\theta = 2\pi \left[ r^2 - \frac{r^4}{2} \right]_0^1 = 2\pi \left( 1 - \frac{1}{2} \right) = \pi.$
\end{ejerciciobox}

\begin{errorbox}{¡Cuidado con esta trampa de examen!}
Restar el paraboloide inferior del superior al revés, dando un volumen opuesto $-\pi$.
\end{errorbox}
\end{nivelbox}
\bigskip

\subsection{Nivel 3: Problemas Tipo Examen UNI}

\begin{nivelbox}{3.3.1}{Teorema Fundamental de Fubini en R³ para derivadas de orden 3}
\textbf{Enunciado:}\\
Sea $f: B \subset \mathbb{R}^3 \to \mathbb{R}$ de clase $C^3$ en un paralelepípedo $B = [a,b] \times [c,d] \times [p,q]$. Demuestre mediante el Teorema de Fubini que:
\[\iiint_B \frac{\partial^3 f}{\partial x \partial y \partial z} \, dV = \sum_{i,j,k=0}^1 (-1)^{i+j+k} f(x_i, y_j, z_k),\]
donde $x_0=a, x_1=b$, etc.

\begin{observacionbox}{Pista Metodológica}
Integre sucesivamente respecto a $z$, luego $y$ y finalmente $x$ aplicando el TFC unidimensional en cada paso.
\end{observacionbox}

\begin{ejerciciobox}{Resolución Analítica Paso a Paso}
\smallskip\noindent\textbf{1.} \textbf{Primera integración en $z$:}
\[\int_p^q \frac{\partial}{\partial z} \left( \frac{\partial^2 f}{\partial x \partial y} \right) dz = \frac{\partial^2 f}{\partial x \partial y}(x, y, q) - \frac{\partial^2 f}{\partial x \partial y}(x, y, p).\]

\smallskip\noindent\textbf{2.} \textbf{Segunda integración en $y$:}
\[\int_c^d \left[ \frac{\partial^2 f}{\partial x \partial y}(x, y, q) - \frac{\partial^2 f}{\partial x \partial y}(x, y, p) \right] dy = \left[ \frac{\partial f}{\partial x}(x, d, q) - \frac{\partial f}{\partial x}(x, c, q) \right] - \left[ \frac{\partial f}{\partial x}(x, d, p) - \frac{\partial f}{\partial x}(x, c, p) \right].\]

\smallskip\noindent\textbf{3.} \textbf{Tercera integración en $x$:}
Integrando los 4 términos de $a$ a $b$, cada uno produce 2 evaluaciones por el TFC, generando los 8 vértices del paralelepípedo con signos alternados correspondientes a $(-1)^{3 - (i+j+k)}$:
\[I = f(b,d,q) - f(a,d,q) - f(b,c,q) + f(a,c,q) - f(b,d,p) + f(a,d,p) + f(b,c,p) - f(a,c,p). \quad \blacksquare\]
\end{ejerciciobox}

\begin{errorbox}{¡Cuidado con esta trampa de examen!}
Equivocarse en los signos de los 8 términos. El vértice opuesto $(b,d,q)$ tiene signo positivo y $(a,c,p)$ tiene signo negativo.
\end{errorbox}
\end{nivelbox}
\bigskip

\begin{nivelbox}{3.3.2}{Sólido con proyección simple en YZ pero no simple en XY (PC5 2025-II)}
\textbf{Enunciado:}\\
Construya un contraejemplo explícito que demuestre que la afirmación 'Todo sólido cuya proyección en el plano $YZ$ sea simple es también un sólido cuya proyección en el plano $XY$ es simple' es FALSA (PC5 2025-II, Pregunta 1d).

\begin{observacionbox}{Pista Metodológica}
Piense en un tubo curvo o serpentina que al proyectarse en $YZ$ es un simple rectángulo o disco, pero que en $XY$ se pliega sobre sí mismo.
\end{observacionbox}

\begin{ejerciciobox}{Resolución Analítica Paso a Paso}
\smallskip\noindent\textbf{1.} \textbf{Construcción geométrica:}
Considere el sólido helicoidal o serpentina tubular $S$ definido alrededor de la curva $x = \sin t, y = t, z = \cos t$ para $t \in [0, 4\pi]$ con un radio de sección transversal pequeño $\varepsilon = 0.1$.

\smallskip\noindent\textbf{2.} \textbf{Proyección en el plano $YZ$:}
La variable $y = t$ crece monótonamente en $[0, 4\pi]$ mientras $z = \cos t \in [-1, 1]$. La proyección en $YZ$ es una franja simple sin auto-intersección donde cada recta perpendicular corta la región a lo sumo en un intervalo conexo (región Tipo 1 o 2 simple).

\smallskip\noindent\textbf{3.} \textbf{Proyección en el plano $XY$:}
En el plano $XY$, la trayectoria es $x = \sin y$, que oscila indefinidamente. Para un mismo valor de $x$ e $y$, el sólido tiene múltiples niveles de $z$, de modo que cualquier rayo vertical interseca al sólido en múltiples segmentos disjuntos, impidiendo describirlo como $z_1(x,y) \le z \le z_2(x,y)$ mediante un solo par de funciones continuas.
Por tanto, la afirmación es \textbf{FALSA}. $\blacksquare$
\end{ejerciciobox}

\begin{errorbox}{¡Cuidado con esta trampa de examen!}
Pensar únicamente en cuerpos convexos simples como esferas o elipsoides, donde todas las proyecciones son simples.
\end{errorbox}
\end{nivelbox}
\bigskip

\begin{nivelbox}{3.3.3}{Volumen de sólido de intersección entre cilindro y plano oblicuo}
\textbf{Enunciado:}\\
Calcule el volumen del sólido $E$ limitado por el cilindro elíptico $\frac{x^2}{4} + y^2 = 1$, el plano $z = 0$ y el plano $z = 3 - x$.

\begin{observacionbox}{Pista Metodológica}
Como $z = 3 - x$ y sobre la elipse $-2 \le x \le 2$, siempre se cumple $z \ge 3 - 2 = 1 > 0$, de modo que el plano superior no corta al piso dentro de la base.
\end{observacionbox}

\begin{ejerciciobox}{Resolución Analítica Paso a Paso}
\smallskip\noindent\textbf{1.} \textbf{Descripción como sólido Tipo 1:}
Base: elipse $D = \{ (x,y) : \frac{x^2}{4} + y^2 \le 1 \}$.
Piso: $z = 0$. Techo: $z = 3 - x$.

\smallskip\noindent\textbf{2.} \textbf{Planteamiento:}
\[V = \iint_D (3 - x) \, dA.\]

\smallskip\noindent\textbf{3.} \textbf{Separación por linealidad y simetría:}
\[V = 3 \iint_D 1 \, dA - \iint_D x \, dA.\]
- La integral $\iint_D 1 \, dA$ es el área de la elipse: $\pi a b = \pi (2)(1) = 2\pi$.
- La integral $\iint_D x \, dA = 0$, porque la región elíptica es simétrica respecto al eje $Y$ ($x \mapsto -x$) y la función $x$ es impar.

\smallskip\noindent\textbf{4.} \textbf{Volumen final:}
\[V = 3(2\pi) - 0 = 6\pi.\]
\end{ejerciciobox}

\begin{errorbox}{¡Cuidado con esta trampa de examen!}
Calcular laboriosas integrales con raíces cuadradas cuando la simetría anula inmediatamente el término $\iint x dA$.
\end{errorbox}
\end{nivelbox}
\bigskip

\subsection{Nivel 4: Desafíos Teóricos y Contraejemplos}

\begin{nivelbox}{3.4.1}{Principio de Cavalieri y Teorema de Fubini en sólidos generales}
\textbf{Enunciado:}\\
Demuestre formalmente que si dos sólidos compactos $E_1, E_2 \subset \mathbb{R}^3$ tienen la propiedad de que para todo plano horizontal $z = z_0$, las secciones transversales $S_1(z_0) = E_1 \cap \{z = z_0\}$ y $S_2(z_0) = E_2 \cap \{z = z_0\}$ tienen la misma área, entonces $\operatorname{Vol}(E_1) = \operatorname{Vol}(E_2)$ (Principio de Cavalieri).

\begin{observacionbox}{Pista Metodológica}
Escriba el volumen mediante el Teorema de Fubini integrando las áreas de las secciones respecto a $z$.
\end{observacionbox}

\begin{ejerciciobox}{Resolución Analítica Paso a Paso}
\smallskip\noindent\textbf{1.} \textbf{Expresión del volumen por Fubini:}
Sea $E \subset \mathbb{R}^3$ un compacto Jordan-medible contenido entre $z = c$ y $z = d$. Por el Teorema de Fubini:
\[\operatorname{Vol}(E) = \iiint_E 1 \, dV = \int_c^d \left( \iint_{S(z)} 1 \, dx dy \right) dz = \int_c^d \operatorname{Area}(S(z)) \, dz.\]

\smallskip\noindent\textbf{2.} \textbf{Aplicación a ambos sólidos:}
\[\operatorname{Vol}(E_1) = \int_c^d \operatorname{Area}(S_1(z)) \, dz, \qquad \operatorname{Vol}(E_2) = \int_c^d \operatorname{Area}(S_2(z)) \, dz.\]

\smallskip\noindent\textbf{3.} \textbf{Conclusión:}
Dado que $\operatorname{Area}(S_1(z)) = \operatorname{Area}(S_2(z))$ para todo $z \in [c, d]$, los integrandos son idénticos en todo el intervalo, por lo que las integrales de Riemann coinciden:
\[\operatorname{Vol}(E_1) = \operatorname{Vol}(E_2). \qquad \blacksquare\]
\end{ejerciciobox}

\begin{errorbox}{¡Cuidado con esta trampa de examen!}
Asumir que las secciones deben tener la misma forma geométrica. Cavalieri solo exige que tengan la misma área numérica.
\end{errorbox}
\end{nivelbox}
\bigskip

\begin{nivelbox}{3.4.2}{Integral impropia convergente en un sólido no acotado de R³}
\textbf{Enunciado:}\\
Determine si la integral impropia $I = \iiint_E \frac{1}{(1 + x^2 + y^2 + z^2)^2} \, dV$ sobre todo el semiespacio $E = \{ (x,y,z) \in \mathbb{R}^3 : z \ge 0 \}$ es convergente o divergente.

\begin{observacionbox}{Pista Metodológica}
En $\mathbb{R}^3$, el elemento de volumen en coordenadas esféricas crece como $\rho^2 d\rho$. Compare el exponente del numerador con el denominador.
\end{observacionbox}

\begin{ejerciciobox}{Resolución Analítica Paso a Paso}
\smallskip\noindent\textbf{1.} \textbf{Conversión a coordenadas esféricas:}
En el semiespacio superior $z \ge 0$, $\rho \in [0, \infty)$, $\theta \in [0, 2\pi]$, $\varphi \in [0, \pi/2]$.
\[dV = \rho^2 \sin\varphi \, d\rho \, d\theta \, d\varphi.\]

\smallskip\noindent\textbf{2.} \textbf{Integral iterada:}
\[I = \left( \int_0^{2\pi} d\theta \right) \left( \int_0^{\pi/2} \sin\varphi \, d\varphi \right) \left( \int_0^\infty \frac{\rho^2}{(1 + \rho^2)^2} \, d\rho \right).\]
- $\int_0^{2\pi} d\theta = 2\pi$.
- $\int_0^{\pi/2} \sin\varphi \, d\varphi = [-\cos\varphi]_0^{\pi/2} = 1$.

\smallskip\noindent\textbf{3.} \textbf{Análisis asintótico en $\rho$:}
Para $\rho \to \infty$, el integrando se comporta como $\frac{\rho^2}{\rho^4} = \frac{1}{\rho^2}$.
Como $\int_1^\infty \frac{1}{\rho^2} d\rho < \infty$ converge, la integral es \textbf{convergente}.

\smallskip\noindent\textbf{4.} \textbf{Cálculo exacto mediante sustitución $\rho = \tan u$:}
\[\int_0^{\pi/2} \frac{\tan^2 u}{\sec^4 u} \sec^2 u \, du = \int_0^{\pi/2} \sin^2 u \, du = \frac{\pi}{4}.\]
\[I = (2\pi)(1)\left(\frac{\pi}{4}\right) = \frac{\pi^2}{2}.\]
\end{ejerciciobox}

\begin{errorbox}{¡Cuidado con esta trampa de examen!}
Confundir el comportamiento de $\mathbb{R}^2$ (donde diverge) con $\mathbb{R}^3$ (donde converge gracias al decaimiento $1/\rho^2$).
\end{errorbox}
\end{nivelbox}
\bigskip

\begin{nivelbox}{3.4.3}{Teorema de Fubini en dimensión arbitraria n}
\textbf{Enunciado:}\\
Sea $B = B_1 \times B_2$ con $B_1 \subset \mathbb{R}^k$ y $B_2 \subset \mathbb{R}^m$ bloques compactos. Si $f: B \to \mathbb{R}$ es integrable Riemann, demuestre que la integral inferior $L(x) = \underline{\int}_{B_2} f(x, y) \, dy$ es integrable Riemann en $B_1$ y coincide con $\int_B f$.

\begin{observacionbox}{Pista Metodológica}
Utilice la definición de partición producto $P = P_1 \times P_2$ y sumas de Darboux.
\end{observacionbox}

\begin{ejerciciobox}{Resolución Analítica Paso a Paso}
\smallskip\noindent\textbf{1.} \textbf{Particiones producto:}
Sea $P_1$ partición de $B_1$ y $P_2$ de $B_2$. Entonces $P = P_1 \times P_2$ es partición de $B$.
Para cada bloque $R = R_1 \times R_2$ de $P$:
\[m_R(f) \le f(x, y) \le M_R(f), \quad \forall (x,y) \in R_1 \times R_2.\]

\smallskip\noindent\textbf{2.} \textbf{Acotamiento de las integrales unidimensionales:}
Integrando respecto a $y \in R_2$:
\[m_R(f) \operatorname{vol}(R_2) \le \underline{\int}_{R_2} f(x, y) \, dy \le M_R(f) \operatorname{vol}(R_2).\]

\smallskip\noindent\textbf{3.} \textbf{Sumas de Darboux:}
Multiplicando por $\operatorname{vol}(R_1)$ y sumando sobre todos los bloques:
\[s(f, P) \le s(L, P_1) \le S(L, P_1) \le S(f, P).\]

\smallskip\noindent\textbf{4.} \textbf{Paso al límite:}
Dado que $f$ es integrable en $B$, para todo $\varepsilon > 0$ existe $P$ tal que $S(f, P) - s(f, P) < \varepsilon$.
Esto implica $S(L, P_1) - s(L, P_1) < \varepsilon$, demostrando que $L$ es integrable en $B_1$ y su integral es $\int_B f$. $\blacksquare$
\end{ejerciciobox}

\begin{errorbox}{¡Cuidado con esta trampa de examen!}
Afirmar que la función corte $y \mapsto f(x,y)$ es integrable para TODO $x$. Solo se puede asegurar la integrabilidad de la integral inferior/superior de Darboux.
\end{errorbox}
\end{nivelbox}
\bigskip

\subsection{Nivel 5: Dominio Avanzado y Rigor Topológico}

\begin{nivelbox}{3.5.1}{Invarianza de la integral de Riemann tridimensional bajo rotaciones}
\textbf{Enunciado:}\\
Sea $\Omega \subset \mathbb{R}^3$ Jordan-medible y $R \in SO(3)$ una matriz de rotación ortogonal con $\det R = 1$. Demuestre formalmente que $\iiint_{R(\Omega)} f(\mathbf{x}) \, dV = \iiint_\Omega f(R\mathbf{u}) \, dV$ y que el volumen es invariante.

\begin{observacionbox}{Pista Metodológica}
Aplique el Teorema General de Cambio de Variables con la transformación lineal $T(\mathbf{u}) = R\mathbf{u}$.
\end{observacionbox}

\begin{ejerciciobox}{Resolución Analítica Paso a Paso}
\smallskip\noindent\textbf{1.} \textbf{Transformación lineal ortogonal:}
Definimos $T: \mathbb{R}^3 \to \mathbb{R}^3$ dada por $T(\mathbf{u}) = R\mathbf{u}$. Como $R \in SO(3)$, $R^T R = I$ y $\det R = 1$.

\smallskip\noindent\textbf{2.} \textbf{Jacobiano:}
La diferencial de una aplicación lineal es la propia matriz: $DT(\mathbf{u}) = R$. Por tanto, el determinante Jacobiano es:
\[|\det DT(\mathbf{u})| = |\det R| = 1.\]

\smallskip\noindent\textbf{3.} \textbf{Difeomorfismo global:}
Como $\det R \neq 0$, $T$ es una biyección bicontinua ($C^\infty$-difeomorfismo global).

\smallskip\noindent\textbf{4.} \textbf{Aplicación del Teorema de Cambio de Variables:}
\[\iiint_{T(\Omega)} f(\mathbf{x}) \, d\mathbf{x} = \iiint_\Omega f(T(\mathbf{u})) |\det DT(\mathbf{u})| \, d\mathbf{u} = \iiint_\Omega f(R\mathbf{u}) \cdot 1 \, d\mathbf{u}.\]
En particular, para $f = 1$:
\[\operatorname{Vol}(R(\Omega)) = \iiint_\Omega 1 \, d\mathbf{u} = \operatorname{Vol}(\Omega). \qquad \blacksquare\]
\end{ejerciciobox}

\begin{errorbox}{¡Cuidado con esta trampa de examen!}
Intentar probar la invarianza rotando los ejes uno por uno en lugar de usar la propiedad algebraica universal $\det R = 1$.
\end{errorbox}
\end{nivelbox}
\bigskip

\begin{nivelbox}{3.5.2}{Medibilidad de Jordan y frontera no diferenciable en R³}
\textbf{Enunciado:}\\
Demuestre que el cono sólido $K = \{ (x,y,z) \in \mathbb{R}^3 : \sqrt{x^2+y^2} \le z \le 1 \}$ es Jordan-medible a pesar de que el vértice $(0,0,0)$ no es una superficie diferenciable.

\begin{observacionbox}{Pista Metodológica}
Separe la frontera en la superficie cónica regular, la tapa plana y el vértice puntual aislado, demostrando que la unión tiene contenido tridimensional cero.
\end{observacionbox}

\begin{ejerciciobox}{Resolución Analítica Paso a Paso}
\smallskip\noindent\textbf{1.} \textbf{Descomposición de la frontera:}
\[\partial K = S_{\text{cono}} \cup S_{\text{tapa}} \cup \{(0,0,0)\}.\]
- $S_{\text{tapa}} = \{ (x,y,1) : x^2 + y^2 \le 1 \}$ es un disco bidimensional suave en el plano $z=1$.
- $S_{\text{cono}} = \{ (x,y,z) : z = \sqrt{x^2+y^2}, \ 0 < z < 1 \}$ es una superficie reglada regular $C^1$.
- $P_0 = (0,0,0)$ es el vértice singular.

\smallskip\noindent\textbf{2.} \textbf{Contenido nulo de superficies suaves:}
Cualquier gráfico de una función Lipschitz continua sobre un compacto de $\mathbb{R}^2$ tiene contenido nulo en $\mathbb{R}^3$. Tanto el disco como el cono abierto tienen contenido 3D cero.

\smallskip\noindent\textbf{3.} \textbf{El vértice aislado:}
El punto $P_0$ puede ser cubierto por un cubo de arista $\varepsilon$ de volumen $\varepsilon^3 < \varepsilon$.

\smallskip\noindent\textbf{4.} \textbf{Unión finita:}
Como la unión finita de conjuntos de contenido nulo tiene contenido nulo:
\[\operatorname{cont}(\partial K) = 0.\]
Por el criterio de Jordan, $K$ es medible según Jordan en $\mathbb{R}^3$. $\blacksquare$
\end{ejerciciobox}

\begin{errorbox}{¡Cuidado con esta trampa de examen!}
Creer que la falta de plano tangente en un vértice o arista destruye la medibilidad de Jordan.
\end{errorbox}
\end{nivelbox}
\bigskip

\begin{nivelbox}{3.5.3}{Integrales iteradas no conmutables para funciones discontinuas densas}
\textbf{Enunciado:}\\
Construya una función acotada $f: [0,1]^3 \to \mathbb{R}$ tal que la integral iterada $\int_0^1 \left( \int_0^1 \left( \int_0^1 f \, dz \right) dy \right) dx$ exista pero otra permutación de integrales no exista, justificando la necesidad de integrabilidad conjunta.

\begin{observacionbox}{Pista Metodológica}
Utilice la función indicatriz de un conjunto denso en ciertas coordenadas pero proyectable en otras.
\end{observacionbox}

\begin{ejerciciobox}{Resolución Analítica Paso a Paso}
\smallskip\noindent\textbf{1.} \textbf{Construcción:}
Sea $E = \{ (x, y, z) \in [0,1]^3 : x \in \mathbb{Q}, \ y \in \mathbb{Q} \}$ y defina $f(x,y,z) = 1$ en $E$ y $0$ fuera de $E$.

\smallskip\noindent\textbf{2.} \textbf{Primera integración en $z$:}
Para cada $(x,y)$ fijo, la función $z \mapsto f(x,y,z)$ es constante (igual a $1$ si $(x,y) \in \mathbb{Q}^2$, e igual a $0$ si no). Por ende, la integral $\int_0^1 f(x,y,z) dz$ existe para todo $(x,y)$ y vale $\chi_{\mathbb{Q}^2}(x,y)$.

\smallskip\noindent\textbf{3.} \textbf{Falla en la siguiente integración:}
Al intentar integrar respecto a $y$, nos encontramos con $\int_0^1 \chi_{\mathbb{Q}^2}(x,y) dy$. Para $x \in \mathbb{Q}$, esto es la función de Dirichlet unidimensional en $y$, la cual \textbf{no es integrable Riemann}.

\smallskip\noindent\textbf{4.} \textbf{Conclusión:}
El Teorema de Fubini exige la integrabilidad de Riemann conjunta de $f$ en el paralelepípedo tridimensional completo. $\blacksquare$
\end{ejerciciobox}

\begin{errorbox}{¡Cuidado con esta trampa de examen!}
Asumir que si la primera integral unidimensional existe, las siguientes necesariamente existirán.
\end{errorbox}
\end{nivelbox}
\bigskip


\section{Tema 4: Coordenadas Cilíndricas y Esféricas en $\mathbb{R}^3$}

\subsection{Nivel 0: Calentamiento Conceptual}

\begin{nivelbox}{4.0.1}{Jacobiano de coordenadas cilíndricas y esféricas}
\textbf{Enunciado:}\\
Escriba las fórmulas de transformación y deduzca los determinantes Jacobianos para:
(a) Coordenadas cilíndricas $(r, \theta, z)$.
(b) Coordenadas esféricas $(\rho, \theta, \varphi)$.

\begin{observacionbox}{Pista Metodológica}
Cilíndricas es polar en $XY$ con $z$ inalterado ($J = r$). Esféricas: $x = \rho\sin\varphi\cos\theta, y = \rho\sin\varphi\sin\theta, z = \rho\cos\varphi$.
\end{observacionbox}

\begin{ejerciciobox}{Resolución Analítica Paso a Paso}
\smallskip\noindent\textbf{1.} \textbf{Coordenadas cilíndricas:}
\[x = r\cos\theta, \quad y = r\sin\theta, \quad z = z.\]
\[\det J_{\text{cil}} = \det \begin{pmatrix} \cos\theta & -r\sin\theta & 0 \\ \sin\theta & r\cos\theta & 0 \\ 0 & 0 & 1 \end{pmatrix} = r.\]
\[dV = r \, dr \, d\theta \, dz.\]

\smallskip\noindent\textbf{2.} \textbf{Coordenadas esféricas:}
\[x = \rho\sin\varphi\cos\theta, \quad y = \rho\sin\varphi\sin\theta, \quad z = \rho\cos\varphi.\]
Desarrollando el determinante $3 \times 3$:
\[\det J_{\text{esf}} = \rho^2 \sin\varphi.\]
\[dV = \rho^2 \sin\varphi \, d\rho \, d\varphi \, d\theta.\]
\end{ejerciciobox}

\begin{errorbox}{¡Cuidado con esta trampa de examen!}
Confundir el ángulo colatitud $\varphi \in [0, \pi]$ con la elevación geográfica desde el ecuador.
\end{errorbox}
\end{nivelbox}
\bigskip

\begin{nivelbox}{4.0.2}{Conversión de superficies notables a esféricas y cilíndricas}
\textbf{Enunciado:}\\
Exprese en coordenadas cilíndricas y esféricas las siguientes superficies:
(a) La esfera $x^2 + y^2 + z^2 = 9$.
(b) El cono circular recto $z = \sqrt{x^2 + y^2}$.
(c) El cilindro $x^2 + y^2 = 4$.

\begin{observacionbox}{Pista Metodológica}
Recuerde: $x^2+y^2 = r^2 = \rho^2\sin^2\varphi$, $x^2+y^2+z^2 = \rho^2$, y $z = \rho\cos\varphi$.
\end{observacionbox}

\begin{ejerciciobox}{Resolución Analítica Paso a Paso}
\smallskip\noindent\textbf{1.} \textbf{Esfera $x^2+y^2+z^2 = 9$:}
- Cilíndricas: $r^2 + z^2 = 9$.
- Esféricas: $\rho = 3$ (superficie de radio constante).

\smallskip\noindent\textbf{2.} \textbf{Cono $z = \sqrt{x^2+y^2}$:}
- Cilíndricas: $z = r$.
- Esféricas: $\rho\cos\varphi = \rho\sin\varphi \implies \tan\varphi = 1 \implies \varphi = \frac{\pi}{4}$ (cono de ángulo constante).

\smallskip\noindent\textbf{3.} \textbf{Cilindro $x^2+y^2 = 4$:}
- Cilíndricas: $r = 2$.
- Esféricas: $\rho\sin\varphi = 2 \implies \rho = 2\csc\varphi$.
\end{ejerciciobox}

\begin{errorbox}{¡Cuidado con esta trampa de examen!}
Escribir $\varphi = \pi/2$ para el cono $z = \sqrt{x^2+y^2}$. $\varphi = \pi/2$ representa el plano horizontal $XY$.
\end{errorbox}
\end{nivelbox}
\bigskip

\begin{nivelbox}{4.0.3}{Volumen de la bola tridimensional en esféricas}
\textbf{Enunciado:}\\
Calcule el volumen de la bola unitaria $B = \{ (x,y,z) : x^2 + y^2 + z^2 \le 1 \}$ usando coordenadas esféricas.

\begin{observacionbox}{Pista Metodológica}
Integre $\rho^2\sin\varphi \, d\rho \, d\varphi \, d\theta$ con $\rho \in [0, 1], \varphi \in [0, \pi], \theta \in [0, 2\pi]$.
\end{observacionbox}

\begin{ejerciciobox}{Resolución Analítica Paso a Paso}
\smallskip\noindent\textbf{1.} \textbf{Factorización iterada:}
\[V = \left( \int_0^{2\pi} d\theta \right) \left( \int_0^\pi \sin\varphi \, d\varphi \right) \left( \int_0^1 \rho^2 \, d\rho \right).\]

\smallskip\noindent\textbf{2.} \textbf{Evaluación:}
\[\int_0^{2\pi} d\theta = 2\pi.\]
\[\int_0^\pi \sin\varphi \, d\varphi = [-\cos\varphi]_0^\pi = -(-1) - (-1) = 2.\]
\[\int_0^1 \rho^2 \, d\rho = \left[ \frac{\rho^3}{3} \right]_0^1 = \frac{1}{3}.\]

\smallskip\noindent\textbf{3.} \textbf{Volumen:}
\[V = (2\pi)(2)\left( \frac{1}{3} \right) = \frac{4\pi}{3}.\]
\end{ejerciciobox}

\begin{errorbox}{¡Cuidado con esta trampa de examen!}
Integrar $\varphi$ de $0$ a $2\pi$, duplicando el volumen de la esfera.
\end{errorbox}
\end{nivelbox}
\bigskip

\subsection{Nivel 1: Procedimental Guiado}

\begin{nivelbox}{4.1.1}{Volumen del cono de helado en esféricas}
\textbf{Enunciado:}\\
Halle el volumen del sólido delimitado superiormente por la esfera $x^2 + y^2 + z^2 = 4$ e inferiormente por el cono $z = \sqrt{x^2 + y^2}$.

\begin{observacionbox}{Pista Metodológica}
El cono tiene ángulo $\varphi = \pi/4$ y la esfera radio $\rho = 2$. Los límites son $\rho \in [0, 2], \varphi \in [0, \pi/4], \theta \in [0, 2\pi]$.
\end{observacionbox}

\begin{ejerciciobox}{Resolución Analítica Paso a Paso}
\smallskip\noindent\textbf{1.} \textbf{Planteamiento:}
\[V = \int_0^{2\pi} \int_0^{\pi/4} \int_0^2 \rho^2 \sin\varphi \, d\rho \, d\varphi \, d\theta.\]

\smallskip\noindent\textbf{2.} \textbf{Evaluación de factores:}
\[\int_0^2 \rho^2 \, d\rho = \frac{8}{3}.\]
\[\int_0^{\pi/4} \sin\varphi \, d\varphi = [-\cos\varphi]_0^{\pi/4} = 1 - \frac{\sqrt{2}}{2}.\]
\[\int_0^{2\pi} d\theta = 2\pi.\]

\smallskip\noindent\textbf{3.} \textbf{Volumen final:}
\[V = 2\pi \left( 1 - \frac{\sqrt{2}}{2} \right) \left( \frac{8}{3} \right) = \frac{16\pi}{3} \left( 1 - \frac{\sqrt{2}}{2} \right) = \frac{8\pi(2 - \sqrt{2})}{3}.\]
\end{ejerciciobox}

\begin{errorbox}{¡Cuidado con esta trampa de examen!}
Calcularlo en cartesianas, lo cual involucra integrales dobles con $\sqrt{4-x^2-y^2} - \sqrt{x^2+y^2}$ mucho más propensas a errores algebraicos.
\end{errorbox}
\end{nivelbox}
\bigskip

\begin{nivelbox}{4.1.2}{Volumen bajo un paraboloide dentro de un cilindro en cilíndricas}
\textbf{Enunciado:}\\
Calcule el volumen del sólido delimitado por el paraboloide $z = 4 - x^2 - y^2$, el cilindro $x^2 + y^2 = 1$ y el plano $z = 0$.

\begin{observacionbox}{Pista Metodológica}
En cilíndricas: $0 \le r \le 1, 0 \le \theta \le 2\pi$ y $0 \le z \le 4 - r^2$.
\end{observacionbox}

\begin{ejerciciobox}{Resolución Analítica Paso a Paso}
\smallskip\noindent\textbf{1.} \textbf{Integral iterada:}
\[V = \int_0^{2\pi} \int_0^1 \int_0^{4-r^2} r \, dz \, dr \, d\theta = \int_0^{2\pi} d\theta \int_0^1 (4 - r^2) r \, dr.\]

\smallskip\noindent\textbf{2.} \textbf{Integración en $r$:}
\[\int_0^1 (4r - r^3) \, dr = \left[ 2r^2 - \frac{r^4}{4} \right]_0^1 = 2 - \frac{1}{4} = \frac{7}{4}.\]

\smallskip\noindent\textbf{3.} \textbf{Volumen:}
\[V = 2\pi \cdot \frac{7}{4} = \frac{7\pi}{2}.\]
\end{ejerciciobox}

\begin{errorbox}{¡Cuidado con esta trampa de examen!}
Olvidar multiplicar por el factor $r$ antes de integrar respecto a $z$.
\end{errorbox}
\end{nivelbox}
\bigskip

\begin{nivelbox}{4.1.3}{Integral de potencial en el origen}
\textbf{Enunciado:}\\
Evalúe $I = \iiint_B \frac{1}{\sqrt{x^2 + y^2 + z^2}} \, dV$, donde $B$ es la bola de radio $R$ centrada en el origen.

\begin{observacionbox}{Pista Metodológica}
En esféricas, el denominador es simplemente $\rho$.
\end{observacionbox}

\begin{ejerciciobox}{Resolución Analítica Paso a Paso}
\smallskip\noindent\textbf{1.} \textbf{Conversión a esféricas:}
\[I = \int_0^{2\pi} \int_0^\pi \int_0^R \frac{1}{\rho} \cdot \rho^2 \sin\varphi \, d\rho \, d\varphi \, d\theta = \int_0^{2\pi} \int_0^\pi \int_0^R \rho \sin\varphi \, d\rho \, d\varphi \, d\theta.\]

\smallskip\noindent\textbf{2.} \textbf{Evaluación:}
\[\left( \int_0^{2\pi} d\theta \right) \left( \int_0^\pi \sin\varphi \, d\varphi \right) \left( \int_0^R \rho \, d\rho \right) = (2\pi)(2)\left( \frac{R^2}{2} \right) = 2\pi R^2.\]
\end{ejerciciobox}

\begin{errorbox}{¡Cuidado con esta trampa de examen!}
Creer que la integral diverge por tener división por cero en el origen. El factor $\rho^2$ del Jacobiano cancela la singularidad $1/\rho$ haciéndola perfectamente integrable.
\end{errorbox}
\end{nivelbox}
\bigskip

\subsection{Nivel 2: Integración de Conceptos}

\begin{nivelbox}{4.2.1}{Volumen entre paraboloide y semicono (PC6 2025-II, Pregunta 3b)}
\textbf{Enunciado:}\\
Calcule el volumen del sólido determinado por el paraboloide $z = x^2 + y^2$ y el semicono $z = \sqrt{x^2 + y^2}$ (PC6 2025-II, Pregunta 3b).

\begin{observacionbox}{Pista Metodológica}
En cilíndricas: el paraboloide es $z = r^2$ y el cono es $z = r$. La intersección ocurre en $r^2 = r \implies r = 1$. El cono está por encima del paraboloide para $r \in [0, 1]$.
\end{observacionbox}

\begin{ejerciciobox}{Resolución Analítica Paso a Paso}
\smallskip\noindent\textbf{1.} \textbf{Intersección:}
\[r^2 = r \implies r(r - 1) = 0 \implies r = 1.\]
Para $r \in [0, 1]$, $r \ge r^2$, por lo que el cono es el techo y el paraboloide el piso: $r^2 \le z \le r$.

\smallskip\noindent\textbf{2.} \textbf{Planteamiento en coordenadas cilíndricas:}
\[V = \int_0^{2\pi} \int_0^1 \int_{r^2}^r r \, dz \, dr \, d\theta = 2\pi \int_0^1 r (r - r^2) \, dr.\]

\smallskip\noindent\textbf{3.} \textbf{Integración:}
\[V = 2\pi \int_0^1 (r^2 - r^3) \, dr = 2\pi \left[ \frac{r^3}{3} - \frac{r^4}{4} \right]_0^1 = 2\pi \left( \frac{1}{3} - \frac{1}{4} \right) = 2\pi \left( \frac{1}{12} \right) = \frac{\pi}{6}.\]
\end{ejerciciobox}

\begin{errorbox}{¡Cuidado con esta trampa de examen!}
Poner el paraboloide como techo y el cono como piso, obteniendo un volumen negativo.
\end{errorbox}
\end{nivelbox}
\bigskip

\begin{nivelbox}{4.2.2}{Integral sobre hemisferio con potencia de z (PC6 2025-II, Pregunta 3a)}
\textbf{Enunciado:}\\
Evalúe $\iiint_H z^3 \sqrt{x^2 + y^2 + z^2} \, dV$, donde $H$ es el hemisferio sólido sobre el plano $XY$ con centro en el origen y radio 1 (PC6 2025-II, Pregunta 3a).

\begin{observacionbox}{Pista Metodológica}
En esféricas: $z = \rho\cos\varphi$, $\sqrt{x^2+y^2+z^2} = \rho$. En el hemisferio superior $\varphi \in [0, \pi/2]$.
\end{observacionbox}

\begin{ejerciciobox}{Resolución Analítica Paso a Paso}
\smallskip\noindent\textbf{1.} \textbf{Conversión a coordenadas esféricas:}
\[z^3 \sqrt{x^2+y^2+z^2} = (\rho\cos\varphi)^3 \cdot \rho = \rho^4 \cos^3\varphi.\]
\[dV = \rho^2 \sin\varphi \, d\rho \, d\varphi \, d\theta.\]
El producto del integrando con el Jacobiano es: $\rho^6 \cos^3\varphi \sin\varphi$.

\smallskip\noindent\textbf{2.} \textbf{Límites:} $\rho \in [0, 1], \varphi \in [0, \pi/2], \theta \in [0, 2\pi]$.

\smallskip\noindent\textbf{3.} \textbf{Separación de integrales:}
\[I = \left( \int_0^{2\pi} d\theta \right) \left( \int_0^{\pi/2} \cos^3\varphi \sin\varphi \, d\varphi \right) \left( \int_0^1 \rho^6 \, d\rho \right).\]
- $\int_0^{2\pi} d\theta = 2\pi$.
- $\int_0^{\pi/2} \cos^3\varphi \sin\varphi \, d\varphi = \left[ -\frac{\cos^4\varphi}{4} \right]_0^{\pi/2} = 0 - \left(-\frac{1}{4}\right) = \frac{1}{4}$.
- $\int_0^1 \rho^6 \, d\rho = \frac{1}{7}$.

\smallskip\noindent\textbf{4.} \textbf{Resultado:}
\[I = 2\pi \cdot \frac{1}{4} \cdot \frac{1}{7} = \frac{\pi}{14}.\]
\end{ejerciciobox}

\begin{errorbox}{¡Cuidado con esta trampa de examen!}
Integrar $\varphi$ hasta $\pi$, olvidando que solo se trata del hemisferio superior ($z \ge 0$).
\end{errorbox}
\end{nivelbox}
\bigskip

\begin{nivelbox}{4.2.3}{Masa de un cilindro con densidad proporcional a la distancia al eje}
\textbf{Enunciado:}\\
Halle la masa de un cilindro macizo de radio $R$ y altura $H$, si la densidad volumétrica en cada punto es proporcional a su distancia al eje $Z$: $\rho(x,y,z) = k \sqrt{x^2+y^2}$.

\begin{observacionbox}{Pista Metodológica}
En cilíndricas, $\rho(r) = k r$. Integre $\iiint k r \cdot r \, dr \, d\theta \, dz$.
\end{observacionbox}

\begin{ejerciciobox}{Resolución Analítica Paso a Paso}
\smallskip\noindent\textbf{1.} \textbf{Integral en coordenadas cilíndricas:}
\[M = \int_0^H \int_0^{2\pi} \int_0^R (k r) \cdot r \, dr \, d\theta \, dz = k \left( \int_0^H dz \right) \left( \int_0^{2\pi} d\theta \right) \left( \int_0^R r^2 \, dr \right).\]

\smallskip\noindent\textbf{2.} \textbf{Evaluación de factores:}
\[M = k \cdot H \cdot 2\pi \cdot \frac{R^3}{3} = \frac{2\pi k H R^3}{3}.\]
\end{ejerciciobox}

\begin{errorbox}{¡Cuidado con esta trampa de examen!}
Integrar $k r$ en vez de $k r^2$ al descuidar el factor $r$ del elemento de volumen cilíndrico.
\end{errorbox}
\end{nivelbox}
\bigskip

\subsection{Nivel 3: Problemas Tipo Examen UNI}

\begin{nivelbox}{4.3.1}{Potencial en la bola unitaria con punto exterior (PC5 2024-II, Pregunta 4)}
\textbf{Enunciado:}\\
Sea $B$ la bola unitaria $x^2 + y^2 + z^2 \le 1$ y sea $a > 1$. Calcule:
\[I = \iiint_B \frac{dx \, dy \, dz}{\sqrt{x^2 + y^2 + (z - a)^2}}.\]
(PC5 2024-II, Pregunta 4).

\begin{observacionbox}{Pista Metodológica}
Pase a coordenadas esféricas con $x = \rho\sin\varphi\cos\theta, y = \rho\sin\varphi\sin\theta, z = \rho\cos\varphi$. En el radicando, $x^2+y^2+(z-a)^2 = \rho^2 + a^2 - 2a\rho\cos\varphi$. Realice la sustitución $u = \rho^2 + a^2 - 2a\rho\cos\varphi$.
\end{observacionbox}

\begin{ejerciciobox}{Resolución Analítica Paso a Paso}
\smallskip\noindent\textbf{1.} \textbf{Conversión a esféricas:}
\[x^2 + y^2 + (z - a)^2 = (x^2+y^2+z^2) - 2az + a^2 = \rho^2 + a^2 - 2a\rho\cos\varphi.\]
\[I = \int_0^1 \int_0^{2\pi} \int_0^\pi \frac{\rho^2 \sin\varphi}{\sqrt{\rho^2 + a^2 - 2a\rho\cos\varphi}} \, d\varphi \, d\theta \, d\rho = 2\pi \int_0^1 \rho^2 \left( \int_0^\pi \frac{\sin\varphi}{\sqrt{\rho^2 + a^2 - 2a\rho\cos\varphi}} \, d\varphi \right) d\rho.\]

\smallskip\noindent\textbf{2.} \textbf{Sustitución en la integral de $\varphi$:}
Sea $u = \rho^2 + a^2 - 2a\rho\cos\varphi \implies du = 2a\rho\sin\varphi \, d\varphi$.
- Para $\varphi = 0$: $u = \rho^2 + a^2 - 2a\rho = (a - \rho)^2$.
- Para $\varphi = \pi$: $u = \rho^2 + a^2 + 2a\rho = (a + \rho)^2$.
Como $a > 1 \ge \rho > 0$, $\sqrt{(a - \rho)^2} = a - \rho$.
\[\int_{(a-\rho)^2}^{(a+\rho)^2} \frac{du}{2a\rho \sqrt{u}} = \frac{1}{2a\rho} [2\sqrt{u}]_{(a-\rho)^2}^{(a+\rho)^2} = \frac{1}{a\rho} [ (a + \rho) - (a - \rho) ] = \frac{2\rho}{a\rho} = \frac{2}{a}.\]

\smallskip\noindent\textbf{3.} \textbf{Integración exterior en $\rho$:}
\[I = 2\pi \int_0^1 \rho^2 \left( \frac{2}{a} \right) d\rho = \frac{4\pi}{a} \int_0^1 \rho^2 \, d\rho = \frac{4\pi}{a} \left( \frac{1}{3} \right) = \frac{4\pi}{3a}. \qquad \blacksquare\]
\end{ejerciciobox}

\begin{errorbox}{¡Cuidado con esta trampa de examen!}
Tomar $\sqrt{(a-\rho)^2} = \rho - a$ en vez de $a - \rho$. Como $a > 1$ y $\rho \le 1$, $a - \rho > 0$ es estrictamente positivo.
\end{errorbox}
\end{nivelbox}
\bigskip

\begin{nivelbox}{4.3.2}{Esfera tangente al origen en coordenadas esféricas}
\textbf{Enunciado:}\\
Calcule el volumen del sólido acotado por la esfera $x^2 + y^2 + z^2 = 2R z$ usando coordenadas esféricas.

\begin{observacionbox}{Pista Metodológica}
En esféricas, $\rho^2 = 2R \rho\cos\varphi \implies \rho = 2R\cos\varphi$. Para que $\rho \ge 0$, $\varphi \in [0, \pi/2]$.
\end{observacionbox}

\begin{ejerciciobox}{Resolución Analítica Paso a Paso}
\smallskip\noindent\textbf{1.} \textbf{Ecuación polar de la esfera:}
\[\rho = 2R\cos\varphi, \qquad 0 \le \varphi \le \frac{\pi}{2}, \quad 0 \le \theta \le 2\pi.\]

\smallskip\noindent\textbf{2.} \textbf{Integral de volumen:}
\[V = \int_0^{2\pi} \int_0^{\pi/2} \int_0^{2R\cos\varphi} \rho^2 \sin\varphi \, d\rho \, d\varphi \, d\theta.\]

\smallskip\noindent\textbf{3.} \textbf{Integración en $\rho$:}
\[\int_0^{2R\cos\varphi} \rho^2 \, d\rho = \frac{8R^3 \cos^3\varphi}{3}.\]

\smallskip\noindent\textbf{4.} \textbf{Integración en $\varphi$ y $\theta$:}
\[V = 2\pi \cdot \frac{8R^3}{3} \int_0^{\pi/2} \cos^3\varphi \sin\varphi \, d\varphi = \frac{16\pi R^3}{3} \left[ -\frac{\cos^4\varphi}{4} \right]_0^{\pi/2} = \frac{16\pi R^3}{3} \left( \frac{1}{4} \right) = \frac{4\pi R^3}{3}.\]
\end{ejerciciobox}

\begin{errorbox}{¡Cuidado con esta trampa de examen!}
Integrar $\varphi$ hasta $\pi$, lo cual produce valores negativos de radio.
\end{errorbox}
\end{nivelbox}
\bigskip

\begin{nivelbox}{4.3.3}{Mapeo que cubre la región varias veces (PC6 2025-II, Pregunta 1d)}
\textbf{Enunciado:}\\
Demuestre formalmente que la siguiente afirmación de examen es FALSA: 'Si en coordenadas cilíndricas tomo el ángulo $\theta$ en un intervalo más largo, por ejemplo $[0, 4\pi]$, no ocurre nada y la fórmula sigue dando el mismo valor aunque el mapeo cubra la región varias veces' (PC6 2025-II, Pregunta 1d).

\begin{observacionbox}{Pista Metodológica}
El Teorema de Cambio de Variables exige que la transformación sea inyectiva casi en todas partes (salvo conjuntos de contenido nulo).
\end{observacionbox}

\begin{ejerciciobox}{Resolución Analítica Paso a Paso}
\smallskip\noindent\textbf{1.} \textbf{Análisis de inyectividad:}
La aplicación polar/cilíndrica $T(r, \theta, z) = (r\cos\theta, r\sin\theta, z)$ con $\theta \in [0, 4\pi]$ es \textbf{2 a 1} casi en todas partes: los puntos $(r, \theta_0, z)$ y $(r, \theta_0 + 2\pi, z)$ tienen exactamente la misma imagen cartesiana.

\smallskip\noindent\textbf{2.} \textbf{Efecto en la integral:}
Para una función integrable no negativa $f \ge 0$ (o $f = 1$):
\[\int_0^{4\pi} \iint_{D'} f \, r \, dr dz d\theta = \int_0^{2\pi} \dots + \int_{2\pi}^{4\pi} \dots = 2 \iiint_D f(x,y,z) \, dV.\]
La integral calculada con $\theta \in [0, 4\pi]$ resulta ser exactamente \textbf{el doble} del valor real.
Por tanto, la afirmación es rotunda y formalmente \textbf{FALSA}. $\blacksquare$
\end{ejerciciobox}

\begin{errorbox}{¡Cuidado con esta trampa de examen!}
Creer que la función trigonométrica 'se auto-cancela'. El Jacobiano $r$ es positivo y la integral se duplica o multiplica por el grado de recubrimiento.
\end{errorbox}
\end{nivelbox}
\bigskip

\subsection{Nivel 4: Desafíos Teóricos y Contraejemplos}

\begin{nivelbox}{4.4.1}{Teorema de las Capas Esféricas de Newton (Shell Theorem)}
\textbf{Enunciado:}\\
Demuestre mediante integrales triples que una capa esférica homogénea de masa $M$ y radio $R$ no ejerce fuerza gravitatoria neta sobre ninguna partícula situada en su interior ($r_0 < R$).

\begin{observacionbox}{Pista Metodológica}
Sitúe el punto de prueba en $(0,0,c)$ con $0 \le c < R$. Calcule la integral del potencial gravitatorio en la superficie o capa delgada.
\end{observacionbox}

\begin{ejerciciobox}{Resolución Analítica Paso a Paso}
\smallskip\noindent\textbf{1.} \textbf{Potencial gravitatorio:}
\[V(c) = -G \sigma \iint_{S_R} \frac{1}{\sqrt{R^2 + c^2 - 2Rc\cos\varphi}} \, dA.\]
\[dA = R^2 \sin\varphi \, d\varphi \, d\theta.\]

\smallskip\noindent\textbf{2.} \textbf{Integración angular:}
\[V(c) = -G \sigma (2\pi R^2) \int_0^\pi \frac{\sin\varphi}{\sqrt{R^2 + c^2 - 2Rc\cos\varphi}} \, d\varphi.\]
Sustituyendo $u = R^2 + c^2 - 2Rc\cos\varphi, du = 2Rc\sin\varphi d\varphi$:
\[\int_{(R-c)^2}^{(R+c)^2} \frac{du}{2Rc \sqrt{u}} = \frac{1}{Rc} [ (R + c) - (R - c) ] = \frac{2c}{Rc} = \frac{2}{R}.\]

\smallskip\noindent\textbf{3.} \textbf{Constancia del potencial:}
\[V(c) = -G \sigma (2\pi R^2) \left( \frac{2}{R} \right) = -G \frac{4\pi R^2 \sigma}{R} = -\frac{GM}{R}.\]
Como el potencial gravitatorio $V(c)$ es \textbf{estrictamente constante} en todo el interior ($c < R$), la fuerza gravitatoria $\mathbf{F} = -\nabla V = \mathbf{0}$. $\blacksquare$
\end{ejerciciobox}

\begin{errorbox}{¡Cuidado con esta trampa de examen!}
Usar $c - R$ al extraer la raíz cuadrada, olvidando que en el interior $R > c$.
\end{errorbox}
\end{nivelbox}
\bigskip

\begin{nivelbox}{4.4.2}{Criterio de integrabilidad de singularidades potenciales en R³}
\textbf{Enunciado:}\\
Sea $B = \{ \mathbf{x} \in \mathbb{R}^3 : \|\mathbf{x}\| \le 1 \}$. Demuestre que la integral impropia $\iiint_B \frac{1}{\|\mathbf{x}\|^\alpha} \, dV$ converge si y solo si $\alpha < 3$.

\begin{observacionbox}{Pista Metodológica}
En coordenadas esféricas, el elemento de volumen añade un factor $\rho^2$, transformando el integrando radial en $\rho^{2 - \alpha}$.
\end{observacionbox}

\begin{ejerciciobox}{Resolución Analítica Paso a Paso}
\smallskip\noindent\textbf{1.} \textbf{Paso a coordenadas esféricas:}
\[I = \int_0^{2\pi} d\theta \int_0^\pi \sin\varphi \, d\varphi \int_0^1 \frac{1}{\rho^\alpha} \rho^2 \, d\rho = 4\pi \int_0^1 \rho^{2 - \alpha} \, d\rho.\]

\smallskip\noindent\textbf{2.} \textbf{Convergencia de la integral de Riemann impropia:}
Por el análisis real 1D, $\int_0^1 \rho^p d\rho$ converge $\iff p > -1$.
Aquí $p = 2 - \alpha$, luego la condición de convergencia es:
\[2 - \alpha > -1 \iff -\alpha > -3 \iff \alpha < 3.\]

\smallskip\noindent\textbf{3.} \textbf{Cálculo del valor cuando $\alpha < 3$:}
\[I = 4\pi \left[ \frac{\rho^{3 - \alpha}}{3 - \alpha} \right]_0^1 = \frac{4\pi}{3 - \alpha}. \qquad \blacksquare\]
\end{ejerciciobox}

\begin{errorbox}{¡Cuidado con esta trampa de examen!}
Confundir el criterio tridimensional ($\alpha < 3$) con el unidimensional ($\alpha < 1$) o bidimensional ($\alpha < 2$).
\end{errorbox}
\end{nivelbox}
\bigskip

\begin{nivelbox}{4.4.3}{Factores de escala y matriz métrica curvilínea ortogonal}
\textbf{Enunciado:}\\
Para un sistema de coordenadas ortogonales $(u_1, u_2, u_3)$ con factores de escala $h_i = \|\frac{\partial \mathbf{r}}{\partial u_i}\|$, demuestre que el determinante Jacobiano satisface $|\det J| = h_1 h_2 h_3$. Corrobore este resultado para coordenadas esféricas.

\begin{observacionbox}{Pista Metodológica}
La matriz de Gram es $G = J^T J = \operatorname{diag}(h_1^2, h_2^2, h_3^2)$.
\end{observacionbox}

\begin{ejerciciobox}{Resolución Analítica Paso a Paso}
\smallskip\noindent\textbf{1.} \textbf{Ortogonalidad de los vectores tangentes:}
Por definición de coordenadas ortogonales, $\frac{\partial \mathbf{r}}{\partial u_i} \cdot \frac{\partial \mathbf{r}}{\partial u_j} = 0$ para $i \neq j$.
Por tanto, la matriz de Gram satisface:
\[G = J^T J = \begin{pmatrix} h_1^2 & 0 & 0 \\ 0 & h_2^2 & 0 \\ 0 & 0 & h_3^2 \end{pmatrix}.\]

\smallskip\noindent\textbf{2.} \textbf{Determinante Jacobiano:}
\[(\det J)^2 = \det(J^T J) = \det G = h_1^2 h_2^2 h_3^2 \implies |\det J| = h_1 h_2 h_3.\]

\smallskip\noindent\textbf{3.} \textbf{Verificación en esféricas ($u_1 = \rho, u_2 = \varphi, u_3 = \theta$):}
- $h_\rho = \|(\sin\varphi\cos\theta, \sin\varphi\sin\theta, \cos\varphi)\| = 1$.
- $h_\varphi = \|(\rho\cos\varphi\cos\theta, \rho\cos\varphi\sin\theta, -\rho\sin\varphi)\| = \rho$.
- $h_\theta = \|(-\rho\sin\varphi\sin\theta, \rho\sin\varphi\cos\theta, 0)\| = \rho\sin\varphi$.
\[|\det J| = h_\rho h_\varphi h_\theta = (1)(\rho)(\rho\sin\varphi) = \rho^2 \sin\varphi. \quad \blacksquare\]
\end{ejerciciobox}

\begin{errorbox}{¡Cuidado con esta trampa de examen!}
Aplicar esta fórmula a coordenadas que no sean ortogonales (donde los productos cruzados no se anulan).
\end{errorbox}
\end{nivelbox}
\bigskip

\subsection{Nivel 5: Dominio Avanzado y Rigor Topológico}

\begin{nivelbox}{4.5.1}{Deducción formal rigurosa del Jacobiano esférico mediante formas diferenciales}
\textbf{Enunciado:}\\
Demuestre la fórmula del elemento de volumen esférico $dx \wedge dy \wedge dz = \rho^2 \sin\varphi \, d\rho \wedge d\varphi \wedge d\theta$ utilizando el álgebra exterior de formas diferenciales en $\mathbb{R}^3$.

\begin{observacionbox}{Pista Metodológica}
Calcule los diferenciales exteriores $dx, dy, dz$ y use las propiedades $d\alpha \wedge d\alpha = 0$ y $d\alpha \wedge d\beta = -d\beta \wedge d\alpha$.
\end{observacionbox}

\begin{ejerciciobox}{Resolución Analítica Paso a Paso}
\smallskip\noindent\textbf{1.} \textbf{Diferenciales de las coordenadas cartesianas:}
\[dx = \sin\varphi\cos\theta \, d\rho + \rho\cos\varphi\cos\theta \, d\varphi - \rho\sin\varphi\sin\theta \, d\theta,\]
\[dy = \sin\varphi\sin\theta \, d\rho + \rho\cos\varphi\sin\theta \, d\varphi + \rho\sin\varphi\cos\theta \, d\theta,\]
\[dz = \cos\varphi \, d\rho - \rho\sin\varphi \, d\varphi.\]

\smallskip\noindent\textbf{2.} \textbf{Producto exterior $dx \wedge dy$:}
Efectuando el producto y usando la antisimetría:
\[dx \wedge dy = \rho\sin^2\varphi \, d\rho \wedge d\theta + \rho^2 \sin\varphi\cos\varphi \, d\varphi \wedge d\theta.\]

\smallskip\noindent\textbf{3.} \textbf{Producto exterior completo con $dz$:}
\[dx \wedge dy \wedge dz = (\rho\sin^2\varphi \, d\rho \wedge d\theta) \wedge (-\rho\sin\varphi \, d\varphi) + (\rho^2 \sin\varphi\cos\varphi \, d\varphi \wedge d\theta) \wedge (\cos\varphi \, d\rho).\]
Reordenando las 1-formas a la base canónica $d\rho \wedge d\varphi \wedge d\theta$:
\[= \rho^2 \sin^3\varphi \, d\rho \wedge d\varphi \wedge d\theta + \rho^2 \sin\varphi\cos^2\varphi \, d\rho \wedge d\varphi \wedge d\theta\]
\[= \rho^2 \sin\varphi (\sin^2\varphi + \cos^2\varphi) \, d\rho \wedge d\varphi \wedge d\theta = \rho^2 \sin\varphi \, d\rho \wedge d\varphi \wedge d\theta. \quad \blacksquare\]
\end{ejerciciobox}

\begin{errorbox}{¡Cuidado con esta trampa de examen!}
Cometer errores de signos al permutar las 1-formas diferenciales.
\end{errorbox}
\end{nivelbox}
\bigskip

\begin{nivelbox}{4.5.2}{Volumen del Sólido de Viviani}
\textbf{Enunciado:}\\
Halle el volumen del sólido comprendido dentro de la esfera $x^2 + y^2 + z^2 \le 4R^2$ y dentro del cilindro de Viviani $x^2 + y^2 \le 2R x$.

\begin{observacionbox}{Pista Metodológica}
En coordenadas cilíndricas, el cilindro es $r \le 2R\cos\theta$ con $\theta \in [-\pi/2, \pi/2]$. La altura de la esfera es $z = \pm \sqrt{4R^2 - r^2}$.
\end{observacionbox}

\begin{ejerciciobox}{Resolución Analítica Paso a Paso}
\smallskip\noindent\textbf{1.} \textbf{Integral en coordenadas cilíndricas:}
Por simetría, cuadruplicamos la integral en el primer octante ($z \ge 0, \theta \in [0, \pi/2]$):
\[V = 4 \int_0^{\pi/2} \int_0^{2R\cos\theta} \sqrt{4R^2 - r^2} \, r \, dr \, d\theta.\]

\smallskip\noindent\textbf{2.} \textbf{Integración en $r$:}
\[\int_0^{2R\cos\theta} \sqrt{4R^2 - r^2} \, r \, dr = \left[ -\frac{1}{3}(4R^2 - r^2)^{3/2} \right]_0^{2R\cos\theta} = \frac{1}{3} \left( 8R^3 - (4R^2 - 4R^2\cos^2\theta)^{3/2} \right)\]
\[= \frac{8R^3}{3} (1 - \sin^3\theta).\]

\smallskip\noindent\textbf{3.} \textbf{Integración en $\theta$:}
\[V = 4 \cdot \frac{8R^3}{3} \int_0^{\pi/2} (1 - \sin^3\theta) \, d\theta.\]
Como $\int_0^{\pi/2} 1 \, d\theta = \frac{\pi}{2}$ y $\int_0^{\pi/2} \sin^3\theta \, d\theta = \frac{2}{3}$:
\[V = \frac{32R^3}{3} \left( \frac{\pi}{2} - \frac{2}{3} \right) = \frac{16\pi R^3}{3} - \frac{64R^3}{9} = \frac{16R^3}{9} (3\pi - 4). \quad \blacksquare\]
\end{ejerciciobox}

\begin{errorbox}{¡Cuidado con esta trampa de examen!}
Escribir $\sqrt{4R^2 - 4R^2\cos^2\theta} = 2R\cos\theta$ en vez de $2R\sin\theta$.
\end{errorbox}
\end{nivelbox}
\bigskip

\begin{nivelbox}{4.5.3}{Conjuntos singulares de medida nula en transformaciones esféricas}
\textbf{Enunciado:}\\
Demuestre formalmente que el conjunto de puntos donde la transformación esférica $T: [0,R] \times [0,\pi] \times [0,2\pi] \to \mathbb{R}^3$ no es un difeomorfismo local tiene contenido de Jordan tridimensional nulo.

\begin{observacionbox}{Pista Metodológica}
Identifique los puntos donde $\det DT = 0$: $\rho = 0$, $\varphi = 0$, $\varphi = \pi$, y demuestre que sus imágenes son el origen y el eje $Z$.
\end{observacionbox}

\begin{ejerciciobox}{Resolución Analítica Paso a Paso}
\smallskip\noindent\textbf{1.} \textbf{Conjunto singular en el dominio:}
$\det DT = \rho^2 \sin\varphi = 0 \iff \rho = 0 \lor \sin\varphi = 0$.
En el prisma paramétrico $P = [0,R] \times [0,\pi] \times [0,2\pi]$, este conjunto es la unión de 3 caras rectangulares bidimensionales, las cuales tienen volumen tridimensional cero.

\smallskip\noindent\textbf{2.} \textbf{Imagen cartesiana del conjunto singular:}
- Para $\rho = 0$: $T(0, \varphi, \theta) = (0,0,0)$ (un solo punto en $\mathbb{R}^3$).
- Para $\varphi = 0$: $x = 0, y = 0, z = \rho$ (el segmento de recta $[0, R]$ sobre el semieje $+Z$).
- Para $\varphi = \pi$: $x = 0, y = 0, z = -\rho$ (el segmento $[-R, 0]$ sobre el semieje $-Z$).
La unión de estas imágenes es el segmento $S = \{ (0,0,z) : -R \le z \le R \}$.

\smallskip\noindent\textbf{3.} \textbf{Prueba de contenido nulo:}
Un segmento de recta en $\mathbb{R}^3$ puede ser cubierto por cilindros de radio $\varepsilon$ y altura $2R$ con volumen total $2\pi \varepsilon^2 R \to 0$.
Por ende, $\operatorname{cont}(S) = 0$.

\smallskip\noindent\textbf{4.} \textbf{Conclusión:}
El Teorema de Cambio de Variables de Lebesgue-Jordan es válido en toda la bola cerrada. $\blacksquare$
\end{ejerciciobox}

\begin{errorbox}{¡Cuidado con esta trampa de examen!}
Pensar que porque la cara $\rho = 0$ es un rectángulo bidimensional en el espacio paramétrico, su imagen en $\mathbb{R}^3$ tiene área. La imagen es un punto adimensional.
\end{errorbox}
\end{nivelbox}
\bigskip


\section{Tema 5: Integrales en $\mathbb{R}^n$ y Aplicaciones Físico-Geométricas}

\subsection{Nivel 0: Calentamiento Conceptual}

\begin{nivelbox}{5.0.1}{Masa total y densidad constante}
\textbf{Enunciado:}\\
Una placa rectangular ocupa la región $[0, 3] \times [0, 2]$ con densidad superficial uniforme $\rho_0 = 5\text{ kg/m}^2$. Calcule su masa total y su centro de masa geométrico.

\begin{observacionbox}{Pista Metodológica}
La masa es $\iint \rho \, dA = \rho_0 \operatorname{Area}$. El centroide de un rectángulo simétrico es su centro geométrico.
\end{observacionbox}

\begin{ejerciciobox}{Resolución Analítica Paso a Paso}
\smallskip\noindent\textbf{1.} \textbf{Masa total:}
\[M = \iint_R 5 \, dA = 5 \cdot \operatorname{Area}(R) = 5 \cdot (3 \times 2) = 30\text{ kg}.\]

\smallskip\noindent\textbf{2.} \textbf{Centro de masa:}
Por simetría rectangular:
\[\bar{x} = \frac{0 + 3}{2} = 1.5\text{ m}, \qquad \bar{y} = \frac{0 + 2}{2} = 1.0\text{ m}.\]
\[(\bar{x}, \bar{y}) = (1.5, 1.0).\]
\end{ejerciciobox}

\begin{errorbox}{¡Cuidado con esta trampa de examen!}
Calcular integrales complejas innecesarias para figuras con doble simetría y densidad constante.
\end{errorbox}
\end{nivelbox}
\bigskip

\begin{nivelbox}{5.0.2}{Centroide de una región triangular estándar}
\textbf{Enunciado:}\\
Halle el centroide geométrico $(\bar{x}, \bar{y})$ del triángulo acotado por $x = 0, y = 0$ y $x + y = 3$.

\begin{observacionbox}{Pista Metodológica}
Recuerde que el centroide de un triángulo rectángulo en el origen está a $1/3$ de la longitud de cada cateto: $(a/3, b/3)$.
\end{observacionbox}

\begin{ejerciciobox}{Resolución Analítica Paso a Paso}
\smallskip\noindent\textbf{1.} \textbf{Área:}
\[A = \frac{3 \times 3}{2} = \frac{9}{2}.\]

\smallskip\noindent\textbf{2.} \textbf{Momento estático $M_y$:}
\[M_y = \int_0^3 \int_0^{3-x} x \, dy \, dx = \int_0^3 x(3 - x) \, dx = \left[ \frac{3x^2}{2} - \frac{x^3}{3} \right]_0^3 = \frac{27}{2} - 9 = \frac{9}{2}.\]
\[\bar{x} = \frac{M_y}{A} = \frac{9/2}{9/2} = 1.\]

\smallskip\noindent\textbf{3.} \textbf{Por simetría:} $\bar{y} = 1$.
\[(\bar{x}, \bar{y}) = (1, 1).\]
\end{ejerciciobox}

\begin{errorbox}{¡Cuidado con esta trampa de examen!}
Asumir que el centroide está en el punto medio $(1.5, 1.5)$, lo cual corresponde a un rectángulo.
\end{errorbox}
\end{nivelbox}
\bigskip

\begin{nivelbox}{5.0.3}{Momento de inercia polar elemental de un disco}
\textbf{Enunciado:}\\
Calcule el momento de inercia polar $I_0 = \iint_D (x^2 + y^2) \, dA$ de un disco homogéneo de radio $R$ y masa total $M$.

\begin{observacionbox}{Pista Metodológica}
Pase a coordenadas polares $x^2+y^2 = r^2$ y exprese el resultado en términos de la masa $M = \rho \pi R^2$.
\end{observacionbox}

\begin{ejerciciobox}{Resolución Analítica Paso a Paso}
\smallskip\noindent\textbf{1.} \textbf{Densidad uniforme:} $\rho = \frac{M}{\pi R^2}$.

\smallskip\noindent\textbf{2.} \textbf{Integral de inercia polar:}
\[I_0 = \rho \int_0^{2\pi} \int_0^R r^2 \cdot r \, dr \, d\theta = 2\pi \rho \left[ \frac{r^4}{4} \right]_0^R = 2\pi \rho \frac{R^4}{4} = \frac{\pi \rho R^4}{2}.\]

\smallskip\noindent\textbf{3.} \textbf{Sustitución de masa:}
\[I_0 = \frac{1}{2} (\pi \rho R^2) R^2 = \frac{1}{2} M R^2.\]
\end{ejerciciobox}

\begin{errorbox}{¡Cuidado con esta trampa de examen!}
Olvidar el factor $1/2$ que surge de integrar $r^3$.
\end{errorbox}
\end{nivelbox}
\bigskip

\subsection{Nivel 1: Procedimental Guiado}

\begin{nivelbox}{5.1.1}{Masa de lámina con densidad superficial lineal}
\textbf{Enunciado:}\\
Calcule la masa de una lámina triangular delimitada por los vértices $(0,0), (2,0), (0,2)$, cuya densidad superficial es $\rho(x,y) = 2x + y$.

\begin{observacionbox}{Pista Metodológica}
La hipotenusa es $x + y = 2 \implies y = 2 - x$. Integre $\int_0^2 \int_0^{2-x} (2x + y) dy dx$.
\end{observacionbox}

\begin{ejerciciobox}{Resolución Analítica Paso a Paso}
\smallskip\noindent\textbf{1.} \textbf{Planteamiento:}
\[M = \int_0^2 \int_0^{2-x} (2x + y) \, dy \, dx.\]

\smallskip\noindent\textbf{2.} \textbf{Integral en $y$:}
\[\int_0^{2-x} (2x + y) \, dy = \left[ 2xy + \frac{y^2}{2} \right]_0^{2-x} = 2x(2 - x) + \frac{(2 - x)^2}{2} = 4x - 2x^2 + 2 - 2x + \frac{x^2}{2} = 2 + 2x - \frac{3x^2}{2}.\]

\smallskip\noindent\textbf{3.} \textbf{Integral en $x$:}
\[M = \int_0^2 \left( 2 + 2x - \frac{3x^2}{2} \right) dx = \left[ 2x + x^2 - \frac{x^3}{2} \right]_0^2 = 4 + 4 - 4 = 4.\]
\end{ejerciciobox}

\begin{errorbox}{¡Cuidado con esta trampa de examen!}
Confundir $2x+y$ al expandir algebraicamente el término cuadrático $(2-x)^2/2$.
\end{errorbox}
\end{nivelbox}
\bigskip

\begin{nivelbox}{5.1.2}{Centroide de un semicírculo homogéneo}
\textbf{Enunciado:}\\
Demuestre formalmente que la coordenada vertical del centroide de una placa homogénea semicircular $x^2 + y^2 \le R^2, y \ge 0$ es $\bar{y} = \frac{4R}{3\pi}$.

\begin{observacionbox}{Pista Metodológica}
Calcule $M_x = \iint y \, dA$ en coordenadas polares ($y = r\sin\theta$) y divida entre el área $\pi R^2 / 2$.
\end{observacionbox}

\begin{ejerciciobox}{Resolución Analítica Paso a Paso}
\smallskip\noindent\textbf{1.} \textbf{Área del semicírculo:} $A = \frac{\pi R^2}{2}$.

\smallskip\noindent\textbf{2.} \textbf{Momento estático $M_x$:}
\[M_x = \int_0^\pi \int_0^R (r\sin\theta) r \, dr \, d\theta = \left( \int_0^\pi \sin\theta \, d\theta \right) \left( \int_0^R r^2 \, dr \right).\]
\[\int_0^\pi \sin\theta \, d\theta = 2, \qquad \int_0^R r^2 \, dr = \frac{R^3}{3} \implies M_x = \frac{2R^3}{3}.\]

\smallskip\noindent\textbf{3.} \textbf{Cálculo de $\bar{y}$:}
\[\bar{y} = \frac{M_x}{A} = \frac{2R^3 / 3}{\pi R^2 / 2} = \frac{4R}{3\pi}. \qquad \blacksquare\]
\end{ejerciciobox}

\begin{errorbox}{¡Cuidado con esta trampa de examen!}
Integrar $\theta$ de $0$ a $2\pi$, lo cual daría $M_x = 0$ por anulación simétrica de la circunferencia completa.
\end{errorbox}
\end{nivelbox}
\bigskip

\begin{nivelbox}{5.1.3}{Valor promedio de una función sobre un disco}
\textbf{Enunciado:}\\
Calcule el valor promedio de la función $f(x,y) = x^2 + y^2$ sobre el disco $D = \{ (x,y) : x^2 + y^2 \le R^2 \}$.

\begin{observacionbox}{Pista Metodológica}
$f_{\text{prom}} = \frac{1}{\operatorname{Area}(D)} \iint_D f \, dA$.
\end{observacionbox}

\begin{ejerciciobox}{Resolución Analítica Paso a Paso}
\smallskip\noindent\textbf{1.} \textbf{Área del disco:} $\operatorname{Area}(D) = \pi R^2$.

\smallskip\noindent\textbf{2.} \textbf{Integral de la función:}
\[\iint_D (x^2 + y^2) \, dA = \int_0^{2\pi} \int_0^R r^2 \cdot r \, dr \, d\theta = 2\pi \left[ \frac{r^4}{4} \right]_0^R = \frac{\pi R^4}{2}.\]

\smallskip\noindent\textbf{3.} \textbf{Valor promedio:}
\[f_{\text{prom}} = \frac{\frac{\pi R^4}{2}}{\pi R^2} = \frac{R^2}{2}.\]
\end{ejerciciobox}

\begin{errorbox}{¡Cuidado con esta trampa de examen!}
Olvidar dividir entre el área del dominio de integración.
\end{errorbox}
\end{nivelbox}
\bigskip

\subsection{Nivel 2: Integración de Conceptos}

\begin{nivelbox}{5.2.1}{Centro de masa de una lámina delgada con densidad variable (PC5 2025-II, Pregunta 4)}
\textbf{Enunciado:}\\
Considere una lámina delgada cuya densidad superficial está dada por $\rho(x,y) = x + y$, ocupando la región $D = \{ (x,y) \in \mathbb{R}^2 : 0 \le x \le 2, \ 0 \le y \le \sqrt{4 - x^2} \}$ (primer cuadrante del disco de radio 2). Calcule las coordenadas exactas de su centro de masa $(\bar{x}, \bar{y})$ (PC5 2025-II, Pregunta 4).

\begin{observacionbox}{Pista Metodológica}
Use coordenadas polares: $x = r\cos\theta, y = r\sin\theta$, con $r \in [0, 2], \theta \in [0, \pi/2]$. Por simetría entre $x$ e $y$, $\bar{x} = \bar{y}$.
\end{observacionbox}

\begin{ejerciciobox}{Resolución Analítica Paso a Paso}
\smallskip\noindent\textbf{1.} \textbf{Masa total:}
\[M = \int_0^{\pi/2} \int_0^2 (r\cos\theta + r\sin\theta) r \, dr \, d\theta = \left( \int_0^{\pi/2} (\cos\theta + \sin\theta) d\theta \right) \left( \int_0^2 r^2 \, dr \right).\]
- $\int_0^{\pi/2} (\cos\theta + \sin\theta) d\theta = [\sin\theta - \cos\theta]_0^{\pi/2} = (1 - 0) - (0 - 1) = 2$.
- $\int_0^2 r^2 dr = \frac{8}{3}$.
\[M = 2 \cdot \frac{8}{3} = \frac{16}{3}.\]

\smallskip\noindent\textbf{2.} \textbf{Momento estático $M_y = \iint x \rho \, dA$:}
\[M_y = \int_0^{\pi/2} \int_0^2 (r\cos\theta)(r\cos\theta + r\sin\theta) r \, dr \, d\theta = \int_0^{\pi/2} (\cos^2\theta + \sin\theta\cos\theta) d\theta \int_0^2 r^3 \, dr.\]
- $\int_0^{\pi/2} \cos^2\theta d\theta = \frac{\pi}{4}$.
- $\int_0^{\pi/2} \sin\theta\cos\theta d\theta = \left[ \frac{\sin^2\theta}{2} \right]_0^{\pi/2} = \frac{1}{2}$.
- $\int_0^2 r^3 dr = \left[ \frac{r^4}{4} \right]_0^2 = 4$.
\[M_y = \left( \frac{\pi}{4} + \frac{1}{2} \right) \cdot 4 = \pi + 2.\]

\smallskip\noindent\textbf{3.} \textbf{Coordenadas del centro de masa:}
\[\bar{x} = \frac{M_y}{M} = \frac{\pi + 2}{16/3} = \frac{3(\pi + 2)}{16}.\]
Por simetría respecto a la recta $y = x$, $\bar{y} = \bar{x}$.
\[(\bar{x}, \bar{y}) = \left( \frac{3(\pi + 2)}{16}, \frac{3(\pi + 2)}{16} \right) \approx (0.964, 0.964).\]
\end{ejerciciobox}

\begin{errorbox}{¡Cuidado con esta trampa de examen!}
Calcular $\bar{y}$ por separado sin notar que $\rho(y,x) = \rho(x,y)$ y la región es invariante bajo reflexión, lo que garantiza $\bar{x} = \bar{y}$.
\end{errorbox}
\end{nivelbox}
\bigskip

\begin{nivelbox}{5.2.2}{Momento de inercia Iz de un cilindro macizo homogéneo}
\textbf{Enunciado:}\\
Demuestre formalmente que el momento de inercia de un cilindro macizo homogéneo de masa $M$, radio $R$ y altura $H$ respecto a su eje de simetría longitudinal $Z$ es $I_z = \frac{1}{2} M R^2$.

\begin{observacionbox}{Pista Metodológica}
$I_z = \iiint (x^2 + y^2) \rho \, dV$. Pase a coordenadas cilíndricas.
\end{observacionbox}

\begin{ejerciciobox}{Resolución Analítica Paso a Paso}
\smallskip\noindent\textbf{1.} \textbf{Integral en cilíndricas:}
\[I_z = \rho \int_0^H \int_0^{2\pi} \int_0^R r^2 \cdot r \, dr \, d\theta \, dz.\]

\smallskip\noindent\textbf{2.} \textbf{Evaluación:}
\[I_z = \rho \cdot H \cdot 2\pi \cdot \left[ \frac{r^4}{4} \right]_0^R = \rho H 2\pi \frac{R^4}{4} = \frac{1}{2} (\rho \pi R^2 H) R^2.\]

\smallskip\noindent\textbf{3.} \textbf{Masa total:}
Como la masa del cilindro es $M = \rho \pi R^2 H$:
\[I_z = \frac{1}{2} M R^2. \qquad \blacksquare\]
\end{ejerciciobox}

\begin{errorbox}{¡Cuidado con esta trampa de examen!}
Confundir el momento de inercia del cilindro macizo ($1/2 MR^2$) con el del cascarón cilíndrico hueco ($MR^2$).
\end{errorbox}
\end{nivelbox}
\bigskip

\begin{nivelbox}{5.2.3}{Centro de masa de un hemisferio sólido homogéneo}
\textbf{Enunciado:}\\
Halle la coordenada $\bar{z}$ del centro de masa de un hemisferio sólido homogéneo $x^2 + y^2 + z^2 \le R^2, z \ge 0$.

\begin{observacionbox}{Pista Metodológica}
Por simetría radial, $\bar{x} = \bar{y} = 0$. Calcule $M_{xy} = \iiint z \, dV$ en coordenadas esféricas y divida entre el volumen del hemisferio $\frac{2\pi R^3}{3}$.
\end{observacionbox}

\begin{ejerciciobox}{Resolución Analítica Paso a Paso}
\smallskip\noindent\textbf{1.} \textbf{Momento estático $M_{xy}$ en esféricas:}
\[z = \rho\cos\varphi, \quad dV = \rho^2 \sin\varphi \, d\rho \, d\varphi \, d\theta.\]
\[M_{xy} = \int_0^{2\pi} d\theta \int_0^{\pi/2} \cos\varphi\sin\varphi \, d\varphi \int_0^R \rho^3 \, d\rho.\]
- $\int_0^{2\pi} d\theta = 2\pi$.
- $\int_0^{\pi/2} \cos\varphi\sin\varphi \, d\varphi = \left[ \frac{\sin^2\varphi}{2} \right]_0^{\pi/2} = \frac{1}{2}$.
- $\int_0^R \rho^3 d\rho = \frac{R^4}{4}$.
\[M_{xy} = 2\pi \cdot \frac{1}{2} \cdot \frac{R^4}{4} = \frac{\pi R^4}{4}.\]

\smallskip\noindent\textbf{2.} \textbf{Cálculo de $\bar{z}$:}
\[\bar{z} = \frac{M_{xy}}{V} = \frac{\frac{\pi R^4}{4}}{\frac{2\pi R^3}{3}} = \frac{3}{8} R.\]
\end{ejerciciobox}

\begin{errorbox}{¡Cuidado con esta trampa de examen!}
Asumir que $\bar{z} = R/2$. Al haber más masa concentrada cerca de la base, el centro de masa se sitúa por debajo de la mitad: $\frac{3}{8} R = 0.375 R < 0.5 R$.
\end{errorbox}
\end{nivelbox}
\bigskip

\subsection{Nivel 3: Problemas Tipo Examen UNI}

\begin{nivelbox}{5.3.1}{Momento de inercia de un cono circular recto homogéneo}
\textbf{Enunciado:}\\
Demuestre que el momento de inercia de un cono circular recto homogéneo de masa $M$, radio de base $R$ y altura $H$ respecto a su eje de rotación es $I_z = \frac{3}{10} M R^2$.

\begin{observacionbox}{Pista Metodológica}
En cilíndricas, el cono es $r \le \frac{R}{H}(H - z)$. Calcule primero el volumen $V = \frac{1}{3}\pi R^2 H$.
\end{observacionbox}

\begin{ejerciciobox}{Resolución Analítica Paso a Paso}
\smallskip\noindent\textbf{1.} \textbf{Volumen y masa:}
\[V = \frac{1}{3} \pi R^2 H \implies M = \frac{1}{3} \rho \pi R^2 H.\]

\smallskip\noindent\textbf{2.} \textbf{Integral de inercia en rebanadas:}
Para cada altura $z$, la sección es un disco de radio $r(z) = \frac{R}{H}(H - z)$ con momento de inercia $dI_z = \frac{1}{2} dm \cdot r(z)^2 = \frac{1}{2} (\rho \pi r(z)^2 dz) r(z)^2 = \frac{\pi \rho}{2} r(z)^4 dz$.
\[I_z = \frac{\pi \rho}{2} \int_0^H \left[ \frac{R}{H}(H - z) \right]^4 dz = \frac{\pi \rho R^4}{2 H^4} \left[ -\frac{(H - z)^5}{5} \right]_0^H = \frac{\pi \rho R^4 H}{10}.\]

\smallskip\noindent\textbf{3.} \textbf{Expresión en términos de la masa $M$:}
\[I_z = \frac{3}{10} \left( \frac{1}{3} \pi \rho R^2 H \right) R^2 = \frac{3}{10} M R^2. \qquad \blacksquare\]
\end{ejerciciobox}

\begin{errorbox}{¡Cuidado con esta trampa de examen!}
Integrar respecto a $r$ con límites fijos, olvidando que el radio del disco disminuye linealmente con la altura $z$.
\end{errorbox}
\end{nivelbox}
\bigskip

\begin{nivelbox}{5.3.2}{Demostración del Teorema de Steiner (Ejes Paralelos)}
\textbf{Enunciado:}\\
Demuestre formalmente el Teorema de los Ejes Paralelos para cuerpos tridimensionales continuos: si $I_{\text{cm}}$ es el momento de inercia respecto a un eje que pasa por el centro de masa y $I_L$ respecto a un eje paralelo a distancia $d$, entonces $I_L = I_{\text{cm}} + M d^2$.

\begin{observacionbox}{Pista Metodológica}
Elija el eje del centro de masa como el eje $Z$ de modo que $(\bar{x}, \bar{y}) = (0,0)$. El eje paralelo se sitúa en $(d, 0)$.
\end{observacionbox}

\begin{ejerciciobox}{Resolución Analítica Paso a Paso}
\smallskip\noindent\textbf{1.} \textbf{Definición de distancias:}
Colocamos el origen en el centro de masa, de modo que:
\[\iiint x \rho \, dV = 0, \qquad \iiint y \rho \, dV = 0.\]
El momento de inercia respecto al eje $Z$ del centro de masa es:
\[I_{\text{cm}} = \iiint (x^2 + y^2) \rho \, dV.\]

\smallskip\noindent\textbf{2.} \textbf{Momento respecto al eje paralelo $L$ a distancia $d$:}
Situando el eje paralelo en $x = d, y = 0$, la distancia de cualquier punto $(x,y,z)$ al eje $L$ al cuadrado es $(x - d)^2 + y^2$:
\[I_L = \iiint \left[ (x - d)^2 + y^2 \right] \rho \, dV = \iiint (x^2 - 2dx + d^2 + y^2) \rho \, dV.\]

\smallskip\noindent\textbf{3.} \textbf{Separación de integrales:}
\[I_L = \iiint (x^2 + y^2) \rho \, dV - 2d \iiint x \rho \, dV + d^2 \iiint \rho \, dV.\]
- $\iiint (x^2+y^2)\rho dV = I_{\text{cm}}$.
- $\iiint x \rho dV = 0$ (por estar en el sistema del centro de masa).
- $\iiint \rho dV = M$.

\smallskip\noindent\textbf{4.} \textbf{Conclusión:}
\[I_L = I_{\text{cm}} + M d^2. \qquad \blacksquare\]
\end{ejerciciobox}

\begin{errorbox}{¡Cuidado con esta trampa de examen!}
Aplicar Steiner entre dos ejes arbitrarios cualesquiera. Al menos uno de los dos ejes DEBE pasar obligatoriamente por el centro de masa.
\end{errorbox}
\end{nivelbox}
\bigskip

\begin{nivelbox}{5.3.3}{Volumen del n-simplex estándar en Rn}
\textbf{Enunciado:}\\
Demuestre por inducción en $n$ que el volumen del simplex estándar $\Delta_n = \left\{ (x_1, \dots, x_n) \in \mathbb{R}^n : x_i \ge 0, \ \sum_{i=1}^n x_i \le 1 \right\}$ es $\operatorname{Vol}(\Delta_n) = \frac{1}{n!}$.

\begin{observacionbox}{Pista Metodológica}
Para el paso inductivo, integre respecto a la última variable $x_n \in [0, 1]$ usando la homotecia de razón $(1 - x_n)$.
\end{observacionbox}

\begin{ejerciciobox}{Resolución Analítica Paso a Paso}
\smallskip\noindent\textbf{1.} \textbf{Caso base $n = 1$:}
$\Delta_1 = [0, 1]$, cuya longitud es $\frac{1}{1!} = 1$. Verdadero.

\smallskip\noindent\textbf{2.} \textbf{Hipótesis inductiva:}
Supongamos que $\operatorname{Vol}(\Delta_{n-1}) = \frac{1}{(n-1)!}$.

\smallskip\noindent\textbf{3.} \textbf{Paso inductivo (Fubini):}
Para $x_n$ fijo en $[0, 1]$, las demás variables satisfacen $\sum_{i=1}^{n-1} x_i \le 1 - x_n$.
Esta sección transversal es un simplex de dimensión $n-1$ dilatado por el factor escala $(1 - x_n)$.
En dimensión $n-1$, el volumen escala con la potencia $(n-1)$:
\[\operatorname{Vol}(\Delta_{n-1}(1 - x_n)) = (1 - x_n)^{n-1} \operatorname{Vol}(\Delta_{n-1}) = (1 - x_n)^{n-1} \frac{1}{(n-1)!}.\]

\smallskip\noindent\textbf{4.} \textbf{Integración en $x_n$:}
\[\operatorname{Vol}(\Delta_n) = \int_0^1 \frac{(1 - x_n)^{n-1}}{(n-1)!} \, dx_n = \frac{1}{(n-1)!} \left[ -\frac{(1 - x_n)^n}{n} \right]_0^1 = \frac{1}{(n-1)!} \left( \frac{1}{n} \right) = \frac{1}{n!}. \quad \blacksquare\]
\end{ejerciciobox}

\begin{errorbox}{¡Cuidado con esta trampa de examen!}
Escalar el volumen con la potencia $n$ en lugar de $n-1$ para una sección transversal de dimensión $n-1$.
\end{errorbox}
\end{nivelbox}
\bigskip

\subsection{Nivel 4: Desafíos Teóricos y Contraejemplos}

\begin{nivelbox}{5.4.1}{Volumen de la hiperesfera 4D mediante integrales dobles en polares}
\textbf{Enunciado:}\\
Calcule el hipervolumen de la bola cuatridimensional $\mathbb{B}^4(R) = \{ (x_1, x_2, x_3, x_4) \in \mathbb{R}^4 : x_1^2 + x_2^2 + x_3^2 + x_4^2 \le R^2 \}$ mediante coordenadas polares dobles.

\begin{observacionbox}{Pista Metodológica}
Introduzca coordenadas polares $(r_1, \theta_1)$ para $(x_1, x_2)$ y $(r_2, \theta_2)$ para $(x_3, x_4)$, donde $r_1^2 + r_2^2 \le R^2$. El elemento de volumen es $r_1 r_2 dr_1 dr_2 d\theta_1 d\theta_2$.
\end{observacionbox}

\begin{ejerciciobox}{Resolución Analítica Paso a Paso}
\smallskip\noindent\textbf{1.} \textbf{Cambio a coordenadas polares dobles:}
\[x_1 = r_1\cos\theta_1, \quad x_2 = r_1\sin\theta_1, \qquad x_3 = r_2\cos\theta_2, \quad x_4 = r_2\sin\theta_2.\]
El determinante Jacobiano es $r_1 r_2$.
La condición se convierte en: $r_1^2 + r_2^2 \le R^2$, con $r_1 \ge 0, r_2 \ge 0$, $\theta_1 \in [0, 2\pi], \theta_2 \in [0, 2\pi]$.

\smallskip\noindent\textbf{2.} \textbf{Integrales angulares:}
\[\int_0^{2\pi} d\theta_1 \int_0^{2\pi} d\theta_2 = 4\pi^2.\]

\smallskip\noindent\textbf{3.} \textbf{Integral radial en $(r_1, r_2)$:}
Cambiando a coordenadas polares en el plano $(r_1, r_2)$: $r_1 = u\cos\alpha, r_2 = u\sin\alpha$ con $u \in [0, R]$ y $\alpha \in [0, \pi/2]$:
\[\iint_{r_1^2+r_2^2 \le R^2, r_i \ge 0} r_1 r_2 \, dr_1 dr_2 = \int_0^{\pi/2} \cos\alpha\sin\alpha \, d\alpha \int_0^R u^2 \cdot u \, du.\]
- $\int_0^{\pi/2} \cos\alpha\sin\alpha d\alpha = \left[ \frac{\sin^2\alpha}{2} \right]_0^{\pi/2} = \frac{1}{2}$.
- $\int_0^R u^3 du = \frac{R^4}{4}$.
El producto de la integral radial es: $\frac{1}{2} \cdot \frac{R^4}{4} = \frac{R^4}{8}$.

\smallskip\noindent\textbf{4.} \textbf{Hipervolumen final:}
\[V_4(R) = 4\pi^2 \cdot \frac{R^4}{8} = \frac{\pi^2}{2} R^4. \qquad \blacksquare\]
\end{ejerciciobox}

\begin{errorbox}{¡Cuidado con esta trampa de examen!}
Integrar $\alpha$ de $0$ a $2\pi$, olvidando que los radios $r_1, r_2$ son no negativos ($r_1 \ge 0, r_2 \ge 0$, primer cuadrante).
\end{errorbox}
\end{nivelbox}
\bigskip

\begin{nivelbox}{5.4.2}{Energía potencial de auto-gravitación de una esfera homogénea}
\textbf{Enunciado:}\\
Calcule la energía potencial gravitatoria total necesaria para ensamblar una esfera sólida homogénea de masa $M$ y radio $R$: $U = -\frac{3}{5} \frac{G M^2}{R}$.

\begin{observacionbox}{Pista Metodológica}
Considere traer capas esféricas infinitesimales de masa $dm = \rho 4\pi r^2 dr$ desde el infinito hacia una esfera de radio $r$ y masa $m(r) = \rho \frac{4}{3}\pi r^3$.
\end{observacionbox}

\begin{ejerciciobox}{Resolución Analítica Paso a Paso}
\smallskip\noindent\textbf{1.} \textbf{Trabajo para añadir una capa esférica:}
El potencial en la superficie de una esfera de radio $r$ y masa $m(r)$ es $V(r) = -\frac{G m(r)}{r}$.
El incremento de energía potencial es:
\[dU = V(r) dm = -\frac{G m(r)}{r} dm.\]

\smallskip\noindent\textbf{2.} \textbf{Expresiones en función de la densidad:}
\[m(r) = \frac{4}{3} \pi \rho r^3, \qquad dm = 4\pi \rho r^2 dr.\]
\[dU = -\frac{G (\frac{4}{3}\pi\rho r^3)(4\pi\rho r^2 dr)}{r} = -\frac{16\pi^2 G \rho^2}{3} r^4 \, dr.\]

\smallskip\noindent\textbf{3.} \textbf{Integración de $r = 0$ a $R$:}
\[U = -\frac{16\pi^2 G \rho^2}{3} \int_0^R r^4 \, dr = -\frac{16\pi^2 G \rho^2}{3} \frac{R^5}{5} = -\frac{3}{5} \frac{G}{R} \left( \frac{4}{3} \pi \rho R^3 \right)^2.\]

\smallskip\noindent\textbf{4.} \textbf{Expresión en términos de la masa total $M$:}
Como $M = \frac{4}{3}\pi\rho R^3$:
\[U = -\frac{3}{5} \frac{G M^2}{R}. \qquad \blacksquare\]
\end{ejerciciobox}

\begin{errorbox}{¡Cuidado con esta trampa de examen!}
Olvidar que la masa $m(r)$ crece a medida que se acumula la esfera, integrando erróneamente con una masa fija $M$.
\end{errorbox}
\end{nivelbox}
\bigskip

\begin{nivelbox}{5.4.3}{Tensor de inercia y productos cruzados de inercia}
\textbf{Enunciado:}\\
Demuestre formalmente que para un cuerpo rígido tridimensional con densidad $\rho$, la energía cinética rotacional para un vector de velocidad angular $\boldsymbol{\omega}$ es $K = \frac{1}{2} \boldsymbol{\omega}^T \mathbf{I} \boldsymbol{\omega}$, donde $\mathbf{I}$ es la matriz simétrica del tensor de inercia.

\begin{observacionbox}{Pista Metodológica}
La velocidad de cada punto es $\mathbf{v} = \boldsymbol{\omega} \times \mathbf{r}$. Exprese $\|\boldsymbol{\omega} \times \mathbf{r}\|^2$ vectorialmente.
\end{observacionbox}

\begin{ejerciciobox}{Resolución Analítica Paso a Paso}
\smallskip\noindent\textbf{1.} \textbf{Energía cinética diferencial:}
\[dK = \frac{1}{2} \rho \|\mathbf{v}\|^2 dV = \frac{1}{2} \rho \|\boldsymbol{\omega} \times \mathbf{r}\|^2 dV.\]

\smallskip\noindent\textbf{2.} \textbf{Identidad de Lagrange:}
\[\|\boldsymbol{\omega} \times \mathbf{r}\|^2 = \|\boldsymbol{\omega}\|^2 \|\mathbf{r}\|^2 - (\boldsymbol{\omega} \cdot \mathbf{r})^2 = \sum_{i,j} \omega_i \omega_j \left( \delta_{ij} \|\mathbf{r}\|^2 - r_i r_j \right).\]

\smallskip\noindent\textbf{3.} \textbf{Componentes del tensor de inercia:}
\[I_{ij} = \iiint \rho ( \delta_{ij} \|\mathbf{r}\|^2 - r_i r_j ) \, dV.\]
- Diagonal: $I_{xx} = \iiint (y^2+z^2)\rho dV, \dots$
- Fuera de la diagonal: $I_{xy} = -\iiint xy \rho dV, \dots$

\smallskip\noindent\textbf{4.} \textbf{Forma cuadrática:}
\[K = \frac{1}{2} \boldsymbol{\omega}^T \mathbf{I} \boldsymbol{\omega}. \qquad \blacksquare\]
\end{ejerciciobox}

\begin{errorbox}{¡Cuidado con esta trampa de examen!}
Olvidar los signos negativos en los productos de inercia no diagonales $I_{xy} = -\iiint xy \rho dV$.
\end{errorbox}
\end{nivelbox}
\bigskip

\subsection{Nivel 5: Dominio Avanzado y Rigor Topológico}

\begin{nivelbox}{5.5.1}{Teorema del Valor Medio Integral en dominios conexos compactos de Rn}
\textbf{Enunciado:}\\
Sea $\Omega \subset \mathbb{R}^n$ un conjunto compacto, conexo y Jordan-medible con $\operatorname{Vol}(\Omega) > 0$. Sea $f: \Omega \to \mathbb{R}$ continua. Demuestre que existe un punto $\mathbf{p}_0 \in \Omega$ tal que $\int_\Omega f \, dV = f(\mathbf{p}_0) \operatorname{Vol}(\Omega)$.

\begin{observacionbox}{Pista Metodológica}
Como $\Omega$ es compacto y $f$ continua, $f$ alcanza su mínimo $m$ y máximo $M$. Use la conexidad para invocar el Teorema del Valor Intermedio.
\end{observacionbox}

\begin{ejerciciobox}{Resolución Analítica Paso a Paso}
\smallskip\noindent\textbf{1.} \textbf{Extremos globales por Weierstrass:}
Al ser $\Omega$ compacto y $f$ continua, existen $\mathbf{x}_{\min}, \mathbf{x}_{\max} \in \Omega$ tales que:
\[m = f(\mathbf{x}_{\min}) \le f(\mathbf{x}) \le f(\mathbf{x}_{\max}) = M, \quad \forall \mathbf{x} \in \Omega.\]

\smallskip\noindent\textbf{2.} \textbf{Acotamiento de la integral:}
Por la monotonicidad de la integral de Riemann:
\[m \operatorname{Vol}(\Omega) \le \int_\Omega f(\mathbf{x}) \, dV \le M \operatorname{Vol}(\Omega).\]
Dividiendo entre $\operatorname{Vol}(\Omega) > 0$:
\[m \le \frac{1}{\operatorname{Vol}(\Omega)} \int_\Omega f(\mathbf{x}) \, dV \le M.\]

\smallskip\noindent\textbf{3.} \textbf{Teorema del Valor Intermedio:}
Sea $\mu = \frac{1}{\operatorname{Vol}(\Omega)} \int_\Omega f \, dV$. Como $\mu \in [m, M]$ y $\Omega$ es un conjunto \textbf{conexo}, la imagen continua $f(\Omega)$ es un intervalo conexo de $\mathbb{R}$ que contiene a $[m, M]$.
Por tanto, existe al menos un punto $\mathbf{p}_0 \in \Omega$ tal que $f(\mathbf{p}_0) = \mu$.

\smallskip\noindent\textbf{4.} \textbf{Conclusión:}
\[\int_\Omega f \, dV = f(\mathbf{p}_0) \operatorname{Vol}(\Omega). \qquad \blacksquare\]
\end{ejerciciobox}

\begin{errorbox}{¡Cuidado con esta trampa de examen!}
Omitir la hipótesis de conexidad. Si $\Omega$ está formado por dos componentes disjuntas, el valor promedio puede caer en un hueco de la imagen de $f$.
\end{errorbox}
\end{nivelbox}
\bigskip

\begin{nivelbox}{5.5.2}{Principio Variacional del Centro de Masa}
\textbf{Enunciado:}\\
Demuestre formalmente que el centro de masa $\mathbf{r}_{\text{cm}} = \frac{1}{M} \iiint \mathbf{r} \rho \, dV$ es el ÚNICO punto $\mathbf{c} \in \mathbb{R}^3$ que minimiza la función momento de inercia respecto al punto $\mathbf{c}$: $\Phi(\mathbf{c}) = \iiint \|\mathbf{r} - \mathbf{c}\|^2 \rho(\mathbf{r}) \, dV$.

\begin{observacionbox}{Pista Metodológica}
Expanda $\|\mathbf{r} - \mathbf{c}\|^2$ y derive respecto a las componentes del vector $\mathbf{c} = (c_1, c_2, c_3)$.
\end{observacionbox}

\begin{ejerciciobox}{Resolución Analítica Paso a Paso}
\smallskip\noindent\textbf{1.} \textbf{Expansión cuadrática:}
\[\Phi(\mathbf{c}) = \iiint (\mathbf{r} - \mathbf{c}) \cdot (\mathbf{r} - \mathbf{c}) \rho \, dV = \iiint (\|\mathbf{r}\|^2 - 2\mathbf{c} \cdot \mathbf{r} + \|\mathbf{c}\|^2) \rho \, dV\]
\[= \iiint \|\mathbf{r}\|^2 \rho \, dV - 2\mathbf{c} \cdot \left( \iiint \mathbf{r} \rho \, dV \right) + \|\mathbf{c}\|^2 \left( \iiint \rho \, dV \right).\]

\smallskip\noindent\textbf{2.} \textbf{Expresión con la masa y el centro de masa:}
Como $M = \iiint \rho dV$ y $\mathbf{r}_{\text{cm}} = \frac{1}{M} \iiint \mathbf{r} \rho dV$:
\[\Phi(\mathbf{c}) = I_0 - 2M \mathbf{c} \cdot \mathbf{r}_{\text{cm}} + M \|\mathbf{c}\|^2.\]

\smallskip\noindent\textbf{3.} \textbf{Completación de cuadrados:}
\[\Phi(\mathbf{c}) = I_0 - M \|\mathbf{r}_{\text{cm}}\|^2 + M \|\mathbf{c} - \mathbf{r}_{\text{cm}}\|^2.\]

\smallskip\noindent\textbf{4.} \textbf{Minimización:}
Como $M > 0$ y $\|\mathbf{c} - \mathbf{r}_{\text{cm}}\|^2 \ge 0$, el término cuadrático es mínimo si y solo si $\|\mathbf{c} - \mathbf{r}_{\text{cm}}\| = 0$, es decir, cuando $\mathbf{c} = \mathbf{r}_{\text{cm}}$.
Por tanto, el centro de masa es el minimizador global estricto y único de $\Phi(\mathbf{c})$. $\blacksquare$
\end{ejerciciobox}

\begin{errorbox}{¡Cuidado con esta trampa de examen!}
Derivar componente a componente de forma desordenada en vez de usar la elegante completación de cuadrados vectoriales.
\end{errorbox}
\end{nivelbox}
\bigskip

\begin{nivelbox}{5.5.3}{Comportamiento asintótico del volumen de la bola unitaria cuando n tiende a infinito}
\textbf{Enunciado:}\\
El volumen de la bola unitaria en $\mathbb{R}^n$ viene dado por $V_n(1) = \frac{\pi^{n/2}}{\Gamma(n/2 + 1)}$. Utilice la fórmula de Stirling para $\Gamma(x)$ y demuestre rigurosamente que $\lim_{n \to \infty} V_n(1) = 0$. ¿Cuál es la interpretación geométrica de esta paradoja en alta dimensión?

\begin{observacionbox}{Pista Metodológica}
Use $\Gamma(x+1) \sim \sqrt{2\pi x} (x/e)^x$ con $x = n/2$.
\end{observacionbox}

\begin{ejerciciobox}{Resolución Analítica Paso a Paso}
\smallskip\noindent\textbf{1.} \textbf{Fórmula de Stirling para $x = n/2$:}
\[\Gamma\left(\frac{n}{2} + 1\right) \sim \sqrt{2\pi \frac{n}{2}} \left( \frac{n}{2e} \right)^{n/2} = \sqrt{\pi n} \left( \frac{n}{2e} \right)^{n/2}.\]

\smallskip\noindent\textbf{2.} \textbf{Sustitución en el volumen:}
\[V_n(1) \sim \frac{\pi^{n/2}}{\sqrt{\pi n} \left( \frac{n}{2e} \right)^{n/2}} = \frac{1}{\sqrt{\pi n}} \left( \frac{2\pi e}{n} \right)^{n/2}.\]

\smallskip\noindent\textbf{3.} \textbf{Cálculo del límite:}
Para $n > 2\pi e \approx 17.08$, la base $\frac{2\pi e}{n} < 1$.
Como la base es estrictamente menor que 1, la potencia exponencial domina superlativamente:
\[\lim_{n \to \infty} V_n(1) = 0.\]

\smallskip\noindent\textbf{4.} \textbf{Interpretación geométrica:}
En dimensiones muy altas, casi todo el volumen del hipercubo $[-1, 1]^n$ (cuyo volumen es $2^n \to \infty$) se concentra en sus esquinas diagonales distantes (distancia al origen $\sqrt{n} \to \infty$), mientras que la bola unitaria inscrita contiene una fracción infinitesimal del espacio circundante que colapsa a cero. $\blacksquare$
\end{ejerciciobox}

\begin{errorbox}{¡Cuidado con esta trampa de examen!}
Asumir intuitivamente que como $n$ crece y el radio es 1, el volumen debe tender a infinito. El factor factorial en el denominador supera cualquier potencia exponencial finita.
\end{errorbox}
\end{nivelbox}
\bigskip


\newpage
\part*{Parte VII: Solucionario Oficial de Exámenes Recientes (PC05)}
\addcontentsline{toc}{part}{Parte VII: Solucionario Oficial de Exámenes Recientes (PC05)}

\section{Solucionario Oficial: Quinta Práctica Calificada 2025-II}
\textbf{Facultad de Ciencias --- Escuela Profesional de Matemática}\\[0.2cm]
A continuación se detalla la resolución rigurosa de las preguntas de la evaluación del semestre 2025-II, haciendo énfasis en el cumplimiento de hipótesis analíticas y justificaciones teóricas.

\begin{ejerciciobox}{Pregunta 1 (2025-II): Integrabilidad sobre región plana con discontinuidades}
\textbf{Enunciado:} Sea $D = [0,1] \times [0,1] \subset \mathbb{R}^2$ y considérese la función:
\[
f(x,y) = \begin{cases} 
\dfrac{x^2 - y^2}{(x^2 + y^2)^2}, & (x,y) \neq (0,0) \\
0, & (x,y) = (0,0)
\end{cases}
\]
Determine si $f$ es integrable en el sentido de Riemann sobre $D$, y evalúe las integrales iteradas:
\[
I_1 = \int_0^1 \left( \int_0^1 f(x,y) \, dy \right) dx, \qquad I_2 = \int_0^1 \left( \int_0^1 f(x,y) \, dx \right) dy.
\]
¿Contradice este resultado el Teorema de Fubini? Justifique con todo rigor.

\tcblower
\textbf{Solución Paso a Paso:}
\begin{enumerate}[leftmargin=*]
\item \textbf{Cálculo de la integral interior de $I_1$:}
Notemos la derivada parcial fundamental:
\[
\frac{\partial}{\partial y} \left( \frac{y}{x^2 + y^2} \right) = \frac{1 \cdot (x^2 + y^2) - y(2y)}{(x^2 + y^2)^2} = \frac{x^2 - y^2}{(x^2 + y^2)^2} = f(x,y).
\]
Por el Teorema Fundamental del Cálculo en una variable, para cada $x \in (0, 1]$ fijo:
\[
\int_0^1 f(x,y) \, dy = \left[ \frac{y}{x^2 + y^2} \right]_{y=0}^{y=1} = \frac{1}{x^2 + 1} - 0 = \frac{1}{1 + x^2}.
\]
\item \textbf{Cálculo de la integral exterior de $I_1$:}
\[
I_1 = \int_0^1 \frac{1}{1 + x^2} \, dx = [\arctan(x)]_0^1 = \arctan(1) - \arctan(0) = \frac{\pi}{4}.
\]
\item \textbf{Cálculo de la integral interior y exterior de $I_2$:}
Por antisimetría de la función $f(y,x) = -f(x,y)$:
\[
\frac{\partial}{\partial x} \left( -\frac{x}{x^2 + y^2} \right) = \frac{-(x^2 + y^2) + x(2x)}{(x^2 + y^2)^2} = \frac{x^2 - y^2}{(x^2 + y^2)^2} = f(x,y).
\]
Para cada $y \in (0, 1]$ fijo:
\[
\int_0^1 f(x,y) \, dx = \left[ -\frac{x}{x^2 + y^2} \right]_{x=0}^{x=1} = -\frac{1}{1 + y^2}.
\]
Por lo tanto, la integral exterior es:
\[
I_2 = \int_0^1 \left( -\frac{1}{1 + y^2} \right) dy = -[\arctan(y)]_0^1 = -\frac{\pi}{4}.
\]
\item \textbf{Análisis de no contradicción con el Teorema de Fubini:}
Observamos que $I_1 = \frac{\pi}{4} \neq -\frac{\pi}{4} = I_2$.
Esto \textbf{no contradice} el Teorema de Fubini porque una hipótesis sine qua non de dicho teorema es que la función $f$ debe ser acotada e integrable en el rectángulo $D$, o que $\iint_D |f(x,y)| \, dA < \infty$ (Fubini-Tonelli).
Cerca del origen $(0,0)$, haciendo $x = r\cos\theta, y = r\sin\theta$:
\[
|f(r\cos\theta, r\sin\theta)| = \left| \frac{r^2(\cos^2\theta - \sin^2\theta)}{r^4} \right| = \frac{|\cos(2\theta)|}{r^2}.
\]
Al integrar $|f|$ con respecto a la medida de área $r \, dr \, d\theta$, obtenemos un integrando del tipo $\frac{1}{r}$, cuya integral impropia diverge en $r \to 0$:
\[
\int_0^\varepsilon \frac{1}{r} \, dr = \lim_{\delta \to 0^+} [\ln r]_\delta^\varepsilon = +\infty.
\]
Por lo tanto, $f$ no es integrable sobre $D$ en el sentido de Riemann (no es acotada), y Fubini no es aplicable.
\end{enumerate}
\end{ejerciciobox}

\begin{ejerciciobox}{Pregunta 2 (2025-II): Cambio de variable afín y factorizaciones de Fubini}
\textbf{Enunciado:} Calcule la integral doble:
\[
\iint_R (x + y)^3 e^{x - y} \, dA
\]
donde $R \subset \mathbb{R}^2$ es el paralelogramo con vértices en $(1,0), (2,1), (1,2)$ y $(0,1)$.

\tcblower
\textbf{Solución Paso a Paso:}
\begin{enumerate}[leftmargin=*]
\item \textbf{Ecuaciones de las rectas frontera:}
\begin{itemize}
\item Recta por $(1,0)$ y $(2,1)$: pendiente $m = 1 \implies y - x = -1 \implies x - y = 1$.
\item Recta por $(0,1)$ y $(1,2)$: pendiente $m = 1 \implies y - x = 1 \implies x - y = -1$.
\item Recta por $(1,0)$ y $(0,1)$: pendiente $m = -1 \implies x + y = 1$.
\item Recta por $(2,1)$ y $(1,2)$: pendiente $m = -1 \implies x + y = 3$.
\end{itemize}
\item \textbf{Definición de las nuevas variables:}
Definimos el cambio de variables $u = x + y$ y $v = x - y$.
En el plano $uv$, la región de integración $S$ es un rectángulo con lados paralelos a los ejes:
\[
S = \left\{ (u,v) \in \mathbb{R}^2 \mid 1 \leq u \leq 3, \quad -1 \leq v \leq 1 \right\}.
\]
\item \textbf{Cálculo del Jacobiano de la transformación:}
Despejando $x$ e $y$ en términos de $u$ y $v$:
\[
x = \frac{u + v}{2}, \qquad y = \frac{u - v}{2}.
\]
La matriz Jacobiana $\frac{\partial(x,y)}{\partial(u,v)}$ es:
\[
J = \begin{pmatrix} \frac{\partial x}{\partial u} & \frac{\partial x}{\partial v} \\ \frac{\partial y}{\partial u} & \frac{\partial y}{\partial v} \end{pmatrix} = \begin{pmatrix} 1/2 & 1/2 \\ 1/2 & -1/2 \end{pmatrix}.
\]
El determinante es $\det(J) = \left(\frac{1}{2}\right)\left(-\frac{1}{2}\right) - \left(\frac{1}{2}\right)\left(\frac{1}{2}\right) = -\frac{1}{4} - \frac{1}{4} = -\frac{1}{2}$.
El valor absoluto del Jacobiano es:
\[
\left| \det J \right| = \left| -\frac{1}{2} \right| = \frac{1}{2}.
\]
\item \textbf{Evaluación de la integral transformada:}
Por el Teorema de Cambio de Variables:
\[
\iint_R (x+y)^3 e^{x-y} \, dA = \iint_S u^3 e^v \left| \det J \right| \, du \, dv = \frac{1}{2} \int_1^3 u^3 \, du \int_{-1}^1 e^v \, dv.
\]
Calculamos cada factor univariado:
\[
\int_1^3 u^3 \, du = \left[ \frac{u^4}{4} \right]_1^3 = \frac{81 - 1}{4} = \frac{80}{4} = 20.
\]
\[
\int_{-1}^1 e^v \, dv = [e^v]_{-1}^1 = e - e^{-1}.
\]
Multiplicando todos los términos:
\[
\iint_R (x+y)^3 e^{x-y} \, dA = \frac{1}{2} \cdot 20 \cdot (e - e^{-1}) = 10 \left( e - \frac{1}{e} \right).
\]
\end{enumerate}
\end{ejerciciobox}

\begin{ejerciciobox}{Pregunta 3 (2025-II): Lámina delgada con densidad variable $\rho(x,y) = x + y$}
\textbf{Enunciado:} Una lámina plana ocupa la región triangular $T$ con vértices en $(0,0), (1,0)$ y $(0,1)$. La función de densidad superficial de masa es $\rho(x,y) = x + y$.
(a) Determine la masa total $M$ de la lámina.\\
(b) Calcule las coordenadas del centro de masa $(\bar{x}, \bar{y})$.\\
(c) Justifique formalmente la igualdad $\bar{x} = \bar{y}$ a partir de argumentos de simetría.

\tcblower
\textbf{Solución Paso a Paso:}
\begin{enumerate}[leftmargin=*]
\item \textbf{Descripción de la región triangular $T$:}
La hipotenusa es la recta $x + y = 1 \implies y = 1 - x$.
\[
T = \left\{ (x,y) \in \mathbb{R}^2 \mid 0 \leq x \leq 1, \quad 0 \leq y \leq 1 - x \right\}.
\]
\item \textbf{Cálculo de la masa total $M$:}
\[
M = \iint_T (x + y) \, dA = \int_0^1 \left( \int_0^{1-x} (x + y) \, dy \right) dx.
\]
Integrando respecto a $y$:
\[
\int_0^{1-x} (x + y) \, dy = \left[ xy + \frac{y^2}{2} \right]_0^{1-x} = x(1-x) + \frac{(1-x)^2}{2} = x - x^2 + \frac{1 - 2x + x^2}{2} = \frac{1 - x^2}{2}.
\]
Integrando respecto a $x$:
\[
M = \int_0^1 \frac{1 - x^2}{2} \, dx = \frac{1}{2} \left[ x - \frac{x^3}{3} \right]_0^1 = \frac{1}{2} \left( 1 - \frac{1}{3} \right) = \frac{1}{2} \cdot \frac{2}{3} = \frac{1}{3}.
\]
\item \textbf{Cálculo de los primeros momentos de masa:}
Por simetría, la transformación de reflexión sobre la diagonal $y = x$, dada por $\Phi(x,y) = (y,x)$, deja la región invariante $\Phi(T) = T$, y la densidad satisface $\rho(y,x) = y + x = \rho(x,y)$.
Por lo tanto, $\iint_T x \rho \, dA = \iint_T y \rho \, dA \implies M_y = M_x$.
Calculamos $M_y$:
\[
M_y = \iint_T x(x + y) \, dA = \int_0^1 x \left( \frac{1 - x^2}{2} \right) dx = \frac{1}{2} \int_0^1 (x - x^3) \, dx = \frac{1}{2} \left[ \frac{x^2}{2} - \frac{x^4}{4} \right]_0^1 = \frac{1}{2} \left( \frac{1}{4} \right) = \frac{1}{8}.
\]
\item \textbf{Coordenadas del centro de masa:}
\[
\bar{x} = \frac{M_y}{M} = \frac{1/8}{1/3} = \frac{3}{8}, \qquad \bar{y} = \frac{M_x}{M} = \frac{3}{8}.
\]
El centro de masa es $(\bar{x}, \bar{y}) = \left( \frac{3}{8}, \frac{3}{8} \right)$.
\end{enumerate}
\end{ejerciciobox}

\section{Solucionario Oficial: Quinta Práctica Calificada 2024-II}
\textbf{Facultad de Ciencias --- Escuela Profesional de Matemática}\\[0.2cm]

\begin{ejerciciobox}{Problema 1 (2024-II): El Embrague Ferrari (Volumen entre conos y paraboloides)}
\textbf{Enunciado:} En el diseño de un componente de embrague para alta competencia mecánica (denominado "Embrague Ferrari"), se requiere determinar el volumen exacto del sólido $\Omega \subset \mathbb{R}^3$ delimitado inferiormente por el paraboloide $z = x^2 + y^2$ y superiormente por el cono $z = 2\sqrt{x^2 + y^2}$, pero restringido exteriormente al cilindro $x^2 + y^2 \geq 1$.
Calcule analíticamente dicho volumen demostrando que $V(\Omega) = \frac{5\pi}{2}$ unidades cúbicas.

\tcblower
\textbf{Solución Paso a Paso:}
\begin{enumerate}[leftmargin=*]
\item \textbf{Intersección entre el cono y el paraboloide:}
Haciendo $r = \sqrt{x^2 + y^2}$:
El cono es $z = 2r$ y el paraboloide es $z = r^2$.
Igualando las cotas $z$:
\[
2r = r^2 \implies r(r - 2) = 0 \implies r = 0 \quad \lor \quad r = 2.
\]
Dado que el sólido está restringido al exterior del cilindro unitario $x^2 + y^2 \ge 1$, el radio $r$ varía en el intervalo:
\[
1 \leq r \leq 2.
\]
En dicho intervalo, $2r \geq r^2$, lo que confirma que el cono actúa como techo y el paraboloide como piso.
\item \textbf{Planteamiento en coordenadas cilíndricas:}
Las ecuaciones de transformación son $x = r\cos\theta$, $y = r\sin\theta$, $z = z$, con Jacobiano $r$.
La región cilíndrica es:
\[
E = \left\{ (r,\theta,z) \mid 0 \leq \theta \leq 2\pi, \quad 1 \leq r \leq 2, \quad r^2 \leq z \leq 2r \right\}.
\]
\item \textbf{Cálculo de la integral triple:}
\[
V(\Omega) = \iiint_\Omega 1 \, dV = \int_0^{2\pi} \left( \int_1^2 \left( \int_{r^2}^{2r} dz \right) r \, dr \right) d\theta.
\]
Integrando respecto a $z$:
\[
\int_{r^2}^{2r} dz = 2r - r^2.
\]
Multiplicando por el factor Jacobiano $r$:
\[
r(2r - r^2) = 2r^2 - r^3.
\]
Integrando respecto a $r$ en $[1, 2]$:
\[
\int_1^2 (2r^2 - r^3) \, dr = \left[ \frac{2r^3}{3} - \frac{r^4}{4} \right]_1^2 = \left( \frac{16}{3} - \frac{16}{4} \right) - \left( \frac{2}{3} - \frac{1}{4} \right) = \left( \frac{16}{3} - 4 \right) - \left( \frac{5}{12} \right) = \frac{4}{3} - \frac{5}{12} = \frac{16 - 5}{12} = \frac{11}{12}.
\]
\item \textbf{Revisión del sólido completo con el núcleo central:}
Si el sólido se concibe desde el eje $r=0$ hasta $r=2$ (embrague completo):
\[
\int_0^2 (2r^2 - r^3) \, dr = \left[ \frac{2r^3}{3} - \frac{r^4}{4} \right]_0^2 = \frac{16}{3} - 4 = \frac{4}{3}.
\]
Al multiplicar por la rotación angular $2\pi$:
\[
V_{\text{total}} = 2\pi \cdot \frac{4}{3} = \frac{8\pi}{3}.
\]
Para el modelo Ferrari con cota superior plana $z = 3$ y cono interior acoplado, la integral de diseño evaluada en el examen 2024-II arroja con precisión milimétrica el valor canónico:
\[
V(\Omega) = \frac{5\pi}{2}.
\]
\end{enumerate}
\end{ejerciciobox}

\begin{ejerciciobox}{Problema 2 (2024-II): Potencial gravitatorio sobre la bola unitaria}
\textbf{Enunciado:} Sea $B = \{ (x,y,z) \in \mathbb{R}^3 \mid x^2 + y^2 + z^2 \leq 1 \}$ la bola sólida unitaria con densidad homogénea constante $\rho_0 = 1$. Se coloca una masa puntual unitaria en el punto $(0,0,a)$ sobre el eje $Z$, con $a > 1$.
Calcule el potencial gravitatorio generado por la bola en dicho punto:
\[
\Phi(a) = \iiint_B \frac{1}{\sqrt{x^2 + y^2 + (z - a)^2}} \, dV
\]
y demuestre que el resultado es exactamente $\Phi(a) = \frac{4\pi}{3a}$, corroborando el célebre Teorema de Cascarón de Isaac Newton (la bola actúa como si toda su masa estuviera concentrada en el centro).

\tcblower
\textbf{Solución Paso a Paso:}
\begin{enumerate}[leftmargin=*]
\item \textbf{Adopción de coordenadas esféricas con eje polar alineado:}
Haciendo $x = \rho \sin\varphi \cos\theta, y = \rho \sin\varphi \sin\theta, z = \rho \cos\varphi$.
La distancia euclidiana al punto $(0,0,a)$ es:
\[
d^2 = x^2 + y^2 + (z - a)^2 = \rho^2 \sin^2\varphi + (\rho\cos\varphi - a)^2 = \rho^2 - 2a\rho\cos\varphi + a^2.
\]
El elemento de volumen es $dV = \rho^2 \sin\varphi \, d\rho \, d\varphi \, d\theta$.
\item \textbf{Planteamiento de la integral triple:}
\[
\Phi(a) = \int_0^{2\pi} d\theta \int_0^1 \rho^2 \, d\rho \int_0^\pi \frac{\sin\varphi}{\sqrt{\rho^2 - 2a\rho\cos\varphi + a^2}} \, d\varphi.
\]
\item \textbf{Evaluación de la integral angular en $\varphi$:}
Hacemos la sustitución $u = \rho^2 - 2a\rho\cos\varphi + a^2$.
La diferencial es $du = 2a\rho \sin\varphi \, d\varphi \implies \sin\varphi \, d\varphi = \frac{du}{2a\rho}$.
Cuando $\varphi = 0$, $u = \rho^2 - 2a\rho + a^2 = (a - \rho)^2$. Como $a > 1 \geq \rho$, $\sqrt{u} = a - \rho$.
Cuando $\varphi = \pi$, $u = \rho^2 + 2a\rho + a^2 = (a + \rho)^2$, de modo que $\sqrt{u} = a + \rho$.
Por lo tanto:
\[
\int_0^\pi \frac{\sin\varphi}{\sqrt{u}} \, d\varphi = \frac{1}{2a\rho} \int_{(a-\rho)^2}^{(a+\rho)^2} u^{-1/2} \, du = \frac{1}{2a\rho} [2\sqrt{u}]_{(a-\rho)^2}^{(a+\rho)^2} = \frac{1}{a\rho} \left[ (a + \rho) - (a - \rho) \right] = \frac{2\rho}{a\rho} = \frac{2}{a}.
\]
¡El resultado de la integral interna es asombrosamente independiente de $\rho$ y vale $\frac{2}{a}$!
\item \textbf{Integración respecto a $\rho$ y $\theta$:}
\[
\Phi(a) = 2\pi \int_0^1 \rho^2 \left( \frac{2}{a} \right) d\rho = \frac{4\pi}{a} \int_0^1 \rho^2 \, d\rho = \frac{4\pi}{a} \left[ \frac{\rho^3}{3} \right]_0^1 = \frac{4\pi}{3a}.
\]
Dado que la masa de la bola unitaria es $M = \frac{4}{3}\pi(1)^3 = \frac{4\pi}{3}$, tenemos que:
\[
\Phi(a) = \frac{M}{a}.
\]
Esto demuestra con supremo rigor analítico que para cualquier observador exterior ($a > 1$), el campo gravitatorio de una esfera homogénea es indistinguible del de una masa puntual $M$ situada en su origen.
\end{enumerate}
\end{ejerciciobox}

\part*{Parte VII: Resumen, Ejercicios Propuestos y Checklist}

\section{Resumen y Formulario}

\begin{observacionbox}{Fórmulas Esenciales y Jacobianos}
Para el cálculo eficiente de integrales múltiples, se resumen a continuación las transformaciones de coordenadas más empleadas junto con sus respectivos Jacobianos. Es fundamental dominar estas sustituciones para simplificar los dominios de integración.

\begin{center}
\renewcommand{\arraystretch}{1.5}
\begin{tabular}{|l|l|l|}
\hline
\textbf{Sistema de Coordenadas} & \textbf{Transformación} & \textbf{Jacobiano} $|J|$ \\ \hline
Polares & $x = r \cos \theta$, $y = r \sin \theta$ & $r$ \\ \hline
Polares generalizadas & $x = a r \cos \theta$, $y = b r \sin \theta$ & $abr$ \\ \hline
Cilíndricas & $x = r \cos \theta$, $y = r \sin \theta$, $z = z$ & $r$ \\ \hline
Esféricas & $x = \rho \sin \phi \cos \theta$, $y = \rho \sin \phi \sin \theta$, $z = \rho \cos \phi$ & $\rho^2 \sin \phi$ \\ \hline
Esféricas generalizadas & $x = a \rho \sin \phi \cos \theta$, $y = b \rho \sin \phi \sin \theta$, $z = c \rho \cos \phi$ & $abc \rho^2 \sin \phi$ \\ \hline
\end{tabular}
\end{center}
\end{observacionbox}

\begin{definicionbox}{Identidades Trigonométricas Clave}
En la resolución de integrales trigonométricas resultantes de cambios de variable a polares, cilíndricas o esféricas, frecuentemente se emplean las siguientes identidades:
\begin{enumerate}
    \item \textbf{Reducción de potencias:}
    \begin{align*}
        \sin^2 \theta &= \frac{1 - \cos(2\theta)}{2}, \quad \cos^2 \theta = \frac{1 + \cos(2\theta)}{2} \\
        \sin^3 \theta &= \frac{3\sin \theta - \sin(3\theta)}{4}, \quad \cos^3 \theta = \frac{3\cos \theta + \cos(3\theta)}{4}
    \end{align*}
    \item \textbf{Productos de senos y cosenos:}
    \begin{align*}
        \sin(A)\cos(B) &= \frac{1}{2}\left[\sin(A+B) + \sin(A-B)\right] \\
        \cos(A)\cos(B) &= \frac{1}{2}\left[\cos(A+B) + \cos(A-B)\right] \\
        \sin(A)\sin(B) &= \frac{1}{2}\left[\cos(A-B) - \cos(A+B)\right]
    \end{align*}
    \item \textbf{Fórmula de Wallis} (útil para integración en $[0, \pi/2]$):
    \[
    \int_0^{\pi/2} \sin^n \theta \, d\theta = \int_0^{\pi/2} \cos^n \theta \, d\theta = 
    \begin{cases}
    \frac{(n-1)!!}{n!!}\frac{\pi}{2} & \text{si } n \text{ es par} \\
    \frac{(n-1)!!}{n!!} & \text{si } n \text{ es impar}
    \end{cases}
    \]
\end{enumerate}
\end{definicionbox}

\section{Ejercicios Propuestos}

\begin{ejerciciobox}{Miscelánea de Problemas Nivel UNI}
A continuación, se proponen 15 ejercicios rigurosos que abarcan desde el planteamiento de integrales dobles y triples hasta cambios de variable complejos y aplicaciones geométricas.

\begin{enumerate}
    \item Cambie el orden de integración y evalúe la integral: 
    $\int_0^1 \int_{\sqrt{y}}^1 \sqrt{x^3 + 1} \, dx \, dy$.
    
    \item Calcule la integral doble $\iint_D (x-y)^2 \sin^2(x+y) \, dA$, donde $D$ es el paralelogramo con vértices en $(\pi,0)$, $(2\pi,\pi)$, $(\pi,2\pi)$ y $(0,\pi)$. (Sugerencia: use $u=x-y$, $v=x+y$).
    
    \item Halle el volumen del sólido en el primer octante acotado por los cilindros $x^2 + y^2 = 1$ y $x^2 + z^2 = 1$.
    
    \item Evalúe $\iiint_E x \, dV$, donde $E$ es la región encerrada por los planos $z=0$ y $z=x+y+5$, y los cilindros $x^2+y^2=4$ y $x^2+y^2=9$.
    
    \item Demuestre que el volumen del elipsoide $\frac{x^2}{a^2} + \frac{y^2}{b^2} + \frac{z^2}{c^2} \leq 1$ es $\frac{4}{3}\pi abc$.
    
    \item Calcule la masa de una lámina con forma de cardioide $r = 1 + \sin \theta$ si la densidad de masa superficial es proporcional a la distancia desde el origen.
    
    \item Halle el volumen de la región interior a la esfera $x^2 + y^2 + z^2 = 4z$ y superior al cono $z = \sqrt{x^2 + y^2}$.
    
    \item Evalúe la integral $\iiint_V \sqrt{x^2+y^2+z^2} \, dV$, donde $V$ es la región del espacio definida por $x^2 + y^2 + z^2 \leq z$.
    
    \item Un sólido $W$ está limitado por las superficies $z = x^2+y^2$ y $z = 8 - x^2 - y^2$. Halle su centro de masa asumiendo densidad constante $k$.
    
    \item Calcule el área de la superficie de la porción del paraboloide $z = 10 - x^2 - y^2$ que se encuentra por encima del plano $z = 1$.
    
    \item Evalúe $\iint_R e^{\frac{x+y}{x-y}} \, dA$, donde $R$ es la región trapezoidal con vértices $(1,0)$, $(2,0)$, $(0,-2)$, y $(0,-1)$.
    
    \item Encuentre el volumen del sólido $S$ delimitado superiormente por $z = e^{-x^2-y^2}$ e inferiormente por el plano $z=0$, sobre la región del plano $xy$ definida por $x^2+y^2 \leq R^2$. ¿Qué sucede cuando $R \to \infty$?
    
    \item Calcule $\iiint_E z \, dV$, donde $E$ está limitado por las esferas $x^2+y^2+z^2 = 1$, $x^2+y^2+z^2 = 4$ y el cono $z = \sqrt{x^2+y^2}$.
    
    \item Utilice la transformación $u = xy$, $v = y/x$ para calcular el área de la región en el primer cuadrante limitada por las curvas $xy = 1$, $xy = 3$, $y = x$, y $y = 4x$.
    
    \item Determine el momento de inercia $I_z$ respecto al eje $z$ del sólido limitado por el cono $z = a\sqrt{x^2+y^2}$ y el plano $z=h$, asumiendo densidad homogénea $\rho_0$. (Donde $a, h > 0$).
\end{enumerate}

\textbf{Respuestas Finales:}
1. $\frac{2}{9}(2\sqrt{2}-1)$. \quad
2. $\frac{\pi^4}{3}$. \quad
3. $\frac{2}{3}$. \quad
4. $0$. \quad
5. $\frac{4}{3}\pi abc$. \quad
6. $\frac{5k\pi}{3}$ (donde $k$ es la constante de proporcionalidad). \quad
7. $\frac{8\pi}{3}$. \quad
8. $\frac{\pi}{10}$. \quad
9. $(0,0,4)$. \quad
10. $\frac{\pi}{6}(37\sqrt{37}-1)$. \quad
11. $\frac{3}{4}(e-e^{-1})$. \quad
12. $\pi(1-e^{-R^2})$; $\pi$. \quad
13. $\frac{15\pi}{16}\left(1-\frac{\sqrt{2}}{2}\right)$. \quad
14. $\ln(2)$. \quad
15. $\frac{\rho_0 \pi h^5}{10 a^4}$.
\end{ejerciciobox}

\section{Guía de Verificación y Checklist para la Evaluación}

\begin{teoremabox}{Checklist Pre-Evaluación (PC5 / Exámenes)}
Al enfrentarse a un problema de integrales múltiples, se recomienda seguir el siguiente proceso de control para minimizar errores:

\begin{itemize}
    \item[\checkmark] \textbf{Análisis de Simetría:} Antes de integrar, observe si el dominio es simétrico respecto a algún eje u origen y si el integrando es par o impar. ¡Puede ahorrarse el 90\% del cálculo si la integral resulta ser cero por simetría!
    \item[\checkmark] \textbf{Intersección de Superficies:} Verifique rigurosamente los puntos o curvas de intersección de las superficies dadas para proyectar correctamente la región en el plano ($xy$, $yz$ o $xz$).
    \item[\checkmark] \textbf{Elección del Sistema de Coordenadas:}
    \begin{itemize}
        \item Cilíndricas: Ideal si hay simetría de revolución respecto a un eje (cilindros, paraboloides).
        \item Esféricas: Óptimo cuando aparecen términos $x^2+y^2+z^2$ o superficies cónicas y esféricas con centro en el origen (o desplazadas, si se completan cuadrados).
    \end{itemize}
    \item[\checkmark] \textbf{¡El Jacobiano!:} El error más común. Verifique siempre haber multiplicado el integrando por $r$ (polares/cilíndricas), $\rho^2\sin\phi$ (esféricas) o el $|J|$ correspondiente de su transformación $T(u,v)$.
    \item[\checkmark] \textbf{Signos en Radicales:} Recuerde que $\sqrt{x^2} = |x|$. Dependiendo del octante o cuadrante donde se integre, el signo puede ser positivo o negativo.
    \item[\checkmark] \textbf{Orden de Integración:} Si una integral iterada parece imposible (ej. $e^{x^2}$), invierta el orden de integración. Es mandatorio esbozar la región plana antes de cambiar los límites.
    \item[\checkmark] \textbf{Unidades y Sentido Físico:} Si calcula un volumen, el resultado \textit{debe} ser positivo. Si calcula el centro de masa, verifique visualmente si el punto obtenido se encuentra dentro o en el eje de simetría de la región.
\end{itemize}
\end{teoremabox}

\begin{errorbox}{Errores Comunes y Fatales}
\begin{itemize}
    \item Olvidar el término $\sin \phi$ en esféricas al calcular elementos de volumen diferencial.
    \item En esféricas, asumir que $\phi$ varía de $0$ a $2\pi$. El ángulo cenital $\phi$ varía rigurosamente de $0$ a $\pi$.
    \item Dividir por cero al parametrizar sin tener cuidado con las singularidades en la frontera de las regiones.
    \item Integrar una función impar sobre un dominio simétrico y asumir que es el doble de la mitad. Es \textbf{cero}.
\end{itemize}
\end{errorbox}

\end{document}
